% -----------------------------------------------------------
\section{1841}
% -----------------------------------------------------------

	% ----------------------
	\subsection{Joseph Smith}
	% ----------------------
		\label{EMS-1841-JS}

		% ----------------------
		\subsubsection{Introduction to 1841}
		% ----------------------
		
			sdf
			
		% -----
		\subsubsection{Discourse, 5 January 1841}
		% -----
			\label{EMS-1841-JS-5-JAN-1841}
			% \label{EMS-1840-JS-DAY-3LETTERMONTH-YEAR}
			
			% --
			\paragraph{As Reported by William P. McIntire}
			% --
				\label{EMS-1841-JS-5-JAN-1841-WM}
			
				\begin{description}
					\item[Print Sources]
				\end{description}

				\begin{tabular}{p{0.5in}p{1in}p{0.5in}p{1in}}
					JSP	 	& ---		& JSR 		& --- \\
					DC	 	& ---		& HC 		& --- \\
					HDDC 	& ---		& TCHC 		& --- \\
				\end{tabular}
				% HDDC = woodford's DC diss; JSR = Marquardt's JS Revelations; TCHC = Vogel HC
				
				\begin{description}
					\item[Online Access] \href{https://www.josephsmithpapers.org/paper-summary/account-of-meeting-and-discourse-5-january-1841-as-reported-by-william-p-mcintire/}{Here}
					% \item[Analysis] \hyperref[XX]{Here}
				\end{description}
				
				\begin{description}
					\item[Background]
				\end{description}
				
					% It might also be useful to give a two sentence context of the period: e.g., "This speech was given in Nauvoo to a small congregation one month prior to the destruction of the Nauvoo Expositor. These and other tensions may have been on the prophet's mind when he gave the speech."
				
				\begin{description}
					\item[Original Source]
				\end{description}
				
					% brief description of original source material, with reference to JSP or somewhere else for more info
					% Background material: it might be helpful to have a note that says whether a given document was presented as a speech, was written in a journal, etc.
					% Also a comment here about how reliable the source might be, or what shape it is in, if there are large omissions or parts that are hard to understand, and so on.
				
				\begin{description}
					\item[Text]
				\end{description}
				
						\hypertarget{EMS-1841-JS-5-JAN-1841-WM-a}%
					a
						\marginnote{a}%
					short minute of the subjects \& the most promonent matter as brought fourth from those subjects at Joseph's office Jan.--- 8\supersu{th}\supers{.} [5th] 1840 [1841]

					subject first.--- Discused by D. C. [Don Carlos] Smith; also this preciple [principle] practized by man; the blessings \& results of the same he said the pri[n]ciple would bind the H[e]arts of man togather \& give them confidence in each other \& as John says thy word is truth; so he says if we keep his word--- we shall all be actuated by the same principle \& be as one man; \& as angels are obedient to the same word we shall have Concorse to them \& also to all the Heavenly throng; Joseph said to D. C. Smith that to be free <from> the Coruption of the Earth that man the speaker should all ways speak in his natureal tone of voice; \& not to keep in one loud strain, but to act without affectation

						\hypertarget{EMS-1841-JS-5-JAN-1841-WM-b}%
					Next
						\marginnote{b}%
					subject was--- did the Lord God <make> the Earth out of nothing;--- By D. Ells--- says he God did not make th[e] earth out of nothing--- for it is contrary to <a> rashanal [rational] mind \& reason. that a something could be brought from--- a nothing;--- also it is contrary to the principle \& means by witch God does work; for instance; when God fo[r]med man, he made him of something; the dust of the Earth, \& he allways took a somthing to afect a something Else; oft he takes man <to> scorge his fellow man, or watter to destroy man--- or fire to distroy man or angels for i[n]stance the angel that went forth to distroyed a hundred thousand one knight [night]
						\hypertarget{EMS-1841-JS-5-JAN-1841-WM-c}%
					Joseph
						\marginnote{c}%
					Smith said to D Ells, \& to the Congragetion that he for a len[g]th of time, thought on phreknoledgee [phrenology]; \& that he had a revalation. the Lord rebukeing him sharply on Creating such a thing; \& further said there was no reality in such a science but was the workes of the Devil; he also said the Lord had told him that bro.--- Law would do well; he would go \& preach the gospel he also said as for his own knowledge the Earth was \sout{make} made out of sumthng for it was imposible for samthng to be made out of nothing fire,--- air, \& watter are eternal existant principles which are the Composition of which the earth--- has been Composed; also this, earth has been organized owt of portions of other globes that has be disorganized; in testimony that this earth was not the first of Gods work;
						\hypertarget{EMS-1841-JS-5-JAN-1841-WM-d}%
					he
						\marginnote{d}%
					quoted a pasage from the testament where Jesus said all things that he had saw the father do he had done \& that he done nothing but what he saw the father do John the 5\supersu{th}\supers{.} he also siad in testimony of the situation the saints in the presence of God. that they had flesh \& bones \& that was the agreement in eter[n]ity to come here \& take on them tabernicles \& the differance between us \& satin [Satan] in that respect is that he fell--- \& had not opertunity to come in the flesh--- \& that he allways is striveing to get others as miserable as himself---
					
			% --
			\paragraph{As Reported by William Clayton}
			% --
				\label{EMS-1841-JS-5-JAN-1841-WC}
			
				\begin{description}
					\item[Print Sources]
				\end{description}

				\begin{tabular}{p{0.5in}p{1in}p{0.5in}p{1in}}
					JSP	 	& ---		& JSR 		& --- \\
					DC	 	& ---		& HC 		& --- \\
					HDDC 	& ---		& TCHC 		& --- \\
				\end{tabular}
				% HDDC = woodford's DC diss; JSR = Marquardt's JS Revelations; TCHC = Vogel HC
				
				\begin{description}
					\item[Online Access] \href{https://www.josephsmithpapers.org/paper-summary/discourse-5-january-1841-as-reported-by-william-clayton/1}{Here}
					% \item[Analysis] \hyperref[XX]{Here}
				\end{description}
				
				\begin{description}
					\item[Background]
				\end{description}
				
					% It might also be useful to give a two sentence context of the period: e.g., "This speech was given in Nauvoo to a small congregation one month prior to the destruction of the Nauvoo Expositor. These and other tensions may have been on the prophet's mind when he gave the speech."
				
				\begin{description}
					\item[Original Source]
				\end{description}
				
					% brief description of original source material, with reference to JSP or somewhere else for more info
					% Background material: it might be helpful to have a note that says whether a given document was presented as a speech, was written in a journal, etc.
					% Also a comment here about how reliable the source might be, or what shape it is in, if there are large omissions or parts that are hard to understand, and so on.
				
				\begin{description}
					\item[Text]
				\end{description}
				
					\begin{tightcenter}
							\hypertarget{EMS-1841-JS-5-JAN-1841-WC-a}%
						By
							\marginnote{a}%
						Joseph, Jany. 5\supersu{th}, 1841, at the organization of a school of instruction.
					\end{tightcenter}

					Description of Paul--- He is about 5 foot high; very dark hair; dark complection; dark skin; large Roman nose; sharp face; small black eyes, penetrating as eternity; round shoulders; a whining voice, exept when elevated and then it almost resembles the roaring of a Lion. He was a good orator, but Doctor [John C.] Bennett is a superior orator, and like Paul is active and diligent, always, employing himself in doing good to his fellow men.

						\hypertarget{EMS-1841-JS-5-JAN-1841-WC-b}%
					By
						\marginnote{b}%
					Joseph, January 5\supersu{th}, 1841. Answer to the question, was the Priesthood of Melchizedeck taken away when Moses died.

					All priesthood is Melchizedeck; but there are different portions or degrees of it. That portion which brought Moses to speak with God face to face was taken away; but that which brought the ministry of angels remained. All the Prophets had the Melchizedeck Priesthood and was ordained by God himself.

						\hypertarget{EMS-1841-JS-5-JAN-1841-WC-c}%
					The
						\marginnote{c}%
					world and earth are not synonymous terms. The world%
						\footnote{%
							% \label{}%
							Compare Christ's remarks about the world in \hyperref[JOHN-15-18]{John 15:18ff}.
							}
					is the human family. This earth was organized or formed out of other planets which were broke up and remodelled and made into the one on which we live.
						\hypertarget{EMS-1841-JS-5-JAN-1841-WC-d}%
					The
						\marginnote{d}%
					elements are eternal. That which has a beginning will surely have an end.

					Take a ring, it is without beginning or end; cut it for a beginning place, and at the same time you \sout{will} have an ending place.

					\uline{A key}, every principle proceeding from God is eternal, and any principle which is not eternal is of the Devil. The sun has no beginning or end, the rays which proceed from himself have no bounds, consequently are eternal. So it is with God. If the soul of man had a beginning it will surely have an end.
						\hypertarget{EMS-1841-JS-5-JAN-1841-WC-e}%
					In
						\marginnote{e}%
					the translation; ``without form and void'' it should read ``empty and desolate'' The word ``created'' should be formed or organized.

					\begin{tightcenter}
							\hypertarget{EMS-1841-JS-5-JAN-1841-WC-f}%
						Observations
							\marginnote{f}%
						on the Sectarian God.
					\end{tightcenter}

					That which is without body or parts is nothing. There is no other God in heaven but that God who has flesh and bones. John 5--- 26, ``As the father hath life in himself, even so hath he given the son to have life in himself.'' God the father took life unto himself precisely as Jesus did. The first step in the salvation of men is the laws of eternal and self-existent principles. Spirits are eternal.
						\hypertarget{EMS-1841-JS-5-JAN-1841-WC-g}%
					At
						\marginnote{g}%
					the first organization in heaven we were all present and saw the Savior chosen and appointed, and the plan of salvation made and we sanctioned it. We came to this earth that we might have a body and present it pure before God in the Celestial Kingdom. The great principle of happiness consists in having a body. The Devil has no body, and herein is his punishment. He is pleased when he can obtain the tabernacle of man and when cast out by the Savior he asked to go into the herd of swine showing that he would prefer a swines body to having none. All beings who have bodies have power over those who have not. The devil has no power over us only as we permit him; the moment we revolt at anything which comes from God the Devil takes power.

						\hypertarget{EMS-1841-JS-5-JAN-1841-WC-h}%
					This
						\marginnote{h}%
					earth will be rolled back into the presence of God and crowned with Celestial Glory.
					
			% --
			\paragraph{As Reported by Unidentified Scribe}
			% --
				\label{EMS-1841-JS-5-JAN-1841-US}
			
				\begin{description}
					\item[Print Sources]
				\end{description}

				\begin{tabular}{p{0.5in}p{1in}p{0.5in}p{1in}}
					JSP	 	& ---		& JSR 		& --- \\
					DC	 	& ---		& HC 		& --- \\
					HDDC 	& ---		& TCHC 		& --- \\
				\end{tabular}
				% HDDC = woodford's DC diss; JSR = Marquardt's JS Revelations; TCHC = Vogel HC
				
				\begin{description}
					\item[Online Access] \href{https://www.josephsmithpapers.org/paper-summary/discourse-5-january-1841-as-reported-by-unidentified-scribe/1}{Here}
					% \item[Analysis] \hyperref[XX]{Here}
				\end{description}
				
				\begin{description}
					\item[Background]
				\end{description}
				
					% It might also be useful to give a two sentence context of the period: e.g., "This speech was given in Nauvoo to a small congregation one month prior to the destruction of the Nauvoo Expositor. These and other tensions may have been on the prophet's mind when he gave the speech."
				
				\begin{description}
					\item[Original Source]
				\end{description}
				
					% brief description of original source material, with reference to JSP or somewhere else for more info
					% Background material: it might be helpful to have a note that says whether a given document was presented as a speech, was written in a journal, etc.
					% Also a comment here about how reliable the source might be, or what shape it is in, if there are large omissions or parts that are hard to understand, and so on.
				
				\begin{description}
					\item[Text]
				\end{description}
				
						\hypertarget{EMS-1841-JS-5-JAN-1841-US-a}%
					\phantom{ }[\textit{previous page(s) missing}]
						\marginnote{a}%
					have an ending place, A. key. Every principle proceeding from God is eternal; and any principle which is not etenal is of the Devil. The sun has no beginning or end, the rays which proceed from himself have no bounds; consequently are eternal; so it is with God, If the soul of man had a beginning it will sureley have an end. In the translation, ``without form and void'' it should read ``empty \& desolate, The word created should be formed or organized.---
					
					\begin{tightcenter}
							\hypertarget{EMS-1841-JS-5-JAN-1841-US-b}%
						Observations
							\marginnote{b}%
						on the sectarian God---
					\end{tightcenter}
					
					That which is without body or parts is nothing. There is no other God in heaven but that God who has flesh \& bones. John 5--- 26 as the Father hath \sout{lives} <life> in himself even so hath he given the son to have life in himself. God the Father took life unto himself, precisely as Jesus did. The first step in the salvation of man, is the laws of eternal and selfexistent principles. Spirits are eternal.
						\hypertarget{EMS-1841-JS-5-JAN-1841-US-c}%
					At
						\marginnote{c}%
					the first organization in heaven we were all present \& saw the Saivior, chosen \& appointed, and the place of salvation made and we san[c]tioned it. We came to this earth that we might have a body, \& present it pure before \sout{before} God in the celestial kingdom. The great principle of happiness consists in having a body. The Devil has no body, \& herein is his punishment he is pleased when he can obtain the tabernacle of man; \& when cast out by the saiviour he asked to go into the herd of swine, showing that he would prefer a swines body to haveing none All beings who have bodies have power over these who have not. The Devil has no power over \sout{those} us only as we permit him. The moment we revolt at any thing which comes from God, the Devil takes power.---
						\hypertarget{EMS-1841-JS-5-JAN-1841-US-d}%
					This
						\marginnote{d}%
					earth will be roled back into the presence of God \& crowned with celestial glory---
					
		% -----
		\subsubsection{Account of Meeting, 12 January 1841}
		% -----
			\label{EMS-1841-JS-12-JAN-1841}
			% \label{EMS-1840-JS-DAY-3LETTERMONTH-YEAR}
			
			% --
			\paragraph{As Reported by William P. McIntire}
			% --
				\label{EMS-1841-JS-12-JAN-1841-WM}
			
				\begin{description}
					\item[Print Sources]
				\end{description}

				\begin{tabular}{p{0.5in}p{1in}p{0.5in}p{1in}}
					JSP	 	& ---		& JSR 		& --- \\
					DC	 	& ---		& HC 		& --- \\
					HDDC 	& ---		& TCHC 		& --- \\
				\end{tabular}
				% HDDC = woodford's DC diss; JSR = Marquardt's JS Revelations; TCHC = Vogel HC
				
				\begin{description}
					\item[Online Access] \href{https://www.josephsmithpapers.org/paper-summary/account-of-meeting-12-january-1841-as-reported-by-william-p-mcintire/1}{Here}
					% \item[Analysis] \hyperref[XX]{Here}
				\end{description}
				
				\begin{description}
					\item[Background]
				\end{description}
				
					% It might also be useful to give a two sentence context of the period: e.g., "This speech was given in Nauvoo to a small congregation one month prior to the destruction of the Nauvoo Expositor. These and other tensions may have been on the prophet's mind when he gave the speech."
				
				\begin{description}
					\item[Original Source]
				\end{description}
				
					% brief description of original source material, with reference to JSP or somewhere else for more info
					% Background material: it might be helpful to have a note that says whether a given document was presented as a speech, was written in a journal, etc.
					% Also a comment here about how reliable the source might be, or what shape it is in, if there are large omissions or parts that are hard to understand, and so on.
				
				\begin{description}
					\item[Text]
				\end{description}
				
						\hypertarget{EMS-1841-JS-12-JAN-1841-WM-a}%
					Tuesday
						\marginnote{a}%
					the \sout{19} <12>\supers{th} at M\supersu{r}\supers{.} Davises---

					1\supersu{st}\supers{.} subject education by M\supersu{r}\supers{.} [\textit{blank}] he said a people whose minds were Culavated [cultivated] \& Man[n]ers Refined By Education that they had Great \& precious E[n]joyment that Ignorant had Not

					2\supersu{d}\supers{.} subject was Vice by Mr. stout [Hosea Stout] he said Murder theft \& the Like come from indulgances in this principle---

						\hypertarget{EMS-1841-JS-12-JAN-1841-WM-b}%
					Joseph---
						\marginnote{b}%
					said that some things ware Eavils that did not come under the Head of vise; for instance one Nation would Come against [an]other it would be an Eavil on the Nation was Come aganst yet it would not be a vice for them to Repell---

					Virtue was 3\supers{d} subject

						\hypertarget{EMS-1841-JS-12-JAN-1841-WM-c}%
					4\supers{th.}
						\marginnote{c}%
					subject the Gospel by M\supers{r.} Badlam [Alexander Badlam Sr.] he lectors on it till he Comes the Laying on the hands for the Holy Ghost--- then Joseph--- takes it up \& ad[d]s the Resurection \& Eternal Judgment in the Eternal Judgment there is many things to know \& to under stand in Gods Judging for instance Peter said David\sout{d} had not yet ascend to heaven \& that he was a Murderer--- \& that his soul was in Hell is plainly told By Peter in acts--- 2\supersu{d}\supers{.} ch--- Petter shews plainer [i]t in the 3\supers{d} of acts that a Murderer could Not be Red [ee]med intill he would send Jesus Christ which before was preached unto you \&c--- that is that faith Repentance \& Baptism would not save them untill the[y] ware scourged in hell or paid the Last farthing---
					
		% -----
		\subsubsection{Proclamation, 15 January 1841}
		% -----
			\label{EMS-1841-JS-15-JAN-1841}
			% \label{EMS-1840-JS-DAY-3LETTERMONTH-YEAR}
			
			\begin{description}
				\item[Print Sources]
			\end{description}

			\begin{tabular}{p{0.5in}p{1in}p{0.5in}p{1in}}
				JSP	 	& ---		& JSR 		& --- \\
				DC	 	& ---		& HC 		& --- \\
				HDDC 	& ---		& TCHC 		& --- \\
			\end{tabular}
			% HDDC = woodford's DC diss; JSR = Marquardt's JS Revelations; TCHC = Vogel HC
			
			\begin{description}
				\item[Online Access] \href{https://www.josephsmithpapers.org/paper-summary/proclamation-15-january-1841/1}{Here}
				% \item[Analysis] \hyperref[XX]{Here}
			\end{description}
			
			\begin{description}
				\item[Background]
			\end{description}
			
				% It might also be useful to give a two sentence context of the period: e.g., "This speech was given in Nauvoo to a small congregation one month prior to the destruction of the Nauvoo Expositor. These and other tensions may have been on the prophet's mind when he gave the speech."
			
			\begin{description}
				\item[Original Source]
			\end{description}
			
				% brief description of original source material, with reference to JSP or somewhere else for more info
				% Background material: it might be helpful to have a note that says whether a given document was presented as a speech, was written in a journal, etc.
				% Also a comment here about how reliable the source might be, or what shape it is in, if there are large omissions or parts that are hard to understand, and so on.
			
			\begin{description}
				\item[Text]
			\end{description}
			
				\begin{tightcenter}
						\hypertarget{EMS-1841-JS-15-JAN-1841-a}%
					A
						\marginnote{a}%
					PROCLAMATION, TO THE SAINTS SCATTERED ABROAD;
				\end{tightcenter}

				\textsc{Greeting:}

				\textit{Beloved Brethren}:---

				The relationship which we sustain to the Church of Jesus Christ of Latter Day Saints, renders it necessary that we should make known from time to time, the circumstances, situation, and prospects of the church, and give such instructions as may be necessary for the well being of the Saints. and for the promotion of those objects, calculated to further their present and everlasting happiness.

				We have to congratulate the Saints on the progress of the great work of the ``last days;'' for not only has it spread through the length and breadth of this vast continent; but on the continent of Europe, and on the Islands of the sea, it is spreading in a manner entirely unprecedented in the annals of time.

					\hypertarget{EMS-1841-JS-15-JAN-1841-b}%
				This
					\marginnote{b}%
				appears the more pleasing when we consider, that but a short time has elapsed, since we were unmercifully driven from the State of Missouri, after suffering cruelties and persecutions in their various, and horrid forms.--- Then our overthrow, to many, seemed inevitable, while the enemies of truth triumphed over us, and by their cruel reproaches endeavored to aggravate our sufferings. But ``the Lord of Hosts was with us, the God of Jacob was our refuge!'' and we were delivered from the hands of bloody and deceitful men; and in the State of Illinois we found an asylum, and were kindly welcomed by persons worthy the characters of \textsc{freemen}.
					\hypertarget{EMS-1841-JS-15-JAN-1841-c}%
				It
					\marginnote{c}%
				would be impossible to enumerate all those who in our time of deep distress, nobly came forward to our relief, and like the good Samaritan poured oil into our wounds, and contributed liberally to our necessities, as the citizens of Quincy \textit{en masse} and the people of Illinois, generally, seemed to emulate each other in this labor of love. We would, however, make honorable mention of Governor [Thomas] Carlin, Judge [Richard M.] Young, General [Samuel] Leech, Judge [James] Ralston, Rev. Mr. Young, Col. Henry, N[ehemiah] Bushnell, John Wood, I[saac] N. Morris, S[ylvester] M. Bartlett, Samuel Holmes, and J[oseph] T. Holmes, Esquires, who will long be remembered by a grateful community for their philanthropy to a suffering people, and whose kindness on that occasion is indelibly engraven on the tablet of our hearts, in golden letters of love.

					\hypertarget{EMS-1841-JS-15-JAN-1841-d}%
				We
					\marginnote{d}%
				would, likewise, make mention of the Legislature of this State, who, \textit{without respect of parties, without reluctance, freely, openly, boldly, and nobly,} have come forth to our assistance, owned us as citizens and friends, and took us by the hand, and extended to us all the blessings of civil, political, and religious liberty, by granting us, under date of Dec. 16, 1840, one of the most liberal charters, with the most plenary powers, ever conferred by a legislative assembly on free citizens, for the ``City of Nauvoo,'' the ``Nauvoo Legion'' and the ``University of the City of Nauvoo.'' The first of these charters, (that for the ``City of Nauvoo,'') secures to us in all time to come, irrevocably, all those great blessings of civil liberty, which of right appertain to all the free citizens of a great civilized republic---'tis all we ever claimed.
					\hypertarget{EMS-1841-JS-15-JAN-1841-e}%
				What
					\marginnote{e}%
				a contrast does the proceedings of the legislature of this State present, when compared with those of Missouri, whose bigotry, jealousy, and superstition, prevailed to such an extent, as to deny us our liberty and our sacred rights---Illinois has set a glorious example, to the whole United States and to the world at large, and has nobly carried out the principles of her constitution, and the constitution of these United States, and while she requires of us implicit obedience to the laws, (which we hope ever to see observed) she affords us the protection of law---the security of life, liberty, and the peaceable pursuit of happiness.

					\hypertarget{EMS-1841-JS-15-JAN-1841-f}%
				The
					\marginnote{f}%
				name of our city (Nauvoo,) is of Hebrew origin, and signifies a beautiful situation, or place, carrying with it, also, the idea of \textit{rest}; and is truly descriptive of this most delightful situation. It is situated on the eastern bank of the Mississippi river, at the head of the Des Moines Rapids, in Hancock County; bounded on the east by an extensive prairie of surpassing beauty, and on the north, west, and south, by the Mississippi. This place has been objected to by some, on account of the sickness which has prevailed in the summer months, but it is the opinion of Doctor [John C.] Bennett, a physician of great experience and medical knowledge, that Hancock Co., and all the eastern and southern portions of the City of Nauvoo, are as healthy as any other portions of the western country, (or the world, to acclimated citizens,) whilst the northwestern portion of the city has experienced much affliction from ague and fever, which, however, he thinks can be easily remedied by draining the sloughs on the adjacent islands in the Mississippi.

					\hypertarget{EMS-1841-JS-15-JAN-1841-g}%
				The
					\marginnote{g}%
				population of our city is increasing with unparralled [unparalleled] rapidity, numbering more than three thousand inhabitants. Every facility is afforded in the city and adjacent country, in Hancock County, for the successful prosecution of the mechanical arts, and the pleasing pursuits of agriculture. The waters of the Mississippi can be successfully used for manufactoring purposes, to an almost unlimited extent.

					\hypertarget{EMS-1841-JS-15-JAN-1841-h}%
				Having
					\marginnote{h}%
				been instrumental in the hands of our heavenly Father in laying a foundation for the gathering of Zion, we would say, let all those who appreciate the blessings of the gospel, and realize the importance of obeying the commandments of heaven, who have been blessed of heaven with the possession of this world's goods, first prepare for the general gathering---let them dispose of their effects as fast as circumstances will possibly admit, without making too great sacrifices, and remove to our city and county---establish and build up manufactories in the city, purchase and cultivate farms in the county---this will secure our permanent inheritance, and prepare the way for the gathering of the poor. 
					\hypertarget{EMS-1841-JS-15-JAN-1841-i}%
				\phantom{ }
					\marginnote{i}%
				\textit{This is agreeable to the order of heaven, and the only principal on which the gathering can be effected}---let the rich, then, and all who can assist in establishing this place, make every preparation to come on without delay, and strengthen our hands, and assist in promoting the happiness of the Saints. This cannot be too forcibly impressed on the minds of all, and the elders are hereby instructed to proclaim this word in all places where the Saints reside, in their public administrations, for this is according to the instructions we have received from the Lord.

					\hypertarget{EMS-1841-JS-15-JAN-1841-j}%
				The
					\marginnote{j}%
				Temple of the Lord is in progress of erection here, where the Saints will come to worship the God of their fathers, according to the order of his house, and the powers of the holy priesthood, and will be so constructed as to enable all the functions of the priesthood to be duly exercised, and where instructions from the Most High will be received, and from this place go forth to distant lands.

				Let us then concentrate all our powers, under the provisions of our \textit{magna charta} granted by the Illinois Legislature, at the ``City of Nauvoo,'' and surrounding country, and strive to emulate the actions of the ancient covenant fathers, and patriarchs, in those things, which are of such vast importance to this and every succeeding generation.

					\hypertarget{EMS-1841-JS-15-JAN-1841-k}%
				The
					\marginnote{k}%
				``Nauvoo Legion,'' embraces all our military power, and will enable us to perform our military duty by ourselves, and thus afford us the power, and privilege, of avoiding one of the most fruitful sources of strife, oppression, and collision with the world. It will enable us to show our attachment to the state and nation as a people, whenever the public service requires our aid---thus proving ourselves obedient to the paramount laws of the land, and ready at all times to sustain and execute them.

					\hypertarget{EMS-1841-JS-15-JAN-1841-l}%
				The
					\marginnote{l}%
				``University of the City of Nauvoo,'' will enable us to teach our children wisdom---to instruct them in all knowledge, and learning, in the Arts, Sciences and Learned Professions. We hope to make this institution one of the great lights of the world, and by and through it, to diffuse that kind of knowledge which will be of practical utility, and for the public good, and also for private and individual happiness. The Regents of the University will take the general supervision of all mat[t]ers appertaining to education from common schools up to the highest branches of a most liberal collegiate course. They will establish a regular system of education, and hand over the pupil from teacher to professor, until the regular gradation is consummated, and the education finished. This corporation contains all the powers and perogatives of any other college or university in this state. The charters for the University and Legion are \textit{addenda} to the city charter, making the whole perfect and complete.

					\hypertarget{EMS-1841-JS-15-JAN-1841-m}%
				Not
					\marginnote{m}%
				only has the Lord given us favor in the eyes of the community, who are happy to see us in the enjoyment of all the rights and privileges of freemen, but we are happy to state that several of the principal men of Illinois, who have listened to the doctrines we promulge, have become obedient to the faith and are rejoicing in the same; among whom is John C. Bennett, M. D., Quarter Master General of Illinois. We mention this gentleman first, because, that during our persecutions in Missouri, he became acquainted with the violence we were suffering, while in that State, on account of our religion---his sympathies for us were aroused, and his indignation kindled against our persecutors for the cruelties practised upon us, and their flagrant violation of both the law and the constitution.
					\hypertarget{EMS-1841-JS-15-JAN-1841-n}%
				Amidst
					\marginnote{n}%
				their heated zeal to put down the truth, he addressed us a letter, tendering to us his assistence in delivering us out of the hands of our enemies, and restoring us again to our privileges, and only required at our hands to point out the way, and he would be forthcoming, with all the forces he could raise for that purpose---He has been one of the principal instruments, in effecting our safety and deliverance from the unjust persecutions and demands of the authorities of Missouri, and also in procuring the city charter---He is a man of enterprize, extensive acquirements, and of independant mind, and is calculated to be a great blessing to our community.

					\hypertarget{EMS-1841-JS-15-JAN-1841-o}%
				Dr.
					\marginnote{o}%
				Isaac Galland, also, who is one of our benefactors, having under his control, a large quantity of land in the immediate vicinity of our city, and a considerable portion of the city plot opened both his heart and his hands, and ``when we were strangers---took us in,'' and bade us welcome to share with him in his abundance; leaving his dwelling house, the most splendid edifice in the vicinity, for our accommodation, and betook himself to a small, uncomfortable dwelling---He sold us his large estates, on very reasonable terms, and on long credit, so that we might have an opportunity of paying for them, without being distressed, and has since taken our lands in Missouri in payment for the whole amount, and has given us a clear and indisputable title for the same.
					\hypertarget{EMS-1841-JS-15-JAN-1841-p}%
				And
					\marginnote{p}%
				in addition to the first purchase, we have exchanged lands with him in Missouri to the amonnt of eighty thousand dollars. He is the honored instrument the Lord used, to prepare a home for us, when we were driven from our inheritances, having given him control of vast bodies of land, and prepared his heart to make the use of it the Lord intended he should. Being a man of extensive information, great talents, and high literary fame, he devoted all his powers and influence to give us a character.

					\hypertarget{EMS-1841-JS-15-JAN-1841-q}%
				After
					\marginnote{q}%
				having thus exerted himself for our salvation and comfort, and formed an intimate acquaintance with many of our people, his mind became wrought up to the greatest feelings, being convinced that our persecutions, were like those of the ancient Saints, and after investigating the doctrines we proclaimed, he became convinced of the truth and of the necessity of obedience thereto, and to the great joy and satisfaction of the church he yielded himself to the waters of baptism, and became a partaker with us in our sufferings. ``choosing rather to suffer afflictions with the people of God than enjoy the pleasures of sin for a season.''
					\hypertarget{EMS-1841-JS-15-JAN-1841-r}%
				In
					\marginnote{r}%
				connexion with these, we would mention the names of Gen. James Adams, Judge of Probate, of Sangamon County, Dr. Green, of Shelby County, R[obert] D. Foster, M. D., a gentleman of great energy of character, late of Adams Co., Sidney Knowlton, of Hancock Co., Dr. [Lenox] Knight, of Putnam County, Indiana, with many others of respectability and high standing in society, with nearly all the old settlers in our immediate neighborhood. We make mention of this, that the Saints may be encouraged, and also that they may see the persecutions we suffered in Missouri, were but the prelude to a far more glorious display of the power of truth, and of the religion we have espoused.

					\hypertarget{EMS-1841-JS-15-JAN-1841-s}%
				From
					\marginnote{s}%
				the kind, uniform, and consistent course pursued by the citizens of Illinois, and the great success which has attended us while here, the natural advantages of this place for every purpose we require, and the necessity of the gathering of the Saints of the Most High, we would say, let the brethren who love the prosperity of Zion, who are anxious that her stakes should be strengthened, and her cords lengthened, and who prefer her prosperity to their chief joy, come, and cast in their lots with us, and cheerfully engage in a work so glorious and sublime, and say with Nehemiah, ``we his servants will arise and build.''

					\hypertarget{EMS-1841-JS-15-JAN-1841-t}%
				It
					\marginnote{t}%
				probably would hardly be necessary to enforce this important subject on the attention of the Saints, as its necessity is obvious, and is a subject of paramount importance; but as watchmen to the house of Israel, as Shepherds over the flock which is now scattered over a vast extent of country, and the anxiety we feel for their prosperity and everlasting welfare, and for the carrying out the great and glorious purposes of our God, to which we have been called, we feel to urge its necessity, and say, let the Saints come \textit{here}---\textsc{This is the word of the Lord}, \textit{and in accordance with the great work of the last days}.

					\hypertarget{EMS-1841-JS-15-JAN-1841-u}%
				It
					\marginnote{u}%
				is true the idea of a general gathering has heretofore been associated with most cruel and oppressing scenes, owing to our unrelenting persecutions at the hands of wicked and unjust men; but we hope that those days of darkness and gloom have gone by, and from the liberal policy of our State government, we may expect a scene of peace and prosperity, we have never before witnessed since the rise of our church, and the happiness and prosperity which now await us, is, in all human probablility, incalculably great.
					\hypertarget{EMS-1841-JS-15-JAN-1841-v}%
				By
					\marginnote{v}%
				a concentration of action, and a unity of effort, we can only accomplish the great work of the last days, which we could not do in our remote and scattered condition, while our interests both spiritual and temporal will be greatly enhanced, and the blessings of heaven must flow unto us in an uninterrupted stream; of this, we think there can be no question. The great profusion of temporal and spiritual blessings, which always flow from faithfulness and concerted effort, never attend individual exertion or enterprize. The history of all past ages abundantly attests this fact.

					\hypertarget{EMS-1841-JS-15-JAN-1841-w}%
				In
					\marginnote{w}%
				addition to all temporal blessings, there is no other way for the Saints to be saved in these last days, as the concurrent testimony of all the holy prophets clearly proves, for it is written---``They shall come from the east and be gathered the west; the north shall give up, and the south shall keep not back''---``the sons of God shall be \textit{gathered} from far, and his daughters from the ends of the earth:'' it is also the concurrent testimony of all the prophets, that this gathering together of all the Saints, must take place before the Lord comes to ``take vengeance upon the ungodly,'' and ``to be glorified and admired by all those who obey his gospel.'' The 50 Psalm from the first to the fifth verses, inclusive, describes the glory and majesty of that event. ``The mighty God even the Lord hath spoken and called the earth from the rising of the sun unto the going down thereof.--- Out of Zion, the perfection of beauty, God hath shined.

					\hypertarget{EMS-1841-JS-15-JAN-1841-x}%
				Our
					\marginnote{x}%
				God shall come, and shall not keep silence: a fire shall devour before him and it shall be very tempestous round about him.

				He shall call to the heavens from above, and to the earth, (that he may judge his people.)

				Gather my Saints together unto me; those that have made a covenant with me by sacrifice.''

				We might offer many other quotations from the scriptures, but believing them to be familiar to the Saints we forbear.

					\hypertarget{EMS-1841-JS-15-JAN-1841-y}%
				We
					\marginnote{y}%
				would wish the Saints to understand that, when they come here they must not expect to find perfection, or that all will be harmony, peace and love; if they indulge these ideas, they will undoubtedly be deceived for here there are persons, not only from different States, but from different nations, who, although they feel a great attachment to the cause of truth, have their predjudices of ed[u]cation, and consequently it requires some time before these things can be overcome: again, there are many that creep in unawares, and endeavor to sow discord, strife and animosity, in our midst, and by so doing bring evil upon the Saints;
					\hypertarget{EMS-1841-JS-15-JAN-1841-z}%
				these
					\marginnote{z}%
				things we have to bear with, and these things will prevail either to a greater or lesser extent until ``the floor be thoroughly purged'' and ``the chaff be burnt up.'' Therefore let those who come up to this place, be determined to keep the commandments of God, and not be discouraged by those things we have enumerated, and then they will be prospered, the intelligence of heaven will be communicated to them, and they will eventually see eye to eye, and rejoice in the full fruition of that glory, which is reserved for the righteous.

					\hypertarget{EMS-1841-JS-15-JAN-1841-aa}%
				In
					\marginnote{aa}%
				order to erect the Temple of the Lord, great exertions will be required on the part of the Saints, so that they may build a house which shall be accepted of by the Almighty, and in which his power and glory shall be manifested. Therefore let those who can, freely make a sacrifice of their time, their talents, and their property, for the prosperity of the kingdom, and for the love they have to the cause of truth, bid adieu to their homes and pleasant places of abode, and unite with us in the great work of the last days, and share in the tribulation, that they may ultimately share in the glory and triumph.

					\hypertarget{EMS-1841-JS-15-JAN-1841-bb}%
				We
					\marginnote{bb}%
				wish it, likewise, to be distinctly uderstood that we claim no privilege but what we feel cheerfully disposed to share with our fellow citizens of every denomination, and every sentiment of religion; and therefore say, that, so far from being restricted to our own faith, let all those who desire to locate themselves in this place, or the vicinity, come, and we will hail them as citizens and friends, and shall feel it not only a duty, but a privilege, to reciprocate the kindness we have received from the benevolent and kind hearted citizens of the State of Illinois.

				\begin{tightright}
					JOSEPH SMITH,

					SIDNEY RIGDON,

					HYRUM SMITH,

					Presidents of the Church.
				\end{tightright}

				\textit{Nauvoo, January} 15, 1841.
				
		% -----
		\subsubsection{Account of Meeting, circa 19 January 1841}
		% -----
			\label{EMS-1841-JS-CIR-19-JAN-1841}
			% \label{EMS-1840-JS-DAY-3LETTERMONTH-YEAR}
			
			% --
			\paragraph{As Reported by William P. McIntire}
			% --
				\label{EMS-1841-JS-CIR-19-JAN-1841-WM}
			
				\begin{description}
					\item[Print Sources]
				\end{description}

				\begin{tabular}{p{0.5in}p{1in}p{0.5in}p{1in}}
					JSP	 	& ---		& JSR 		& --- \\
					DC	 	& ---		& HC 		& --- \\
					HDDC 	& ---		& TCHC 		& --- \\
				\end{tabular}
				% HDDC = woodford's DC diss; JSR = Marquardt's JS Revelations; TCHC = Vogel HC
				
				\begin{description}
					\item[Online Access] \href{https://www.josephsmithpapers.org/paper-summary/account-of-meeting-circa-19-january-1841-as-reported-by-william-p-mcintire/1}{Here}
					% \item[Analysis] \hyperref[XX]{Here}
				\end{description}
				
				\begin{description}
					\item[Background]
				\end{description}
				
					% It might also be useful to give a two sentence context of the period: e.g., "This speech was given in Nauvoo to a small congregation one month prior to the destruction of the Nauvoo Expositor. These and other tensions may have been on the prophet's mind when he gave the speech."
				
				\begin{description}
					\item[Original Source]
				\end{description}
				
					% brief description of original source material, with reference to JSP or somewhere else for more info
					% Background material: it might be helpful to have a note that says whether a given document was presented as a speech, was written in a journal, etc.
					% Also a comment here about how reliable the source might be, or what shape it is in, if there are large omissions or parts that are hard to understand, and so on.
				
				\begin{description}
					\item[Text]
				\end{description}
				
					Next Meeting--- Joseph said that before foundation of the Earth in the Grand-Counsel that the spirits of all \uline{men} ware subject to opression \& the express purpose of God in Giveing it a tabernicle was to arm it against the power of Darkness--- for instance Jesus said get behind me Satan also the apostle said Resist the Devil \& he will flee from you
					
		% -----
		\subsubsection{Revelation, 19 January 1841 [D\&C 124]}
		% -----
			\label{EMS-1841-JS-19-JAN-1841}
			% \label{EMS-1840-JS-DAY-3LETTERMONTH-YEAR}
			
			\begin{description}
				\item[Print Sources]
			\end{description}

			\begin{tabular}{p{0.5in}p{1in}p{0.5in}p{1in}}
				JSP	 	& ---		& JSR 		& --- \\
				DC	 	& ---		& HC 		& --- \\
				HDDC 	& ---		& TCHC 		& --- \\
			\end{tabular}
			% HDDC = woodford's DC diss; JSR = Marquardt's JS Revelations; TCHC = Vogel HC
			
			\begin{description}
				\item[Online Access] \href{https://www.josephsmithpapers.org/paper-summary/revelation-19-january-1841-dc-124/1}{Here}
				% \item[Analysis] \hyperref[XX]{Here}
			\end{description}
			
			\begin{description}
				\item[Background]
			\end{description}
			
				% It might also be useful to give a two sentence context of the period: e.g., "This speech was given in Nauvoo to a small congregation one month prior to the destruction of the Nauvoo Expositor. These and other tensions may have been on the prophet's mind when he gave the speech."
			
			\begin{description}
				\item[Original Source]
			\end{description}
			
				% brief description of original source material, with reference to JSP or somewhere else for more info
				% Background material: it might be helpful to have a note that says whether a given document was presented as a speech, was written in a journal, etc.
				% Also a comment here about how reliable the source might be, or what shape it is in, if there are large omissions or parts that are hard to understand, and so on.
			
			\begin{description}
				\item[Text]
			\end{description}
			
				A revelation given to Joseph Smith, January 19\supers{th.} 1841

				Verily thus saith the Lord unto you my Servant Joseph Smith, I am well pleased with your offering, and acknowledgements which you have made; for unto this end have I raised you up, that I might shew forth my wisdom through the weak things of the earth.

				Your prayers are acceptable before me, and in answer to them, I say unto you, that you are now called, immediately to make a solemn proclamation of my gospel, and of this stake, which I have planted to be a corner stone of Zion, which shall be polished with that refinement which is after the similitude of a palace. This proclamation shall be made to all the Kings of the world, to the four corners thereof; to the Honorable President Elect, and the high minded Governors of the nation in which you live, and to all the nations of the earth, scattered abroad. Let it be written in the spirit of meekness and by the power of the holy ghost, which shall be in you, at the time of the writing of the same; for it shall be given you by the holy ghost to know my will concerning those Kings and Authorities, even what shall befall them in a time to come. For behold I am about to call upon them to give heed to the light \sout{of} <and> glory of Zion, for the set time has come, to favor her.

				Call ye, therefore, upon them with loud proclamation and with your testimony, fearing them not, for they are as grass, and all their glory as the flower thereof, which soon falleth away; that they may be left also without excuse, and that I may visit them in the day of visitation when I shall unveil the face of my covering, to appoint the portion of the oppressor, among hypocrites, where there is gnashing of teeth; if they reject my servants, and my testimony, which I have revealed unto them. And again, I will visit and soften their hearts--- many of them for your good, that ye may find grace in their eyes, that they may come to the light of truth, and the gentiles to the exaltation or lifting up of Zion; for the day of my visitation cometh speedily, in an hour when you think not of; and where shall be the safety of my people? and refuge for those who shall be left of them?

				Awake! O Kings of the earth! Come ye, O! come ye with your gold and your silver, to the help of my people, to the house of the daughter of Zion.

				And, again, verily I say unto you, let my servant, Robert Blashel Thompson, help you to write this proclamation, for I am well pleased with him, and that he should be with you, let him therefore hearken to your \sout{council} counsel and I will bless him with a multiplicity of blessings, let him be faithful and true in all things from henceforth and he shall be great in mine eyes; but let him remember that his stewardship will I require at his hands.

				And again, verily, I say unto you, blessed is my servant Hyrum Smith, for I the Lord loveth him, because of the integrity of his heart, and because, he loveth that which is right before me saith the Lord.

				Again, let my servant John C. Bennett help you in your labor, in sending my word to the Kings and peoples of the Earth, and stand by you, even you, my servant Joseph Smith in the hour of affliction, and his reward shall not fail, if he receive \sout{council} counsel, and for his love, he shall be great, for he shall be mine if he do this, saith the Lord. I have seen the work he hath done, which I accept, if he continue, and will crown him with blessings and great glory.

				And, again, I say unto you, that it is my will that my servant Lyman Wight, should continue in preaching for Zion, in the spirit of meekness, confessing me before the world, and I will bear him up, as on eagle's wings, and he shall beget glory and honor to himself, and unto my name, that when he shall finish his work, that I may receive him unto myself, even as I did my servant David [W.] Patten, who is with me at this time, and also my servant Edward Partridge, and also, my aged servant Joseph Smith Sen\supersu{r}\supers{.}, who sitteth with Abraham at his right hand, and blessed and holy is he, for he is mine.

				And, again, verily I say unto you, my servant George Miller is without guile, he may be trusted because of the integrity of his heart; and <for> the love he has to my testimony, I the Lord loveth him. I therefore say unto you I seal upon his head the office of a ``bishoprick'' like unto my servant Edward Partridge, that he may receive the consecrations of mine house, that he may administer blessings upon the heads of the poor of my people saith the Lord; let no man despise my servant George for he shall honor me. Let my servant George, and my servant Lyman, and my servant John Snider, and others, build a house unto my name, such an one as my servant Joseph shall shew unto them, upon the place which he shall shew unto them also; and it shall be for a house for boarding; a house that strangers may come from afar to lodge therein--- therefore let it be a good house, worthy of all acceptation, that the weary traveller may find health and safety, while he shall contemplate the word of the Lord, and the corner stone I have appointed for Zion. This house shall be a healthy habitation, if it be built unto my name, and if the Governor which shall be appointed unto it, shall not suffer any pollution to come upon it--- It shall be holy, or the Lord your God will not dwell therein.

				And, again, verily I say unto you, let all my saints, from afar; `And send ye swift messengers, yea, chosen messengers and say unto them, Come ye, with all your gold, and your silver, and your preceious stones, and with all your antiquitiees, and with all who have knowledge of antiquities, that will come, may come, and bring the box tree, and the fir tree, and the pine tree, together with all the precious trees of the earth, and with iron, with copper and with brass, and with zink <and with> all your precious things of the earth, and build a house unto \sout{mine} my name, for the Most High to dwell therein, for there is not place found on the earth; that he may come and restore again that which was lost unto you, or which he hath taken away, even the fulness of the Priesthood; for a baptismal font there is not upon the earth; that they, my saints may be baptized for those who are dead, for this ordinance belongeth to my house, and cannot be acceptable to me, only in the days of your poverty, wherein ye are not able to build a house unto me; but I command you, all ye my saints to build a house unto me, and I grant unto you a sufficien[t] time to build a house unto me; and during this time your baptisms, shall be acceptable unto me. But behold, at the end of this appointment, your baptisms for your dead, shall not be acceptable unto me, and if you do not these things, at the end of the appointment, ye shall be rejected as a church, with your dead, saith the Lord your God. For verily I say unto you, that after you have had sufficient time to build a house unto me, wherein the ordinance of baptizing for the dead belongeth, and for which the same was instituted from before the foundation of the world, your baptisms for your dead cannot be acceptable unto me, for therein are the keys of the Holy Priesthood ordained, that you may receive honor and glory. And after this time, your baptisms for the dead, by those who are scattered abroad are not acceptable unto me, saith the Lord; for it is ordained that in Zion and in her stakes, and in Jerusalem those places which I have appointed <for refuge> shall be the places for \sout{your} the baptisms for your dead.

				And, again, verily I say unto you, how shall your washings be acceptable unto me, except, ye perform them in a house which you have built to my name<?>, for, for this cause I commanded Moses that he should build a tabernacle, that they should bear it with them in the wilderness, and to build a house in the land of promise, that those ordinances might be revealed, which had been hid from before the world was; therefore, verily I say unto you, that your anointings and your washings, and your babtisms for the dead and your solemn assemblys, and your memorials for your sacrifices, by the sons of Levi, and for your oracles in your most \uline{holy} places, wherein you receive conversations, and your statutes, and judgements, for the beginning of the revelations and foundation of Zion, and for the glory and honor and endowment of all her municiples, are ordained by the ordinance of my holy house, which my people are allways Commanded to build unto my holy name.

				And verily I say unto you, let this house be built unto my name, that I may reveal \sout{my} <mine> ordinances therein, unto my people, for I deign to reveal unto my church, things which have been kept hid from before the foundation of the world, things that pertain to the dispensation of the fulness of times, and I will shew unto my servant Joseph all things pertaining to this house and the priesthood thereof, and the place whereon it shall be built, and ye shall build it on the place where you have contemplated building <it,> for that is the spot which I have chosen for you to build it. If ye labor with all your mights I will consecrate that spot, that it shall be made holy; and if my people will hearken unto my voice and unto the voice of my servants whom I have appointed to lead my people, behold, verily I say unto you, they shall not be moved out of their place, but if they will not hearken to my voice, nor unto the voice of these men whom I have appointed, they shall not be blest, because they pollute \sout{my} mine holy grounds, and mine holy ordinances, and charters, and my holy words which I give unto them; And it shall come to pass, that if you build a house unto my name, and do not do the things that I say, I will not perform the oath which I make unto you, neither fulfill the promises which ye expect at my hands saith the Lord, for instead of blessings, ye, by your own works, bring cursings, wrath, indignation, and judgments upon your own heads by your follies and by all your abominations which you practise before me saith the Lord. Verily, verily, I say unto you, that when I give a commandment unto any of the sons of men, to do a work unto my name, and those sons of men, go with all their mights, and with all they have, to perform that work and cease not their diligence and their enemies come upon them, and hinder them from performing that work, behold, it behoveth me to require that work no more at the hands of those sons of men; but to accept of their offering. and the iniquity, and transgression of my holy laws and commandments I will visit upon <the heads of> those who hindered my work unto the third and fourth generation, so long; as they <repent not> \& hate me, saith the Lord. God. Therefore for this cause have I accepted of the offerings of those whom I commanded to build up a city, and an house unto my name in Jackson County, Missouri and were hind[e]red by their enemies saith the Lord your God. And I will answer judgment, wrath and i[n]dignation; wailing and anguish, and gnashing of teeth upon \sout{them} their heads unto the third and fourth generation, so long as they repent not, and hate me saith the Lord your God. And this, I make an ensample unto you, for your consolation concerning all those who have been commanded to do a work, and have been hindered by the hands of their enemies, and by oppression, saith the Lord your God; for I am the Lord your God and will save all those <of your brethren> who have been pure in heart, and have been slain in the land of Missouri saith the Lord.

				And again, verily I say unto you, I command you again, to build an house to my name even in this place, that ye may prove yourself unto me, that ye are faithful in all things whatsoever I command you, that I may bless you, and crown you with honor, immortality, and eternal life.

				And now I say unto you, as pertaining to my boarding house, which I have commanded you to build for the boarding of strangers, let it be built unto my name, and let my name, \sout{and} be named upon it, and let my servant Joseph and his house have place therein from generation to generation. For this anointing have I put upon his head, that his blessings shall also be put upon the heads of his posterity after him; and as I said unto Abraham, concerning the kindreds of the earth, even so, I say unto my servant Joseph, in thee, and in thy seed shall the Kindreds of the earth be blessed

				Therefore let my servant Joseph and his seed after him, have place in that house from generation to generation, for ever and ever, saith the Lord, and let the name of that house be called the Nauvoo House, and let it be a delightful habitation for man, and a resting place for the weary traveller, that he may comtemplate the glory of Zion, and the glory of this, the corner stone thereof; that he may receive also, the council from those whom I have set to be as plants of renown, and as watchmen upon her walls.

				Behold, verily I say unto you, let my servant George Miller, and my servant Lyman Wight, and my servant John Snider, and my servant Peter Hawes [Haws], organize themselves, and appoint one of them to be a president over their quorum for the purpose of building that house, and they shall form a constitution whereby they may receive stock, for the building of that house. And they shall not receive less than \uline{fifty} \uline{dollars} for a share of stock in that house, and they shall be permitted to receive fifteen thousand dollars from any one man for stock in that house; but they shall not be permitted to receive over fifteen thousand dollars stock, from any one man; and they shall not be permitted to receive under fifty dollars for a share of stock from any one man in that house, and they shall not be permitted to receive any man as a stockholder in that house, except the same shall pay his stock into their hands at the time he receives stock, and in proportion \sout{to the amount} to the amount of stock he pays into their hands, he shall receive stock in that house; but if he pay nothing into their hands, he shall not receive any stock in that house, and if any man pay stock into their hands it shall be for stock in that house, for himself and for his generation after him, from generation to generation, so long as he, and his heirs shall hold that stock, and do not sell, or convey that stock away out of their hands by their own free will and act--- if you will do my will saith the Lord your God.

				And again verily I say unto you if my servant George Miller, and my servant Lyman Wight, and my servant John Snider and my servant Peter Hawes, receive any stock into their hands, in monies or in \sout{property} properties wherein they receive the real value of monies, they shall not appropriate any portion of that stock to any other purpose only in that house, and if they do appropriate any portion of that stock any where else, only in that house, without the consent of the stockholder, and do not repay four fold for the stock which they appropiate any where else, only in that house, they shall be accursed, and shall be moved out of their place saith the Lord God, for I the Lord am God and cannot be mocked in any of these things.

				Verily: I say unto you let my servant Joseph pay stock into their hands for the building of that house as seemeth him good but my servant Joseph cannot pay over fifteen thousand dollars stock in that house, nor under fifty dollars, neither can any other man saith the Lord.

				And there are others also, who wish to know my will concerning them, for they have asked it at my hands. Therefore I say unto you concerning my servant Vinson Knight, if he will do my will let him put stock into that house for himself and for his generation after him, from generation <to generation> and let him lift up his voice long and loud in the midst of the people to plead the cause of the poor and the needy, and let him not fail neither let his heart faint, and I will accept of his offerings, for they shall not be unto me as the offerings of Cain, for he shall be mine saith the Lord; let his family rejoice and turn away their hearts from affliction, for I have chosen him, and anointed him, and he shall be honored in the midst of his house, for I will forgive all his sins saith the Lord. Amen.

				Verily I say unto you, let my servant Hyrum put stock into that house as seemeth him good for himself and his generation after him, from generation to generation \sout{after him from generation to generation}.

				Let my servant Isaac Galland put stock into that house for I the Lord loveth him for the work he hath done and will forgive all his sins, therefore, let him be remembered for an interest in that house from generation to generation: Let my servant \sout{William Marks} <Isaac Galland> be appointed among you and be ordained by my servant William Marks, and be blest of him to go with my servant Hyrum to accomplish the work that my servant Joseph shall point out unto them and they shall be greatly blessed.

				Let my servant William Marks pay stock into that house as seemeth him good for himself and his generation from generation to generation.

				Let my servant Henry G. Sherwood pay stock into that house as seemeth him good for himself and his seed after him from generation to generation.

				Let my servant William Law pay stock into that house for himself and his seed after him from generation to generation; If he will do my will let him not take his family unto the eastern lands, even unto Kirtland: nevertherless I the Lord will build up Kirtland, but I the Lord have a scourge prepared for the inhabitants thereof; and with my servant Alman Babbit [Almon Babbitt] there are many things with which I am not well pleased; behold he aspireth to establish his \sout{council} counsel instead of the \sout{council} counsel which I have ordained, even the presidency of my Church, and he setteth up a golden calf, for the worship of my people. Let no man go from this place who \sout{hath} has come here, assaying to keep my commandments; if they live here let them live unto me, and if they die here let them die unto me, for they shall rest from all their labors here, and shall continue their works. Therefore let my servant William put his trust in me, and cease to fear concerning his family because of the sickness of the land, If ye love me Keep my commandments and the sickness of the land shall redound to your glory.

				Let my Servant William go and proclaim mine everlasting gospel with a loud voice and with great joy as he shall be moved upon by my spirit unto the inhabitants of Warsaw and also unto the inhabitants of Carthage, and also unto the inhabitants of Madison [Fort Madison, Iowa Territory], and also unto the inhabitants of Burlington and await patiently and diligently for further instructions at my general conference saith the Lord. <If he will do my will> let him from henceforth hearken to the \sout{council} counsel of my servant Joseph and with his interest support the cause of the poor and publish the new translation of my holy word unto the inhabitants of the earth, and if he will do this I will bless him with a multiplicity of blessings, that he shall not be forsaken, nor his seed be found begging bread.

				And, again, verily I say unto you that my servant William, be appointed, ordained, and anointed as a \sout{councillor} counsellor unto my servant Joseph in the room of my servant Hyrum, that my servant Hyrum may take the office of priesthood and patriarch, which was appointed unto him by his father by blessing and also by right, that from henceforth he shall hold the keys of the patriarchal blessings upon the heads of all my people, that who ever he blesses shall be blessed, and who ever he curseth shall be cursed, that whatsoever he shall bind on the earth shall be bound in heaven, and, that whatsoever he shall loose on earth shall be loosed in heaven, and from this time forth I appoint unto him, that he may be a prophet and a seer and a revelator unto my church as well as my servant Joseph, that he may act in concert also, with my servant Joseph, and that he shall receive \sout{council} counsel from my servant Joseph, who shall shew unto him the keys whereby he may ask and receive, and be crowned with the same blessings. I crown upon his head, the bishoprick and blessing and glory, and honor and priesthood and gifts of the priesthood, that once \sout{was} were put upon him that was my servant Oliver Cowdery; that my servant Hyrum may bear record of the things which I shall shew unto him, that his name may be had in honorable remembrance from generation to generation for ever and ever. Let my servant William <Law> also receive the Keys by which he may ask and receive blessings, let him be humble before me and be without guile and he shall receive of my spirit, even the comforter, which shall manifest unto him the truth of all things, and shall give him in the very hour, what he shall say, and these signs shall follow him, he shall heal the sick, he shall cast out Devils, and shall be delivered from those who would administer unto him deadly poison, and shall be led in paths where the poisonous serpents cannot lay hold upon his heel; and he shall mount up in the imagination of his thoughts as upon eagles wings, and what if I will that he should raise the dead, let him not withold his voice, therefore let my servant William cry aloud and spare not, with joy and rejoicing, and with hozannas to him that sitteth upon the throne for ever and ever saith the Lord your God.

				Behold I say unto you I have a mission in store for my servant William and my servant Hyrum and for them alone, and let my servant Joseph tarry at home, for he is needed, the remainder I will shew <unto> you hereafter, even so, Amen.

				And, again, verily I say unto you, if my servant Sidney [Rigdon] will serve me and be a \sout{councillor} counsellor unto my servant Joseph, let him arise and come up and stand in the office of his calling, and humble himself before me and if he will offer unto me an acceptable offering and acknowledgements and remain with my people, behold I the Lord your God will heal him that he shall be healed and shall lift up his voice again on the mountains, and be a spokesman before my face; let him come and locate his family in the neighbourhood in which my servant Joseph resides; and in all his journeyings let him lift up his voice as with the sound of a trump and warn the inhabitants of the earth to flee the wrath to come; let him assist my servant Joseph, and also let my servant William Law assist my servant Joseph in making a solemn proclamation unto the kings of the earth even as I have before said unto you.

				If my servant Sidney will do my will, let him not move his family unto the eastern lands, but let him change their habitation even as I have said: Behold, it is not my will that he should seek to find safety and refuge out of the city which I have appointed unto you, even the city of Nauvoo. Verily I say unto you, even now, that if he will hearken unto my voice it shall be well with him, even so, Amen.

				And again verily I say unto you, let my servant Amos Davis pay stock into the hands of those whom I have appointed to build a house for boarding, even the Nauvoo House, this let him do if he will have an interest and let him hearken unto the \sout{council} counsel of my servant Joseph, and labor with his own hands that he may obtain the confidence of men, and when he shall prove himself faithful in all things that shall be entrusted unto his care, yea even a few things, he shall be made ruler over many; let him therefore abase himself, that he may be exalted, even so, Amen.

				And again, verily I say unto you if my servant Robert D. Foster will obey my voice, let him build a house for my servant Joseph according to the contract which <he> has made with him, as the door shall be open to him from time to time, and let him repent of all his folly and clothe himself with charity, and cease to do evil and lay aside all his hard speeches, and pay stock also into the hands of the quorum of the Nauvoo House for himself and for his generation after him, from generation to generation and hearken unto the \sout{council} counsel of my servant Joseph and Hyrum and William Law, and unto the authorities which I have called to lay the foundation of Zion, and it shall be well with him, forever and ever, even so Amen.

				And again, verily I say unto you let \sout{my} no man pay stock to the quorum of the Nauvoo House, unless he shall be a believer in the Book of Mormon and the revelations I have given unto you saith the Lord your God for that which is more or less than this cometh of evil, and shall be attended with cursings and not blessing saith the Lord your God, even so Amen

				And again, verily I say unto you let the quorum of the Nauvoo House, have a just recompence of wages, for all their labors which they do in building the Nauvoo House, and let their wages be as shall be agreed among themselves as pertaining to the price thereof, and let every man who pays stock bear his proportion of their wages, if it must needs be, for their support saith the Lord, otherwise their labors shall be accounted unto them for stock in that house even so, Amen.

				Verily I say unto you, I now give unto you the offices belonging to my priesthood that ye may hold the keys thereof even the Priesthood which is after the order of Melchisadeck, which is after the order of my only begotten son. First, I give unto you Hyrum Smith to be a patriarch unto you to hold the sealing blessings of my church, even the holy spirit of promise, whereby ye are sealed up unto the day of redemption, that ye may not fall, notwithstanding the hour of temptation that may come upon you.

				I give unto you my servant Joseph to be a presiding Elder over all my Church, to be a Translater, a Revelator, a Seer, and Prophet; I give unto him for \sout{Councillors} Counsellors my servant Sidney Rigdon and my servant William Law that these may constitute a quorum and first presidency to receive the oracles for the whole church.

				I give unto you my servant Brigham Young to be a president over the twelve travelling council which twelve hold the keys to open up the authority of my kingdom upon the four corners of the earth and after that to send my word to every creature, they are Heber C Kimbal[l], Parley P. Pratt, Orson Pratt, Orson Hyde, William Smith, John Taylor, John E. Page Wilford Woodruff, Willard Richards, George A. Smith; David Patten I have taken unto myself, behold his priesthood no man taketh from him, but verily I say unto you, another may be appointed unto the same calling.

				And, again, I say unto you, I give unto you I give unto you a a High Council for the corner stone of Zion viz Samuel Bent, H. G. Sherwood, George W Harris, Charles C. Rich, Thomas Grover, Newel Knight, David Dort, [Lewis] Dunbar Wilson, Seymour Brunson I have taken unto myself, no man taketh his priesthood, but another may be appointed unto the same priesthood in his stead (and verily I say unto <you> let my servant Aaron Johnson be ordained unto this calling in his stead) David Ful[l]mer, Alpheus Cutler, William Huntington.

				And again I give unto you Don C[arlos] Smith to be a president over a quorum of high priests, which ordinance is instituted for the prupose of qualifying those who shall be appointed as standing presidents or servants \sout{of} over different stakes scattered abroad and they may travel also if they choose, but rather be ordained for standing presidents, this is the office of their calling sath the Lord your God, I give unto him Amasa Lyman and Noah Packard for councillors that they may preside over the quorum of high priests of my church saith the Lord.

				And again I say unto you I give unto you I give unto you John A Hicks, Samuel Williams, and \sout{over the} <Jesse Baker> which priesthood is to preside over the quorum of Elders, which quorum is instituted for standing ministers, nevertherless they may travel, yet they are ordained to be standing ministers to my church saith the Lord.

				And, again, I give unto you Joseph Young Josiah Butterfield, Daniel Miles, Henry Herriman [Harriman] Zera Pulsipher, Levi Hancock, James Foster to preside over the quorum of seventies, which quorum is instituted for travelling Elders to bear record of my name in all the world, wherever the travelling high council, my apostles shall send them to prepare a\sout{n} way before my face, the difference between this quorum and the quorum of Elders is, that one is to travel continually and the other is to preside over the churches from time to time, the one has the responsibility of presiding from time to time, but the other has no responsibiltiy of presiding \sout{from time to time} saith the Lord your God

				And, again, I say unto you, I give unto you Vinson Knight, Samuel H. Smith and Shadrick [Shadrach] Roundy if he will receive it, to preside over the Bishoprick, a knowledge of said bishoprick is given unto you in the Book of Doctrine and Covenants.

				And, again, I \sout{say} give unto you, Samuel Rolfe and his councillors for priests, and the president of the Teachers and his councillors and also the president of the Deacons and his councillors, and also the president of the stake and his councillors the above offices I have given unto you and the keys thereof for helps and for governments for the work of the ministry, and the perfecting of my saints and a commandment I give unto you that you should fill all these offices and approve of those names which I have mentioned, or else disapprove of them at my general conference; and that ye should prepare rooms for all these offices in my house when you build it unto my name saith the Lord your God. even so. Amen.
				
		% -----
		\subsubsection{Letter to Oliver Granger, 26 January 1841}
		% -----
			\label{EMS-1841-JS-26-JAN-1841}
			% \label{EMS-1840-JS-DAY-3LETTERMONTH-YEAR}
			
			\begin{description}
				\item[Print Sources]
			\end{description}

			\begin{tabular}{p{0.5in}p{1in}p{0.5in}p{1in}}
				JSP	 	& ---		& JSR 		& --- \\
				DC	 	& ---		& HC 		& --- \\
				HDDC 	& ---		& TCHC 		& --- \\
			\end{tabular}
			% HDDC = woodford's DC diss; JSR = Marquardt's JS Revelations; TCHC = Vogel HC
			
			\begin{description}
				\item[Online Access] \href{https://www.josephsmithpapers.org/paper-summary/letter-to-oliver-granger-26-january-1841/1}{Here}
				% \item[Analysis] \hyperref[XX]{Here}
			\end{description}
			
			\begin{description}
				\item[Background]
			\end{description}
			
				% It might also be useful to give a two sentence context of the period: e.g., "This speech was given in Nauvoo to a small congregation one month prior to the destruction of the Nauvoo Expositor. These and other tensions may have been on the prophet's mind when he gave the speech."
			
			\begin{description}
				\item[Original Source]
			\end{description}
			
				% brief description of original source material, with reference to JSP or somewhere else for more info
				% Background material: it might be helpful to have a note that says whether a given document was presented as a speech, was written in a journal, etc.
				% Also a comment here about how reliable the source might be, or what shape it is in, if there are large omissions or parts that are hard to understand, and so on.
			
			\begin{description}
				\item[Text]
			\end{description}
			
				\begin{tightright}
						\hypertarget{EMS-1841-JS-26-JAN-1841-a}%
					City
						\marginnote{a}%
					of Nauvoo Jan\supers{y.} 26, 1841
				\end{tightright}

				Dear Brother [Oliver] Granger,

				I wrote you a few days ago in answer to a letter recently received from you, expressing my surprise that I had not received the letters you mentioned; Since that time, I have to inform you that I have received two letters giving me an account of the proceeding of the Church in Kirtland and some of your business transactions. I am extreemely sorry that the person by whom you sent those letters did not hand them over to me sooner, however, I assure you, I was glad to \sout{pe$\diamond\diamond\diamond\diamond$} <read> them and felt very much satisfied with the perusal. I was likewise very much pleased with the spirit which was manifest by the saints in Kirtland and for their desire to promote the interests of the Kingdom.

					\hypertarget{EMS-1841-JS-26-JAN-1841-b}%
				If
					\marginnote{b}%
				those letters had been received in their propper time, \sout{the} probably might have acted differently at the last conference, but not having the information we desired, we acted to the best of our understanding, which I hope will prove advantageous to all parties. I hope you will not let any of the proceeding of last Conference disturb your mind, for we understood you were intending to come here last fall, and consequently thought it advisable to appoint some one to preside in Kirtland. I should be glad if you would co-opperate with Elder Babbit [Almon Babbitt] and both lay your Shoulders to the work, and if you do so I think you will be a blessing to the Church, and prosperity will Smile upon you.

					\hypertarget{EMS-1841-JS-26-JAN-1841-c}%
				I
					\marginnote{c}%
				should be much pleased to receive a letter from you as oft as you can make it convenient, and give me all the intelligence in your power. I am yet in the dark respecting the New York debts I should be much pleased to hear that they were settled.

				Since writing the above I received yours of the 9\supersu{th} inst, and I assure you I was very much gratified to hear of your success in redeeming the ``Lords House'' \&c I hope Dear Brother that success will attend all your efforts for the prosperity of the cause, so dear to the saints and that you will be abundantly rewarded by that God who\sout{se} has called us to be co-workers with him and his holy spirit in these last days

					\hypertarget{EMS-1841-JS-26-JAN-1841-d}%
				Be
					\marginnote{d}%
				assured of my continued regard for your welfare and for the prosperity of the Saints in Kirtland, and I pray that they may prosper in every good word and work and after the afflictions and tribulations of mortality be crowned with everlasting joy in the celestial kingdom of our God.

				We continue to prosper in this place and expect to have a large increase of inhabitants the next Season, we are making preperation for some large and extensive buildings which we intend to \sout{build} Erect soon.

				With sentiments of respect, I am very respectfully yours Affectionately

				\begin{tightright}
					\uline{Joseph Smith}
				\end{tightright}

					\hypertarget{EMS-1841-JS-26-JAN-1841-e}%
				P.S
					\marginnote{e}%
				With respect to giving advise to the Saints moving to Kirtland I would say, that I feel desireous that the Eastern brethren should come to this place, but at the same time, those who had rather move to Kirtland than to this place are at liberty to do so.

				I am pleased you have secured the keys of the Lords House and should advise you and you are hereby requested to hold them until I \sout{home} come, \sout{whe} I cannot say when I shall pay you a visit but I think not before the New York debts be settled. So if you are desireous to see me, you will have to exert yourselves to get the debts settled.

				Please to write me all the p[ar]ticulars of your transactions and let me know when we shall have the pleasu[re] of Seeing you here

				\begin{tightright}
					J. S.
				\end{tightright}

					\hypertarget{EMS-1841-JS-26-JAN-1841-f}%
				D\supers{r}
					\marginnote{f}%
				\phantom{ }[Isaac] Galland accompanied by some one of the brethren (perhaps Hyrum [Smith]) intend to leave this place in a few days for the East, on business for the church, they will call at Kirtland and will be glad to see you as you may possibly give them such information as \sout{may} will be necessary for them in their business transactions.

				\begin{tightright}
					J. S.
				\end{tightright}

				<NAUVOO ILL

				Ja. 29

				\begin{tightright}
					<25>
				\end{tightright}

				M\supersu{r}\supers{.} Oliver Granger

				Kirtland

				Lake Co

				Ohio
				
		% -----
		\subsubsection{Account of Meeting and Discourse, circa 2 February 1841}
		% -----
			\label{EMS-1841-JS-CIR-2-FEB-1841}
			% \label{EMS-1840-JS-DAY-3LETTERMONTH-YEAR}
			
			% --
			\paragraph{As Reported by William P. McIntire}
			% --
				\label{EMS-1841-JS-CIR-2-FEB-1841-WM}
			
				\begin{description}
					\item[Print Sources]
				\end{description}

				\begin{tabular}{p{0.5in}p{1in}p{0.5in}p{1in}}
					JSP	 	& ---		& JSR 		& --- \\
					DC	 	& ---		& HC 		& --- \\
					HDDC 	& ---		& TCHC 		& --- \\
				\end{tabular}
				% HDDC = woodford's DC diss; JSR = Marquardt's JS Revelations; TCHC = Vogel HC
				
				\begin{description}
					\item[Online Access] \href{https://www.josephsmithpapers.org/paper-summary/account-of-meeting-and-discourse-circa-2-february-1841/1}{Here}
					% \item[Analysis] \hyperref[XX]{Here}
				\end{description}
				
				\begin{description}
					\item[Background]
				\end{description}
				
					% It might also be useful to give a two sentence context of the period: e.g., "This speech was given in Nauvoo to a small congregation one month prior to the destruction of the Nauvoo Expositor. These and other tensions may have been on the prophet's mind when he gave the speech."
				
				\begin{description}
					\item[Original Source]
				\end{description}
				
					% brief description of original source material, with reference to JSP or somewhere else for more info
					% Background material: it might be helpful to have a note that says whether a given document was presented as a speech, was written in a journal, etc.
					% Also a comment here about how reliable the source might be, or what shape it is in, if there are large omissions or parts that are hard to understand, and so on.
				
				\begin{description}
					\item[Text]
				\end{description}
				
					Next Meeting--- Joseph said the Lord said that we should build our house to his Name that we might be Baptized for the Dead--- But if we Did it Not we should be Rejected \& our Dead with us \& this Church should Not be excepted [accepted]
				
					Bro [William] Law Lectured on the 11\supersu{th} of Rom of the powers that be ordained of God the Law of Order in the heavenly Bodys as also in all animal Creation
					
		% -----
		\subsubsection{Account of Meeting and Discourse, circa 9 February 1841}
		% -----
			\label{EMS-1841-JS-CIR-9-FEB-1841}
			% \label{EMS-1840-JS-DAY-3LETTERMONTH-YEAR}
			
			% --
			\paragraph{As Reported by William P. McIntire}
			% --
				\label{EMS-1841-JS-CIR-9-FEB-1841-WM}
			
				\begin{description}
					\item[Print Sources]
				\end{description}

				\begin{tabular}{p{0.5in}p{1in}p{0.5in}p{1in}}
					JSP	 	& ---		& JSR 		& --- \\
					DC	 	& ---		& HC 		& --- \\
					HDDC 	& ---		& TCHC 		& --- \\
				\end{tabular}
				% HDDC = woodford's DC diss; JSR = Marquardt's JS Revelations; TCHC = Vogel HC
				
				\begin{description}
					\item[Online Access] \href{https://www.josephsmithpapers.org/paper-summary/account-of-meeting-and-discourse-circa-9-february-1841/1}{Here}
					% \item[Analysis] \hyperref[XX]{Here}
				\end{description}
				
				\begin{description}
					\item[Background]
				\end{description}
				
					% It might also be useful to give a two sentence context of the period: e.g., "This speech was given in Nauvoo to a small congregation one month prior to the destruction of the Nauvoo Expositor. These and other tensions may have been on the prophet's mind when he gave the speech."
				
				\begin{description}
					\item[Original Source]
				\end{description}
				
					% brief description of original source material, with reference to JSP or somewhere else for more info
					% Background material: it might be helpful to have a note that says whether a given document was presented as a speech, was written in a journal, etc.
					% Also a comment here about how reliable the source might be, or what shape it is in, if there are large omissions or parts that are hard to understand, and so on.
				
				\begin{description}
					\item[Text]
				\end{description}
					
						\hypertarget{EMS-1841-JS-CIR-9-FEB-1841-WM-a}%
					Next
						\marginnote{a}%
					Meeting

					I Hicks on prejedice \& its Results first he Refered to the Maner of Mothers prejudiceing their Children in Infancy By telling faulseties \& Cheating them---

					Next Lorenzo the Benefits of such a Meeting or Schools--- he treated on the Bumps also---

						\hypertarget{EMS-1841-JS-CIR-9-FEB-1841-WM-b}%
					Next
						\marginnote{b}%
					father Cole on the fullness of the Gospel--- Born Benette [John C. Bennett] on Cretecizms said the Body aught to be perfect Easy in a Natureal way free from being Cramped By Clinching the hands, \& Gestures suitable aught to acompany, so that it May be More pleasing \& of Wait--- the speaker aught to give time between sentences tell his hearers Comprehends--- \& if they do not, Repeat it again--- the Speaker aught to have that Much decernment, to know by the Countenances of his hearers when the do Comprehend his word or sentance--- he says the eye aught to follow all apropreate Gestures By the hand or hands, highth width or Depth---

						\hypertarget{EMS-1841-JS-CIR-9-FEB-1841-WM-c}%
					Joseph
						\marginnote{c}%
					said in answer to <M\supersu{r}.> [Hosea] stout that adam Did Not Comit sin in [e]ating the fruit, for God had Dec[r]eed that he should Eat \& fall--- But in complyance with the Decree he should Die--- only he should Die was the saying of the Lord therefore the Lord apointed us to fall \& also Redeemed us--- for where sin abounded Grace did Much More abound--- for Paul says Rom.--- 5--- 10 for if--- when were Enemys we were Reconciled to God by <the Death of> his Son, much more, being Reconciled, we shall be saved by his Life---
					
		% -----
		\subsubsection{Account of Meeting and Discourse, circa 16 February 1841}
		% -----
			\label{EMS-1841-JS-CIR-16-FEB-1841}
			% \label{EMS-1840-JS-DAY-3LETTERMONTH-YEAR}
			
			% --
			\paragraph{As Reported by William P. McIntire}
			% --
				\label{EMS-1841-JS-CIR-16-FEB-1841-WM}
			
				\begin{description}
					\item[Print Sources]
				\end{description}

				\begin{tabular}{p{0.5in}p{1in}p{0.5in}p{1in}}
					JSP	 	& ---		& JSR 		& --- \\
					DC	 	& ---		& HC 		& --- \\
					HDDC 	& ---		& TCHC 		& --- \\
				\end{tabular}
				% HDDC = woodford's DC diss; JSR = Marquardt's JS Revelations; TCHC = Vogel HC
				
				\begin{description}
					\item[Online Access] \href{https://www.josephsmithpapers.org/paper-summary/account-of-meeting-and-discourse-circa-16february-1841/1}{Here}
					% \item[Analysis] \hyperref[XX]{Here}
				\end{description}
				
				\begin{description}
					\item[Background]
				\end{description}
				
					% It might also be useful to give a two sentence context of the period: e.g., "This speech was given in Nauvoo to a small congregation one month prior to the destruction of the Nauvoo Expositor. These and other tensions may have been on the prophet's mind when he gave the speech."
				
				\begin{description}
					\item[Original Source]
				\end{description}
				
					% brief description of original source material, with reference to JSP or somewhere else for more info
					% Background material: it might be helpful to have a note that says whether a given document was presented as a speech, was written in a journal, etc.
					% Also a comment here about how reliable the source might be, or what shape it is in, if there are large omissions or parts that are hard to understand, and so on.
				
				\begin{description}
					\item[Text]
				\end{description}
				
						\hypertarget{EMS-1841-JS-CIR-16-FEB-1841-WM-a}%
					Next
						\marginnote{a}%
					Meetting Badlam [Alexander Badlam Sr.] on Mechanical Powers \& utility--- S[amuel] H. Smith on Virtue--- T[heodore] Turley on the Results of Eivel---
					
						\hypertarget{EMS-1841-JS-CIR-16-FEB-1841-WM-b}%
					Joseph
						\marginnote{b}%
					said Concerning the God-head it was Not as many imagined--- three Heads \& but one body; he said the three were separate bodys, God the first \& Jesus the Mediator the 2\supersu{d} \& the Holy Ghost \& these three agree in one \& this is the man[n]er we should aproach God in order to get his blessings \& he also said Every man \sout{ware} is stimulated by a certain motive, \sout{to act} motive preceeds action \& if we want to know ourselves this is the key to Examine the motive what it is, \& the fact will be manifest
					
		% -----
		\subsubsection{Account of Meeting and Discourse, circa 23 February 1841}
		% -----
			\label{EMS-1841-JS-CIR-23-FEB-1841}
			% \label{EMS-1840-JS-DAY-3LETTERMONTH-YEAR}
			
			% --
			\paragraph{As Reported by William P. McIntire}
			% --
				\label{EMS-1841-JS-CIR-23-FEB-1841-WM}
			
				\begin{description}
					\item[Print Sources]
				\end{description}

				\begin{tabular}{p{0.5in}p{1in}p{0.5in}p{1in}}
					JSP	 	& ---		& JSR 		& --- \\
					DC	 	& ---		& HC 		& --- \\
					HDDC 	& ---		& TCHC 		& --- \\
				\end{tabular}
				% HDDC = woodford's DC diss; JSR = Marquardt's JS Revelations; TCHC = Vogel HC
				
				\begin{description}
					\item[Online Access] \href{https://www.josephsmithpapers.org/paper-summary/account-of-meeting-and-discourse-circa-23february-1841/1}{Here}
					% \item[Analysis] \hyperref[XX]{Here}
				\end{description}
				
				\begin{description}
					\item[Background]
				\end{description}
				
					% It might also be useful to give a two sentence context of the period: e.g., "This speech was given in Nauvoo to a small congregation one month prior to the destruction of the Nauvoo Expositor. These and other tensions may have been on the prophet's mind when he gave the speech."
				
				\begin{description}
					\item[Original Source]
				\end{description}
				
					% brief description of original source material, with reference to JSP or somewhere else for more info
					% Background material: it might be helpful to have a note that says whether a given document was presented as a speech, was written in a journal, etc.
					% Also a comment here about how reliable the source might be, or what shape it is in, if there are large omissions or parts that are hard to understand, and so on.
				
				\begin{description}
					\item[Text]
				\end{description}
				
					Next Metting first Shearwood [Henry G. Sherwood] on the Authantisitty of the book of Mormon 2 V Night [Vinson Knight] on Education--- \supers{d} Davis on Charity--- 4 [George] Miller on Man[n]ers---
					
					Joseph said he Never wanted to hear a man snore Louder than he coul[d] shout in battle--- he Did Not want a man say O Joseph how I Love you \&cc \& when the time of Danger come forsake him
					
		% -----
		\subsubsection{Letter to Vilate Murray Kimball, 2 March 1841}
		% -----
			\label{EMS-1841-JS-2-MAR-1841}
			% \label{EMS-1840-JS-DAY-3LETTERMONTH-YEAR}
			
			\begin{description}
				\item[Print Sources]
			\end{description}

			\begin{tabular}{p{0.5in}p{1in}p{0.5in}p{1in}}
				JSP	 	& ---		& JSR 		& --- \\
				DC	 	& ---		& HC 		& --- \\
				HDDC 	& ---		& TCHC 		& --- \\
			\end{tabular}
			% HDDC = woodford's DC diss; JSR = Marquardt's JS Revelations; TCHC = Vogel HC
			
			\begin{description}
				\item[Online Access] \href{https://www.josephsmithpapers.org/paper-summary/letter-to-vilate-murray-kimball-2march-1841/1}{Here}
				% \item[Analysis] \hyperref[XX]{Here}
			\end{description}
			
			\begin{description}
				\item[Background]
			\end{description}
			
				% It might also be useful to give a two sentence context of the period: e.g., "This speech was given in Nauvoo to a small congregation one month prior to the destruction of the Nauvoo Expositor. These and other tensions may have been on the prophet's mind when he gave the speech."
			
			\begin{description}
				\item[Original Source]
			\end{description}
			
				% brief description of original source material, with reference to JSP or somewhere else for more info
				% Background material: it might be helpful to have a note that says whether a given document was presented as a speech, was written in a journal, etc.
				% Also a comment here about how reliable the source might be, or what shape it is in, if there are large omissions or parts that are hard to understand, and so on.
			
			\begin{description}
				\item[Text]
			\end{description}
			
				\begin{tightright}
						\hypertarget{EMS-1841-JS-2-MAR-1841-a}%
					City
						\marginnote{a}%
					of Nauvoo March 2, 1841
				\end{tightright}

				Dear Sister [Vilate Murray] Kimball.

				I can in some measure enter into your feelings respecting the occurrence which has lately taken place in the church, which is indeed painful to every lover of Truth and Holiness, and probably to none more so than myself. I am indeed sorry that any thing should have transpired which should have caused such a stir in the Church, and bro't [brought] disgrace upon persons who are otherwise respectable. The course I have taken in the matter was such as I felt warranted to take from the testimony which was adduced. Whether they were guilty of crime or not I do not say, but this I must say that their imprudence was carried to an unwarrantable extent.

					\hypertarget{EMS-1841-JS-2-MAR-1841-b}%
				I
					\marginnote{b}%
				do not desire that you should turn the young woman out of doors, far be it from me to advise any such course I think it would be well for her to remain with you at least until Bro [Heber C.] Kimball comes home, because I think that your advise \sout{to her}, may be a blessi[n]g to \uline{her}, and your council and advise such as will tend to her future welfare and happiness. I have no doubt but you will act in wisdom in this matter \& remain \sout{very} Yours in the Gospel

				\begin{tightright}
					Joseph Smith
				\end{tightright}

				[\textit{page [2] blank}]
				
		% -----
		\subsubsection{Account of Meeting and Discourse, circa 2 March 1841}
		% -----
			\label{EMS-1841-JS-CIR-2-MAR-1841}
			% \label{EMS-1840-JS-DAY-3LETTERMONTH-YEAR}
			
			% --
			\paragraph{As Reported by William P. McIntire}
			% --
				\label{EMS-1841-JS-CIR-2-MAR-1841-WM}
			
				\begin{description}
					\item[Print Sources]
				\end{description}

				\begin{tabular}{p{0.5in}p{1in}p{0.5in}p{1in}}
					JSP	 	& ---		& JSR 		& --- \\
					DC	 	& ---		& HC 		& --- \\
					HDDC 	& ---		& TCHC 		& --- \\
				\end{tabular}
				% HDDC = woodford's DC diss; JSR = Marquardt's JS Revelations; TCHC = Vogel HC
				
				\begin{description}
					\item[Online Access] \href{https://www.josephsmithpapers.org/paper-summary/account-of-meeting-and-discourse-circa-2march-1841/1}{Here}
					% \item[Analysis] \hyperref[XX]{Here}
				\end{description}
				
				\begin{description}
					\item[Background]
				\end{description}
				
					% It might also be useful to give a two sentence context of the period: e.g., "This speech was given in Nauvoo to a small congregation one month prior to the destruction of the Nauvoo Expositor. These and other tensions may have been on the prophet's mind when he gave the speech."
				
				\begin{description}
					\item[Original Source]
				\end{description}
				
					% brief description of original source material, with reference to JSP or somewhere else for more info
					% Background material: it might be helpful to have a note that says whether a given document was presented as a speech, was written in a journal, etc.
					% Also a comment here about how reliable the source might be, or what shape it is in, if there are large omissions or parts that are hard to understand, and so on.
				
				\begin{description}
					\item[Text]
				\end{description}
				
					Next Meeting W Soby on true friendship--- he said true Friendship was such as a man Laying Down his Life for his Friend. Joseph said there was a diferance Between the veng[e]ance that Belongeth to the Lord, \& a man Defending himself or friend--- of the Mamon of unrighteousness Joseph said the majority of man king [mankind] was $\diamond\diamond$ist \& true a must Lavish his Goods on all around him \& out of them he will find freinds in the hour of Distress for those that have been made Ritch By the Benevalance of their Rich friend of them there will be some that will do you Good in affliction---
					
		% -----
		\subsubsection{Account of Meeting and Discourses, circa 9 March 1841}
		% -----
			\label{EMS-1841-JS-CIR-9-MAR-1841}
			% \label{EMS-1840-JS-DAY-3LETTERMONTH-YEAR}
			
			% --
			\paragraph{As Reported by William P. McIntire}
			% --
				\label{EMS-1841-JS-CIR-9-MAR-1841-WM}
			
				\begin{description}
					\item[Print Sources]
				\end{description}

				\begin{tabular}{p{0.5in}p{1in}p{0.5in}p{1in}}
					JSP	 	& ---		& JSR 		& --- \\
					DC	 	& ---		& HC 		& --- \\
					HDDC 	& ---		& TCHC 		& --- \\
				\end{tabular}
				% HDDC = woodford's DC diss; JSR = Marquardt's JS Revelations; TCHC = Vogel HC
				
				\begin{description}
					\item[Online Access] \href{https://www.josephsmithpapers.org/paper-summary/account-of-meeting-and-discourses-circa-9march-1841/1}{Here}
					% \item[Analysis] \hyperref[XX]{Here}
				\end{description}
				
				\begin{description}
					\item[Background]
				\end{description}
				
					% It might also be useful to give a two sentence context of the period: e.g., "This speech was given in Nauvoo to a small congregation one month prior to the destruction of the Nauvoo Expositor. These and other tensions may have been on the prophet's mind when he gave the speech."
				
				\begin{description}
					\item[Original Source]
				\end{description}
				
					% brief description of original source material, with reference to JSP or somewhere else for more info
					% Background material: it might be helpful to have a note that says whether a given document was presented as a speech, was written in a journal, etc.
					% Also a comment here about how reliable the source might be, or what shape it is in, if there are large omissions or parts that are hard to understand, and so on.
				
				\begin{description}
					\item[Text]
				\end{description}
				
						\hypertarget{EMS-1841-JS-CIR-9-MAR-1841-WM-a}%
					Subject
						\marginnote{a}%
					1\supersu{st}\supers{.} on the Gospel By father Cole he said that \sout{th} Some thought that He difered from president Smith Concerning the time of the Giving of the Holy Ghost--- as teach that all men Receive the Holy Ghost before Baptizem--- Joseph said men do not take Notice of things as they Read them--- or they might know things as they Read them--- he quotes Petter 2\supersu{d}\supers{.} Repent \& be Baptized \&c--- \& ye shall Receive the Gift of the Holly Ghost--- Now said he (taking up his Cap \& presenting to Pre\supersu{d} [William] Law) in Giveing you this Gift is not Giving myself--- 
						\hypertarget{EMS-1841-JS-CIR-9-MAR-1841-WM-b}%
					However
						\marginnote{b}%
					there is a peistHood [priesthood] \sout{in} With <the> Holy Ghost \& Key--- the Holy Ghost over shadows you \& witness unto you of the authority \& the Gifts of the Holy Ghost--- he said was the provence of the Father to preside as the Chief or president--- Jesus as the Mediator \& Holy Ghost as the testator or witness--- the Son Had a Tabernicle \& so had the father But the Holly Ghost is a personage of spirit without tabernicle the Great God has a Name By wich he will be Called Which is Ahman--- also in asking have Referance to a personage Like Adam for God made Adam Just in his own Image Now this a key for you to know how to ask \& obtain

					Next was Robinson on the Gospel he spoke of adam pertaking of the forbiden fruit \& so-on---

						\hypertarget{EMS-1841-JS-CIR-9-MAR-1841-WM-c}%
					Joseph
						\marginnote{c}%
					said a person ought always take his subject along--- for instance Jesus in all his Doctrines \& pariables set it forth Like a tree or mustard stock Comencing at the foundation \& Root thence up to the first Branch \& out along it then back to the body of the tree or subject \& thence to the Next Branch \& so-on Now as to Adam the Lord said in the day thou shalt Eat there of thou shalt shurely Die Now the Day the Lord has Refferance too is spoken of By Petter a thousand of our years is with the Lord as one Day \&c at the time the Lord said this to adam there was No mode of Counting time By man, as man Now Counts time

					Next was Bron Benet [John C. Bennett] on Phisickal Education for the Health \& preservation of the Lungs---
					
		% -----
		\subsubsection{Revelation, circa Early March 1841 [D\&C 125]}
		% -----
			\label{EMS-1841-JS-MAR-1841}
			% \label{EMS-1840-JS-DAY-3LETTERMONTH-YEAR}
			
			\begin{description}
				\item[Print Sources]
			\end{description}

			\begin{tabular}{p{0.5in}p{1in}p{0.5in}p{1in}}
				JSP	 	& ---		& JSR 		& --- \\
				DC	 	& ---		& HC 		& --- \\
				HDDC 	& ---		& TCHC 		& --- \\
			\end{tabular}
			% HDDC = woodford's DC diss; JSR = Marquardt's JS Revelations; TCHC = Vogel HC
			
			\begin{description}
				\item[Online Access] \href{https://www.josephsmithpapers.org/paper-summary/revelation-circa-early-march-1841-dc-125/1}{Here}
				% \item[Analysis] \hyperref[XX]{Here}
			\end{description}
			
			\begin{description}
				\item[Background]
			\end{description}
			
				% It might also be useful to give a two sentence context of the period: e.g., "This speech was given in Nauvoo to a small congregation one month prior to the destruction of the Nauvoo Expositor. These and other tensions may have been on the prophet's mind when he gave the speech."
			
			\begin{description}
				\item[Original Source]
			\end{description}
			
				% brief description of original source material, with reference to JSP or somewhere else for more info
				% Background material: it might be helpful to have a note that says whether a given document was presented as a speech, was written in a journal, etc.
				% Also a comment here about how reliable the source might be, or what shape it is in, if there are large omissions or parts that are hard to understand, and so on.
			
			\begin{description}
				\item[Text]
			\end{description}
			
				A Revelation given in the City of Nauvoo in answer to the following interrogatory.
				
				What is the will of the Lord concerning the Saints in the Territory of Iowa?
				
				Verily thus saith the Lord, I say unto you if those who call themselves by <my> name, and are assaying to be my saints, if they will do my will and keep my commandments concerning them; let them gather themselves together unto the places which I shall appoint unto them by my servant Joseph, and build up cities unto my name, that they may be prepared for that which is in store for a time to come. Let them build up a city unto my name upon the land opposite to the city of Nauvoo and let the name of Zarahemla be named upon it. And let all those who come from the east and the west, and the north and the south that have desires to dwell therein, take up their inheritances in the same, as well as in the City of Nashville or in the City of Nauvoo, and in all the stakes which I have appointed saith the Lord.
				
		% -----
		\subsubsection{Account of Meeting and Discourse, circa 16 March 1841}
		% -----
			\label{EMS-1841-JS-CIR-16-MAR-1841}
			% \label{EMS-1840-JS-DAY-3LETTERMONTH-YEAR}
			
			% --
			\paragraph{As Reported by William P. McIntire}
			% --
				\label{EMS-1841-JS-CIR-16-MAR-1841-WM}
			
				\begin{description}
					\item[Print Sources]
				\end{description}

				\begin{tabular}{p{0.5in}p{1in}p{0.5in}p{1in}}
					JSP	 	& ---		& JSR 		& --- \\
					DC	 	& ---		& HC 		& --- \\
					HDDC 	& ---		& TCHC 		& --- \\
				\end{tabular}
				% HDDC = woodford's DC diss; JSR = Marquardt's JS Revelations; TCHC = Vogel HC
				
				\begin{description}
					\item[Online Access] \href{https://www.josephsmithpapers.org/paper-summary/account-of-meeting-and-discourse-circa-16march-1841/1}{Here}
					% \item[Analysis] \hyperref[XX]{Here}
				\end{description}
				
				\begin{description}
					\item[Background]
				\end{description}
				
					% It might also be useful to give a two sentence context of the period: e.g., "This speech was given in Nauvoo to a small congregation one month prior to the destruction of the Nauvoo Expositor. These and other tensions may have been on the prophet's mind when he gave the speech."
				
				\begin{description}
					\item[Original Source]
				\end{description}
				
					% brief description of original source material, with reference to JSP or somewhere else for more info
					% Background material: it might be helpful to have a note that says whether a given document was presented as a speech, was written in a journal, etc.
					% Also a comment here about how reliable the source might be, or what shape it is in, if there are large omissions or parts that are hard to understand, and so on.
				
				\begin{description}
					\item[Text]
				\end{description}
				
						\hypertarget{EMS-1841-JS-CIR-16-MAR-1841-WM-a}%
					Next
						\marginnote{a}%
					Meeting--- first Bishop [Vinson] Knight on Kno[w]ledge--- 2 [Elisha] Groves on Punktuality 3 wilson [Law] on Phlathrep 4 on the melenium--- Joseph Said that they wi[c]ked will Not all be Destroyed at the Coming of Christ \& also there will be wiked During the Melenum [Millennium]--- for instance Isaiah says the Days of an infant shall be as the age of a tree also Zaich. [Zechariah] says all who Does Not Come up year by year with their Gifts to the feasts of the tabernicle that No Rain shall fall upon them--- \& that Jesus will be a Resident on the Earth a thousand [years] with the Saints is Not the Case
						\hypertarget{EMS-1841-JS-CIR-16-MAR-1841-WM-b}%
					but
						\marginnote{b}%
					will Raign over the saints \& come Down \& instruct as he Did the 5 hundred Brethern (1\supers{st} Cor.--- 15) \& those of the first Resurerection will also Raign with him over the saints--- then after the Little season is expired \& the Earth underGoes its Last Change \& is Gloryfyed then will all the meek inherit the Earth wherein Dwelleth Righteous--- he says satan cannot seduce us By his enticements unless we in our h[e]arts Consent \& yeald [yield]--- our organization such that we can Resist the Devil If we were Not organized so we would Not be free agents,
					
		% -----
		\subsubsection{Revelation, 20 March 1841}
		% -----
			\label{EMS-1841-JS-20-MAR-1841}
			% \label{EMS-1840-JS-DAY-3LETTERMONTH-YEAR}
			
			\begin{description}
				\item[Print Sources]
			\end{description}

			\begin{tabular}{p{0.5in}p{1in}p{0.5in}p{1in}}
				JSP	 	& ---		& JSR 		& --- \\
				DC	 	& ---		& HC 		& --- \\
				HDDC 	& ---		& TCHC 		& --- \\
			\end{tabular}
			% HDDC = woodford's DC diss; JSR = Marquardt's JS Revelations; TCHC = Vogel HC
			
			\begin{description}
				\item[Online Access] \href{https://www.josephsmithpapers.org/paper-summary/revelation-20march-1841/1}{Here}
				% \item[Analysis] \hyperref[XX]{Here}
			\end{description}
			
			\begin{description}
				\item[Background]
			\end{description}
			
				% It might also be useful to give a two sentence context of the period: e.g., "This speech was given in Nauvoo to a small congregation one month prior to the destruction of the Nauvoo Expositor. These and other tensions may have been on the prophet's mind when he gave the speech."
			
			\begin{description}
				\item[Original Source]
			\end{description}
			
				% brief description of original source material, with reference to JSP or somewhere else for more info
				% Background material: it might be helpful to have a note that says whether a given document was presented as a speech, was written in a journal, etc.
				% Also a comment here about how reliable the source might be, or what shape it is in, if there are large omissions or parts that are hard to understand, and so on.
			
			\begin{description}
				\item[Text]
			\end{description}
			
				\begin{tightright}
						\hypertarget{EMS-1841-JS-20-MAR-1841-a}%
					City
						\marginnote{a}%
					of Nauvoo March [\textit{blank}] 1841
				\end{tightright}

				Brother William Allred, Bishop of the Stake at Pleasent Vale, and also Brother Henry W. Miller President of the Stake at Freedom, desire President Joseph Smith to enquire of the Lord, his will concerning them.

				\begin{tightright}
						\hypertarget{EMS-1841-JS-20-MAR-1841-b}%
					City
						\marginnote{b}%
					of Nauvoo March 20\supersu{th}\supers{.} 1841
				\end{tightright}

				Let my servants William Allred, and Henry W Miller have an agency for the selling of stock for the Nauvoo House, and assist my servants Lyman Wight, Peter Hawes [Haws], George Miller and John Snider in building said house, and let my servants William Allred and Henry W. Miller take stock in the house, that the poor of my people may have employment, and that accomodations may be made for the strangers who shall come to visit this place, and for this purpose let them devote all their properties saith the Lord.
				
		% -----
		\subsubsection{Discourse, circa 21 March 1841}
		% -----
			\label{EMS-1841-JS-CIR-21-MAR-1841}
			% \label{EMS-1840-JS-DAY-3LETTERMONTH-YEAR}
			
			% --
			\paragraph{As Reported by William P. McIntire}
			% --
				\label{EMS-1841-JS-CIR-21-MAR-1841-WM}
			
				\begin{description}
					\item[Print Sources]
				\end{description}

				\begin{tabular}{p{0.5in}p{1in}p{0.5in}p{1in}}
					JSP	 	& ---		& JSR 		& --- \\
					DC	 	& ---		& HC 		& --- \\
					HDDC 	& ---		& TCHC 		& --- \\
				\end{tabular}
				% HDDC = woodford's DC diss; JSR = Marquardt's JS Revelations; TCHC = Vogel HC
				
				\begin{description}
					\item[Online Access] \href{https://www.josephsmithpapers.org/paper-summary/discourse-circa-21march-1841-as-reported-by-williamp-mcintire/1}{Here}
					% \item[Analysis] \hyperref[XX]{Here}
				\end{description}
				
				\begin{description}
					\item[Background]
				\end{description}
				
					% It might also be useful to give a two sentence context of the period: e.g., "This speech was given in Nauvoo to a small congregation one month prior to the destruction of the Nauvoo Expositor. These and other tensions may have been on the prophet's mind when he gave the speech."
				
				\begin{description}
					\item[Original Source]
				\end{description}
				
					% brief description of original source material, with reference to JSP or somewhere else for more info
					% Background material: it might be helpful to have a note that says whether a given document was presented as a speech, was written in a journal, etc.
					% Also a comment here about how reliable the source might be, or what shape it is in, if there are large omissions or parts that are hard to understand, and so on.
				
				\begin{description}
					\item[Text]
				\end{description}
				
						\hypertarget{EMS-1841-JS-CIR-21-MAR-1841-WM-a}%
					Sabath
						\marginnote{a}%
					Meeting--- Joseph Read the 2 \& 3\supersu{d}\supers{.} Chapters of Malichi \& stated that there was a priestHood Confered upon all the sons of Levi throughout the Jeneration of the Jews \& he also said they beCome heirs to that pristHood By Linage or Decent [descent] \& held the Keys of the first principles of the Gospel--- for this--- he quoted concerning Jesus coming to John to be Baptized of John that he (Jesus) might Enter into the Kingdom as John held the Keys suffer it to be so now \&c he also Brought up Zacharias pleading with the Lord in the temple that he might have seed so that the preistHood might be presurved
					
						\hypertarget{EMS-1841-JS-CIR-21-MAR-1841-WM-b}%
					he
						\marginnote{b}%
					also prophesyed that all those that made Light of the Revelations that was Given and him \& his words--- would are [ere] Long Cry \& Lement (when the servent of God would be imprizened) say O!! that we had harkened to the words of God \& the Revelattions Given But all opertunity is Cut of[f] from them
					
			% --
			\paragraph{As Reported by Martha Jane Knowlton Coray}
			% --
				\label{EMS-1841-JS-CIR-21-MAR-1841-MC}
			
				\begin{description}
					\item[Print Sources]
				\end{description}

				\begin{tabular}{p{0.5in}p{1in}p{0.5in}p{1in}}
					JSP	 	& ---		& JSR 		& --- \\
					DC	 	& ---		& HC 		& --- \\
					HDDC 	& ---		& TCHC 		& --- \\
				\end{tabular}
				% HDDC = woodford's DC diss; JSR = Marquardt's JS Revelations; TCHC = Vogel HC
				
				\begin{description}
					\item[Online Access] \href{https://www.josephsmithpapers.org/paper-summary/discourse-circa-21march-1841-as-reported-by-martha-jane-knowlton-coray/1}{Here}
					% \item[Analysis] \hyperref[XX]{Here}
				\end{description}
				
				\begin{description}
					\item[Background]
				\end{description}
				
					% It might also be useful to give a two sentence context of the period: e.g., "This speech was given in Nauvoo to a small congregation one month prior to the destruction of the Nauvoo Expositor. These and other tensions may have been on the prophet's mind when he gave the speech."
				
				\begin{description}
					\item[Original Source]
				\end{description}
				
					% brief description of original source material, with reference to JSP or somewhere else for more info
					% Background material: it might be helpful to have a note that says whether a given document was presented as a speech, was written in a journal, etc.
					% Also a comment here about how reliable the source might be, or what shape it is in, if there are large omissions or parts that are hard to understand, and so on.
				
				\begin{description}
					\item[Text]
				\end{description}
				
					\begin{tightcenter}
							\hypertarget{EMS-1841-JS-CIR-21-MAR-1841-MC-a}%
						Sermon
							\marginnote{a}%
						2\supers{nd}
					\end{tightcenter}

					\sout{on} \sout{the} \sout{death} \sout{of} \sout{Judge} \sout{[Elias]} \sout{Higbee} delivered in the latter part of the winter \sout{about} just before he took Laws store at the House of Bishop [Vinson] Knight year 1841

					Joseph Smith read the 3 chap Malachi dwelt <with emphasis> upon the Levitical priesthood and the promise concerning them

						\hypertarget{EMS-1841-JS-CIR-21-MAR-1841-MC-b}%
					Took
						\marginnote{b}%
					up Jhon [John] shewed him to be a Levite and proceeded The voice of one crying in the wilderness prepare ye the way of the Lord and make his paths strait Now it was written that the priests lips should keep know[l]edge, and to them should the people seek for understanding and above al[l] the law binds them and us to receive the word of the Lord at the hands of the Levites therefore Jhon being Lawful heir to the Levitical Priesthood the people were bound to receive his testimony.
						\hypertarget{EMS-1841-JS-CIR-21-MAR-1841-MC-b}%
					Hence
						\marginnote{b}%
					the saying the Kingdom of heaven suffereth violence and the violent take it by force after Jhon had been testifying of Jesus for some time Jesus came unto him for baptism Jhon felt that the honor of baptizing his master was too great a thing grater than he could claim and said I have need to be baptized of thee and comest thou to me Jesus repl[i]ed thus it behooveth us to fulfill all righteousness thus signifying to Jhon the claim of the Aronic priesthood which holds the keys of enterance into the Kingdom.
						\hypertarget{EMS-1841-JS-CIR-21-MAR-1841-MC-c}%
					Then
						\marginnote{c}%
					the three signs which were given were conclusive The dove which sat upon his shoulder was a sure testimony that he was of God Brethren be not deceived \sout{an} nor doubtful of this fact a spirit of a good man or an angel from heaven who has not a body will never undertake to shake hands with you for he knows you cannot peceive his touch and never will extend his hand but any spirit or body that is attended by a dove you may know to be a pure spirit thus you may is [in] some measure detect \sout{them} the spirits who may come unto you

						\hypertarget{EMS-1841-JS-CIR-21-MAR-1841-MC-d}%
					Jhon
						\marginnote{d}%
					was great in that he baptized Jesus an[d] for this cause Jesus saith ``Among them that are born of \sout{woman} <women> there hath not risen a greater than Jhon the baptist--- But again he says \sout{the} from the \sout{coming} days of Jhon the Batist [John the Baptist] till the Kingdom of Heaven suffereth violence and the violent taketh it by force

					Jesus as I said could not enter except \sout{to} <by> the administration of Jhon. Although Jhon was not a restorer but a forerunner

						\hypertarget{EMS-1841-JS-CIR-21-MAR-1841-MC-e}%
					It
						\marginnote{e}%
					was not the Lawful priests who rejected jesus but the self made priests Those who were preests lawfully received the Saviour in his station which was given him by the Law

					All the Authority that we have is from Jhon The Law is not changed nor the ordinances The keys of ushering into the Kingdom were given to Peter James \& Jhon.

						\hypertarget{EMS-1841-JS-CIR-21-MAR-1841-MC-f}%
					Malachi
						\marginnote{f}%
					4\supers{th} chap And he shall purify the sons of Levi \&c yes brethren the Lord will purify the sons of Levi good or bad for it is through them that blessings flow to Israel and as Israel once was baptized in the cloud and in the sea so shall <God> as a refiners fire and a fullers soap \sout{Purge} purify the sons of Levi and purge them as Gold and as silver \& then and not till then shall the offering of Judah \& Jerusalem be pleasant \sout{into} <unto> the Lord as in days of old and as in former years

						\hypertarget{EMS-1841-JS-CIR-21-MAR-1841-MC-g}%
					And
						\marginnote{g}%
					he shall witness against all iniquity as saith Malachi and shall sorely chastize those who are gone astray, but still he saith I am God I change not therefore ye sons of Jacob are not consumed return thou unto me and I will return unto thee. The Lord will begin by revealing the House of Israel among the gentiles and those who have gone from the ordinances of God shall return unto the keeping of all the law and observing his judgement and statutes to do them Then shall the Law of the Lord go forth from Zion and the word of the Lord to the priests and through them from Jerusalem

						\hypertarget{EMS-1841-JS-CIR-21-MAR-1841-MC-h}%
					And
						\marginnote{h}%
					I prophecy that the day will come when you will say Oh that we had given heed but look now upon our public works the <stone> school house for instance the Simoon [simoom] of the Desert has passed over it the people will not hearken nor hear And bondage Death and destruction are close at our heels The Kingdom will not be broken up but \sout{judgements awaits man} we shall be scattered and driven gathered again \& then dispersed reestablished \& driven abroad and so on until the Ancient of days shall sit and the Kingom and power thereof shall be given to the saints and they shall possess \sout{the} it forever and ever28 which may God hasten for Christs sake Amen
					
		% -----
		\subsubsection{Discourse, circa 28 March 1841}
		% -----
			\label{EMS-1841-JS-CIR-28-MAR-1841}
			% \label{EMS-1840-JS-DAY-3LETTERMONTH-YEAR}
			
			% --
			\paragraph{As Reported by William P. McIntire}
			% --
				\label{EMS-1841-JS-CIR-28-MAR-1841-WM}
			
				\begin{description}
					\item[Print Sources]
				\end{description}

				\begin{tabular}{p{0.5in}p{1in}p{0.5in}p{1in}}
					JSP	 	& ---		& JSR 		& --- \\
					DC	 	& ---		& HC 		& --- \\
					HDDC 	& ---		& TCHC 		& --- \\
				\end{tabular}
				% HDDC = woodford's DC diss; JSR = Marquardt's JS Revelations; TCHC = Vogel HC
				
				\begin{description}
					\item[Online Access] \href{https://www.josephsmithpapers.org/paper-summary/discourse-circa-28march-1841/1}{Here}
					% \item[Analysis] \hyperref[XX]{Here}
				\end{description}
				
				\begin{description}
					\item[Background]
				\end{description}
				
					% It might also be useful to give a two sentence context of the period: e.g., "This speech was given in Nauvoo to a small congregation one month prior to the destruction of the Nauvoo Expositor. These and other tensions may have been on the prophet's mind when he gave the speech."
				
				\begin{description}
					\item[Original Source]
				\end{description}
				
					% brief description of original source material, with reference to JSP or somewhere else for more info
					% Background material: it might be helpful to have a note that says whether a given document was presented as a speech, was written in a journal, etc.
					% Also a comment here about how reliable the source might be, or what shape it is in, if there are large omissions or parts that are hard to understand, and so on.
				
				\begin{description}
					\item[Text]
				\end{description}
				
						\hypertarget{EMS-1841-JS-CIR-28-MAR-1841-WM-a}%
					Next
						\marginnote{a}%
					Meeting Meeting on sunday Joseph Reads the 38\supersu{th}\supers{.} Ch--- of Job. in the book he says is a great Display of human Nature--- it is very Natureal for a man when he sees his fellow man afflicted his Natureal conclusion is that he is suffering the Rath of an angry God \& turn from him in haste not knowing the purposes of God
					
						\hypertarget{EMS-1841-JS-CIR-28-MAR-1841-WM-b}%
					he
						\marginnote{b}%
					says the spirit or the inteligence of men are self Existant principles before the foundation this Earth--- \& quotes the Lords question to Job ``where wast thou when I Laid the foundation of the Earth'' Evedencing that Job was <in> Existing somewhere at that time he says God is Good \& all his acts is for the benefit of infereir [inferior] inteligences---
						\hypertarget{EMS-1841-JS-CIR-28-MAR-1841-WM-c}%
					God
						\marginnote{c}%
					saw that those inteligences had Not power to Defend themselves against those that had a tabernicle therefore the Lord Calls them together in Counsel \& agrees to form them tabernicles so that he might Gender the spirit \& the tabernicle togather so as to create sympathy for their fellow man--- for it is a Natureal thing with those spirits that has the most power to bore down on those of Lesser power so we see the Devil is without a tabernicle \& the Lord has set bands to all Spirits.---
						\hypertarget{EMS-1841-JS-CIR-28-MAR-1841-WM-d}%
					\&
						\marginnote{d}%
					hence came the saying Jesus thou son of David why art thou come to torment us before the time, \& Jesus Comanded him to Come out of the Man \& the Devil besought him that he might enter in a herd of swine Near by (for the Devil knew they were a Coveitous people \& if he could Kill their Hogs that would Drive Jesus out of their coasts \& he then <would> have tabernicle enough) \& Jesus permitted him to Enter into the swine
					
		% -----
		\subsubsection{Account of Meeting, circa 30 March 1841}
		% -----
			\label{EMS-1841-JS-CIR-30-MAR-1841}
			% \label{EMS-1840-JS-DAY-3LETTERMONTH-YEAR}
			
			% --
			\paragraph{As Reported by William P. McIntire}
			% --
				\label{EMS-1841-JS-CIR-30-MAR-1841-WM}
			
				\begin{description}
					\item[Print Sources]
				\end{description}

				\begin{tabular}{p{0.5in}p{1in}p{0.5in}p{1in}}
					JSP	 	& ---		& JSR 		& --- \\
					DC	 	& ---		& HC 		& --- \\
					HDDC 	& ---		& TCHC 		& --- \\
				\end{tabular}
				% HDDC = woodford's DC diss; JSR = Marquardt's JS Revelations; TCHC = Vogel HC
				
				\begin{description}
					\item[Online Access] \href{https://www.josephsmithpapers.org/paper-summary/account-of-meeting-circa-30march-1841/1}{Here}
					% \item[Analysis] \hyperref[XX]{Here}
				\end{description}
				
				\begin{description}
					\item[Background]
				\end{description}
				
					% It might also be useful to give a two sentence context of the period: e.g., "This speech was given in Nauvoo to a small congregation one month prior to the destruction of the Nauvoo Expositor. These and other tensions may have been on the prophet's mind when he gave the speech."
				
				\begin{description}
					\item[Original Source]
				\end{description}
				
					% brief description of original source material, with reference to JSP or somewhere else for more info
					% Background material: it might be helpful to have a note that says whether a given document was presented as a speech, was written in a journal, etc.
					% Also a comment here about how reliable the source might be, or what shape it is in, if there are large omissions or parts that are hard to understand, and so on.
				
				\begin{description}
					\item[Text]
				\end{description}
				
						\hypertarget{EMS-1841-JS-CIR-30-MAR-1841-WM-a}%
					Next
						\marginnote{a}%
					Meeting Lecyum 1\supersu{st}\supers{.}---st[e]ward an Equality in property \& in knoledge Joseph said that an Equality would Not answer for he says if we would equel in property at present in six months we would be worse than ever for there is too many dis[h]onest men amongst us who has more injenity [ingenuity] to cheat the Rest \&c
					
						\hypertarget{EMS-1841-JS-CIR-30-MAR-1841-WM-b}%
					E[benezer]
						\marginnote{b}%
					Robinson spoke on the other Co[m]forter in the 14 \& 16 of John \& that to all man kind for he shall prove the world of sin \& Righteou[s] \& of Judgment \&c Joseph said he would corect in the translation it ought to Read thus \& he shall Remind the world of sin \& of Righteous \& of Judgmen[t] \& this Comforter Reminds of the these things through the servents of the Lord--- But the other Comforter spoken of By John is Jesus himself that is to come \& take up his aboad with them
					
		% -----
		\subsubsection{Benediction, 6 April 1841}
		% -----
			\label{EMS-1841-JS-6-APR-1841}
			% \label{EMS-1840-JS-DAY-3LETTERMONTH-YEAR}
			
			\begin{description}
				\item[Print Sources]
			\end{description}

			\begin{tabular}{p{0.5in}p{1in}p{0.5in}p{1in}}
				JSP	 	& ---		& JSR 		& --- \\
				DC	 	& ---		& HC 		& --- \\
				HDDC 	& ---		& TCHC 		& --- \\
			\end{tabular}
			% HDDC = woodford's DC diss; JSR = Marquardt's JS Revelations; TCHC = Vogel HC
			
			\begin{description}
				\item[Online Access] \href{https://www.josephsmithpapers.org/paper-summary/benediction-6april-1841/1}{Here}
				% \item[Analysis] \hyperref[XX]{Here}
			\end{description}
			
			\begin{description}
				\item[Background]
			\end{description}
			
				% It might also be useful to give a two sentence context of the period: e.g., "This speech was given in Nauvoo to a small congregation one month prior to the destruction of the Nauvoo Expositor. These and other tensions may have been on the prophet's mind when he gave the speech."
			
			\begin{description}
				\item[Original Source]
			\end{description}
			
				% brief description of original source material, with reference to JSP or somewhere else for more info
				% Background material: it might be helpful to have a note that says whether a given document was presented as a speech, was written in a journal, etc.
				% Also a comment here about how reliable the source might be, or what shape it is in, if there are large omissions or parts that are hard to understand, and so on.
			
			\begin{description}
				\item[Text]
			\end{description}
			
				This principal corner stone, in representation of the First Presidency, is now duly laid in honor of the great God; and may it there remain until the whole fabric is completed; and may the same be accomplished speedily; that the Saints may have a place to worship God, and the Son of Man have where to lay his head.
				
		% -----
		\subsubsection{Report of the First Presidency to the Church, circa 7 April 1841}
		% -----
			\label{EMS-1841-JS-CIR-7-APR-1841}
			% \label{EMS-1840-JS-DAY-3LETTERMONTH-YEAR}
			
			\begin{description}
				\item[Print Sources]
			\end{description}

			\begin{tabular}{p{0.5in}p{1in}p{0.5in}p{1in}}
				JSP	 	& ---		& JSR 		& --- \\
				DC	 	& ---		& HC 		& --- \\
				HDDC 	& ---		& TCHC 		& --- \\
			\end{tabular}
			% HDDC = woodford's DC diss; JSR = Marquardt's JS Revelations; TCHC = Vogel HC
			
			\begin{description}
				\item[Online Access] \href{https://www.josephsmithpapers.org/paper-summary/report-of-the-first-presidency-to-the-church-circa-7april-1841/1}{Here}
				% \item[Analysis] \hyperref[XX]{Here}
			\end{description}
			
			\begin{description}
				\item[Background]
			\end{description}
			
				% It might also be useful to give a two sentence context of the period: e.g., "This speech was given in Nauvoo to a small congregation one month prior to the destruction of the Nauvoo Expositor. These and other tensions may have been on the prophet's mind when he gave the speech."
			
			\begin{description}
				\item[Original Source]
			\end{description}
			
				% brief description of original source material, with reference to JSP or somewhere else for more info
				% Background material: it might be helpful to have a note that says whether a given document was presented as a speech, was written in a journal, etc.
				% Also a comment here about how reliable the source might be, or what shape it is in, if there are large omissions or parts that are hard to understand, and so on.
			
			\begin{description}
				\item[Text]
			\end{description}
			
					\hypertarget{EMS-1841-JS-CIR-7-APR-1841-a}%
				The
					\marginnote{a}%
				Presidency of the Church of Jesus Christ of Latter Day Saints, feel great pleasure in assembling with the Saints at another general conference, under circumstances so auspicious and cheering; and with grateful hearts to Almighty God for his providential regard, they cordially unite with the Saints, on this occasion, in ascribing honor, and glory, and blessing to his holy name.

				It is with unfeigned pleasure that they have to make known, the steady and rapid increase of the church in this State, the United States, and in Europe. The anxiety to become acquainted with the principles of the gospel, on every hand, is intense, and the cry of, ``come over and help us,'' is reaching the elders on the wings of every wind, while thousands who have heard the gospel, have become obedient thereto, and are rejoicing in its gifts and blessings.--- Prejudice with its attendant train of evils, is giving way before the force of truth, whose benign rays are penetrating the nations afar off.

					\hypertarget{EMS-1841-JS-CIR-7-APR-1841-b}%
				The
					\marginnote{b}%
				reports from the Twelve in Europe are very satisfactory, and state that the work continues to progress with unparalleled rapidity and that the harvest is truly great.

				In the eastern states, the faithful laborers are successful, and many are flocking to the standard of truth. Nor is the south keeping back---churches have been raised up in the southern and western states, and a very pressing invitation has been received from New Orleans for some of the elders to visit that city, which has been complied with.

				In our own State and immediate neighborhood, many are avowing their attachment to the principles of our holy religion, and have become obedient to the faith.

				Peace and prosperity attend us; and we have favor in the sight of God and virtuous men.

					\hypertarget{EMS-1841-JS-CIR-7-APR-1841-c}%
				The
					\marginnote{c}%
				time was, when we were looked upon as deceivers, and that Mormonism would soon pass away, come to nought, and be forgotten. But the time has gone by when it was looked upon as a trancient matter, or a bubble on the wave, and it is now taking a deep hold in the hearts and affections of all those who are noble minded enough to lay aside the prejudice of education, and investigate the subject with candor and honesty.

					\hypertarget{EMS-1841-JS-CIR-7-APR-1841-d}%
				The
					\marginnote{d}%
				truth, like the sturdy oak, has stood unhurt amid the contending elements, which have beat upon it with tremendous force. The floods have rolled, wave after wave, in quick succession; and have not swallowed it up. ``They have lifted up their voice, O Lord, the floods have lifted up their voice; but the Lord of Hosts is mightier than the mighty waves of the sea.'' Nor, have the flames of persecution, with all the influence of mobs, been able to destroy it; but like Moses' bush it has stood unconsumed, and now at this moment presents an important spectacle both to men and angels.---
					\hypertarget{EMS-1841-JS-CIR-7-APR-1841-e}%
				Where
					\marginnote{e}%
				can we turn our eyes to behold such another? We contemplate a people who have embraced a system of religion unpopular, and the adherence to which has brought upon them repeated persecutions---a people who for their love to God and attachment to his cause, have suffered hunger, nakedness, perils, and almost every privation---a people, who, for the sake of their religion, have had to mour[n] the premature deaths of parents, husbands, wives, and children---a people who have prefered death to slavery and hypocracy, and have honorably maintained their characters, and stood firm and immovable, in times that have tried men's souls.

					\hypertarget{EMS-1841-JS-CIR-7-APR-1841-f}%
				Stand
					\marginnote{f}%
				fast, ye Saints of God, hold on a little while longer, and the storms of life will be past, and you will be rewarded by that God whose servants you are, and who will duly appreciate all your toils and afflictions for Christ's sake and the gospel's. Your names will be handed down to posterity as saints of God, and virtuous men.

				But we hope that those scenes of blood and gore will never more occur, but that many, very many such scenes as the present will be witnessed by the saints, and that in the Temple, the foundation of which has been so happily laid, will the saints of the Most High continue to congregate from year to year, in peace and safety.

					\hypertarget{EMS-1841-JS-CIR-7-APR-1841-g}%
				From
					\marginnote{g}%
				the kind and generous feelings manifest, by the citizens of this State, since our sojourn among them, we may continue to expect the enjoyment of all the blessings of civil and religious liberty, guaranteed by the constitution. The citizens of Illinois have done themselves honor in throwing the mantle of the constitution over a persecuted and afflicted people; and have given evident proof, that they are not only in the enjoyment of the privileges of freemen themselves, but, that they willingly and cheerfully extend that invaluable blessing to others, and that they freely award to faithfulness and virtue their due.

					\hypertarget{EMS-1841-JS-CIR-7-APR-1841-h}%
				The
					\marginnote{h}%
				proceedings of the Legislature in regard to the citizens of this place have been marked with philanthropy and benevolence; and they have laid us under great and lasting obligations, in granting us the several liberal charters we now enjoy, and by which we hope to prosper, until our City becomes the most splendid, our University the most learned, and our Legion the most effective, of any in the Union. In the language of one of our own poets, we would say,

					\begin{quote}
							\hypertarget{EMS-1841-JS-CIR-7-APR-1841-i}%
						In
							\marginnote{i}%
						Illinois we've found a safe retreat,

						A home, a shelter from oppressions dire;

						Where we can worship God as we think right,

						And mobbers come not to disturb our peace;

						Where we can live and hope for better days,

						Enjoy again our liberty, our rights:

						That social intercourse which freedom grants,

						And charity requires of man to man.

						And long may charity pervade each breast,

						And long may \textsc{Illinois} remain the scene

						Of rich prosperity by \textit{peace secured!}
					\end{quote}

					\hypertarget{EMS-1841-JS-CIR-7-APR-1841-j}%
				In
					\marginnote{j}%
				consequence of the impoverished condition of the saints, the buildings which are in progress of erection do not progress as fast as could be desired; but from the interest which is generally manifested by the saints at large, we hope to accomplish much by a combination of effort, and a concentration of action, and erect the Temple and other buildings, which we so much need for our mutual instruction and the education of our children.

				From the reports which have been received, we may expect a large emigration this season. The proclamation which was sent some time ago to the churches abroad, has been responded to, and great numbers are making preparations to come and locate themselves in this city and vicinity.

					\hypertarget{EMS-1841-JS-CIR-7-APR-1841-k}%
				From
					\marginnote{k}%
				what we now witness, we are led to look forward with pleasing anticipation to the future, and soon expect to see the thousands of Israel flocking to this region, in obedience to the heavenly command; numerous habitations of the saints thickly studding the flowery and wide spread prairies of Illinois; temples for the worship of our God erecting in various parts; and great peace resting upon Israel.

				We would call the attention of the saints more particularly to the erection of the Temple, for on its speedy erection great blessings depend. The zeal which is manifested by the saints in this city is indeed praise worthy, and we hope will be imitated by the saints in the various stakes and branches of the church, and that those who cannot contribute labor, will bring their gold and their silver, their brass, and their iron, with the pine tree and box tree, to beautify the same.

					\hypertarget{EMS-1841-JS-CIR-7-APR-1841-l}%
				We
					\marginnote{l}%
				are glad to hear of the organization of the different quorums in this city, and hope that the organization will be attended to in every stake and branch of the church, for the Almighty is a lover of order and good government.

				From the faith and enterprise of the saints generally, we feel greatly encouraged, and cheerfully attend to the important duties devolving upon us, knowing that we not only have the approval of Heaven, but that our efforts for the establishing of Zion and the spread of truth, are cheerfully seconded by the thousands of Israel.

				In conclusion we would say, brethren, be faithful; let your love and moderation be known unto all men; be patient; be mindful to observe all the commandments of your heavenly Father; and the God of all grace shall blesss you, even so, Amen.

				\begin{tightright}
					R[obert] B. THOMPSON, \textit{Clerk}.
				\end{tightright}
				
		% -----
		\subsubsection{Minutes, 7--11 April 1841}
		% -----
			\label{EMS-1841-JS-7-11-APR-1841}
			% \label{EMS-1840-JS-DAY-3LETTERMONTH-YEAR}
			
			\begin{description}
				\item[Print Sources]
			\end{description}

			\begin{tabular}{p{0.5in}p{1in}p{0.5in}p{1in}}
				JSP	 	& ---		& JSR 		& --- \\
				DC	 	& ---		& HC 		& --- \\
				HDDC 	& ---		& TCHC 		& --- \\
			\end{tabular}
			% HDDC = woodford's DC diss; JSR = Marquardt's JS Revelations; TCHC = Vogel HC
			
			\begin{description}
				\item[Online Access] \href{https://www.josephsmithpapers.org/paper-summary/minutes-7-11april-1841/1}{Here}
				% \item[Analysis] \hyperref[XX]{Here}
			\end{description}
			
			\begin{description}
				\item[Background]
			\end{description}
			
				% It might also be useful to give a two sentence context of the period: e.g., "This speech was given in Nauvoo to a small congregation one month prior to the destruction of the Nauvoo Expositor. These and other tensions may have been on the prophet's mind when he gave the speech."
			
			\begin{description}
				\item[Original Source]
			\end{description}
			
				% brief description of original source material, with reference to JSP or somewhere else for more info
				% Background material: it might be helpful to have a note that says whether a given document was presented as a speech, was written in a journal, etc.
				% Also a comment here about how reliable the source might be, or what shape it is in, if there are large omissions or parts that are hard to understand, and so on.
			
			\begin{description}
				\item[Text]
			\end{description}
			
					\hypertarget{EMS-1841-JS-7-11-APR-1841-a}%
				...Pres.
					\marginnote{a}%
				Jos. Smith rose and made some observations in explanation of the same, and likewise of the necessity which existed of building the Temple, that the saiints might have a suitable place for worshiping the Almighty, and also the building of the Nauvoo Boarding House, that suitable accomodations might be afforded for the strangers who might visit this city...
				
				Pres't. J. Smith declared the rule of voting, to be a majority in each quorum, exhorted them to deliberation, faith and prayer, and that they should be strict, and impartial in their examinations. He then told them that the presidents of the different quorums would be presented before them for their acceptance or rejection...
				
					\hypertarget{EMS-1841-JS-7-11-APR-1841-b}%
				Pres't.
					\marginnote{b}%
				Joseph Smith presented the building Committee of the ``House of the Lord,'' to the several quorums collectively, who were unanimously received.
				
				Pres't. Smith observed, that it was necessary that some one should be appointed to fill the quorum of the twelve, in the room of the late Elder David W. Patten, whereupon, Pres't. Rigdon nominated Elder Lyman Wight to that office, which was unanimously accepted. Elder Wight stated, that it was an office of great honor and responsibility, and he felt inadequate to the task, but inasmuch as it was the wish of the authorites of the church, that he should take that office, he would endeavor to magnify it....
				
					\hypertarget{EMS-1841-JS-7-11-APR-1841-c}%
				Pres't.
					\marginnote{c}%
				Smith likewise followed on the same subject, threw considerable light on the doctrine which had been investigated...
				
				Pres'r. J. Smith made some observations respecting the duty of the several quorums, in sending their members into the vineyard, and also stated, that labor on the Temple would be as acceptable to the Lord as preaching in the world.

				Pres't. Smith then stated that it was necessary that some one should be appointed to collect funds for building the Temple...
				
					\hypertarget{EMS-1841-JS-7-11-APR-1841-d}%
				Pres't.
					\marginnote{d}%
				Joseph Smith then addressed the assembly and stated, that in consequence of the severety of the weather, the saints had not received as much instruction as he desired and that some things would have to be laid over until the next conference---as there were many who wished to be baptized, they would now go to the water and give opportunity to any who wished to be baptized of doing so. The procession was then organized and proceeded down to the water.
				
		% -----
		\subsubsection{Discourse, between 6 and 9 April 1841}
		% -----
			\label{EMS-1841-JS-6-9-APR-1841}
			% \label{EMS-1840-JS-DAY-3LETTERMONTH-YEAR}
			
			% --
			\paragraph{As Reported by William P. McIntire}}
			% --
				\label{EMS-1841-JS-6-9-APR-1841-WM}
			
				\begin{description}
					\item[Print Sources]
				\end{description}

				\begin{tabular}{p{0.5in}p{1in}p{0.5in}p{1in}}
					JSP	 	& ---		& JSR 		& --- \\
					DC	 	& ---		& HC 		& --- \\
					HDDC 	& ---		& TCHC 		& --- \\
				\end{tabular}
				% HDDC = woodford's DC diss; JSR = Marquardt's JS Revelations; TCHC = Vogel HC
				
				\begin{description}
					\item[Online Access] \href{https://www.josephsmithpapers.org/paper-summary/discourse-between-6-and-9april-1841-as-reported-by-williamp-mcintire/1}{Here}
					% \item[Analysis] \hyperref[XX]{Here}
				\end{description}
				
				\begin{description}
					\item[Background]
				\end{description}
				
					% It might also be useful to give a two sentence context of the period: e.g., "This speech was given in Nauvoo to a small congregation one month prior to the destruction of the Nauvoo Expositor. These and other tensions may have been on the prophet's mind when he gave the speech."
				
				\begin{description}
					\item[Original Source]
				\end{description}
				
					% brief description of original source material, with reference to JSP or somewhere else for more info
					% Background material: it might be helpful to have a note that says whether a given document was presented as a speech, was written in a journal, etc.
					% Also a comment here about how reliable the source might be, or what shape it is in, if there are large omissions or parts that are hard to understand, and so on.
				
				\begin{description}
					\item[Text]
				\end{description}
				
					Conferance April the 6\supersu{th}\supers{.} 18411 Joseph said to the presedents of the quorums that they should see that No man go out to preach unless their Family, was amply Provided for; and if not provided for stay untill they are, this is Done to put up the Borrs against those that, mearly Join the church \& Get ordained to Get their familys struck in her \& Go out \& Glut themselves while the Church has their familys to Keep.
					
		% -----
		\subsubsection{Letter to John M. Bernhisel, 13 April 1841}
		% -----
			\label{EMS-1841-JS-13-APR-1841}
			% \label{EMS-1840-JS-DAY-3LETTERMONTH-YEAR}
			
			\begin{description}
				\item[Print Sources]
			\end{description}

			\begin{tabular}{p{0.5in}p{1in}p{0.5in}p{1in}}
				JSP	 	& ---		& JSR 		& --- \\
				DC	 	& ---		& HC 		& --- \\
				HDDC 	& ---		& TCHC 		& --- \\
			\end{tabular}
			% HDDC = woodford's DC diss; JSR = Marquardt's JS Revelations; TCHC = Vogel HC
			
			\begin{description}
				\item[Online Access] \href{https://www.josephsmithpapers.org/paper-summary/letter-to-johnm-bernhisel-13april-1841/1}{Here}
				% \item[Analysis] \hyperref[XX]{Here}
			\end{description}
			
			\begin{description}
				\item[Background]
			\end{description}
			
				% It might also be useful to give a two sentence context of the period: e.g., "This speech was given in Nauvoo to a small congregation one month prior to the destruction of the Nauvoo Expositor. These and other tensions may have been on the prophet's mind when he gave the speech."
			
			\begin{description}
				\item[Original Source]
			\end{description}
			
				% brief description of original source material, with reference to JSP or somewhere else for more info
				% Background material: it might be helpful to have a note that says whether a given document was presented as a speech, was written in a journal, etc.
				% Also a comment here about how reliable the source might be, or what shape it is in, if there are large omissions or parts that are hard to understand, and so on.
			
			\begin{description}
				\item[Text]
			\end{description}
			
				...The Brethren in New York wrote to me sometime ago on the subject of Babtizm [baptism] for the dead please to inform them I well attend to it as soon as I possibly can...
				
		% -----
		\subsubsection{Letter to Editors, 6 May 1841}
		% -----
			\label{EMS-1841-JS-6-MAY-1841}
			% \label{EMS-1840-JS-DAY-3LETTERMONTH-YEAR}
			
			\begin{description}
				\item[Print Sources]
			\end{description}

			\begin{tabular}{p{0.5in}p{1in}p{0.5in}p{1in}}
				JSP	 	& ---		& JSR 		& --- \\
				DC	 	& ---		& HC 		& --- \\
				HDDC 	& ---		& TCHC 		& --- \\
			\end{tabular}
			% HDDC = woodford's DC diss; JSR = Marquardt's JS Revelations; TCHC = Vogel HC
			
			\begin{description}
				\item[Online Access] \href{https://www.josephsmithpapers.org/paper-summary/letter-to-editors-6may-1841/1}{Here}
				% \item[Analysis] \hyperref[XX]{Here}
			\end{description}
			
			\begin{description}
				\item[Background]
			\end{description}
			
				% It might also be useful to give a two sentence context of the period: e.g., "This speech was given in Nauvoo to a small congregation one month prior to the destruction of the Nauvoo Expositor. These and other tensions may have been on the prophet's mind when he gave the speech."
			
			\begin{description}
				\item[Original Source]
			\end{description}
			
				% brief description of original source material, with reference to JSP or somewhere else for more info
				% Background material: it might be helpful to have a note that says whether a given document was presented as a speech, was written in a journal, etc.
				% Also a comment here about how reliable the source might be, or what shape it is in, if there are large omissions or parts that are hard to understand, and so on.
			
			\begin{description}
				\item[Text]
			\end{description}
			
				\begin{tightcenter}
						\hypertarget{EMS-1841-JS-6-MAY-1841-a}%
					\phantom{ }
						\marginnote{a}%
					\textit{City of Nauvoo. May 6}, 1841.
				\end{tightcenter}

				\textsc{To the Editors of the Times \& Seasons,}

				Gentlemen:---

				I wish, through the medium of your paper, to make known, that on Sunday last, I had the honor of receiving a visit from the Hon. Stephen A. Douglass [Douglas], Justice of the Supreme Court and Judge of the fifth Judicial Circuit of the State of Illinois, and Cyrus Walker Esq. of Macomb, who expressed great pleasure in visiting our city, and were astonished at the improvements which were made. They were officially introduced to the congregation who had assembled on the meeting ground, by the Mayor; and they severally addressed the assembly.
					\hypertarget{EMS-1841-JS-6-MAY-1841-b}%
				Judge
					\marginnote{b}%
				Douglass, expressed his satisfaction of what he had seen and heard respecting our people and took that opportunity of returning thanks to the citizens of Nauvoo, for confering upon him the freedom of the city, stating that he was not aware of rendering us any service, sufficiently important to deserve such marked honor; and likewise spoke in high terms of our location and the improvements we had made, and that our enterprise and industry were highly creditable to us indeed.

					\hypertarget{EMS-1841-JS-6-MAY-1841-c}%
				Mr.
					\marginnote{c}%
				Walker spoke much in favor of the place, the industry of the citizens \&c. and hoped they would continue to enjoy all the blessings and priveleges of our free and glorious Constitution, and as a patriot and a freeman he was willing at all times to stand boldly in defence of liberty and law.

				It must indeed be satisfactory to this community to know, that kind and generous feelings exist in the hearts of men of such high reputation and moral and intellectual worth.

					\hypertarget{EMS-1841-JS-6-MAY-1841-d}%
				Judge
					\marginnote{d}%
				Douglass has ever proved himself friendly to this people; and interested himself to obtain for us our several charters, holding at that time the office of Secretary of State. Mr. Walker also ranks high, and has long held a standing at the bar, which few attain, and is considered one of the most able and profound jurists in the state.

				The sentiments they expressed on the occasion, were highly honorable to them as American citizens, and as gentlemen.

					\hypertarget{EMS-1841-JS-6-MAY-1841-e}%
				How
					\marginnote{e}%
				different their conduct, from that of the official characters in the state of Missouri, whose minds were prejudiced to such an extent, that instead of mingling in our midst and ascertaining for themselves our character, kept entirely aloof, but were ready at all times to listen to those who had the ``poison of adders under their tongues,'' and who sought our overthrow.

				Let every person who may have inbibed sentiments prejudicial to us, imitate the honorable example of our distinguished visitors, (Douglass \& Walker) and I believe they will find much less to condemn then they anticipated, and probably a great deal to commend.

					\hypertarget{EMS-1841-JS-6-MAY-1841-f}%
				What
					\marginnote{f}%
				makes the late visit more pleasing, is the fact, that Messrs. Douglass \& Walker, have long been held in high estimation as politicians, being champions of the two great parties that exist in the State; but laying aside all party strife, like brothers, citizens, and friends, they mingle with us, mutually disposed to extend to us courtesy, respect and friendship, which I hope, we shall ever be proud to reciprocate.

				I am, very respectfully, yours \&c.

				\begin{tightright}
					JOSEPH SMITH.
				\end{tightright}
				
		% -----
		\subsubsection{Discourse, 16 May 1841}
		% -----
			\label{EMS-1841-JS-16-MAY-1841}
			% \label{EMS-1840-JS-DAY-3LETTERMONTH-YEAR}
			
			% --
			\paragraph{As Published in \textit{Times and Seasons}}
			% --
				\label{EMS-1841-JS-16-MAY-1841-TS}
				
				\begin{description}
					\item[Print Sources]
				\end{description}

				\begin{tabular}{p{0.5in}p{1in}p{0.5in}p{1in}}
					JSP	 	& ---		& JSR 		& --- \\
					DC	 	& ---		& HC 		& --- \\
					HDDC 	& ---		& TCHC 		& --- \\
				\end{tabular}
				% HDDC = woodford's DC diss; JSR = Marquardt's JS Revelations; TCHC = Vogel HC
				
				\begin{description}
					\item[Online Access] \href{https://www.josephsmithpapers.org/paper-summary/discourse-16may-1841-as-published-in-times-and-seasons/1}{Here}
					% \item[Analysis] \hyperref[XX]{Here}
				\end{description}
				
				\begin{description}
					\item[Background]
				\end{description}
				
					% It might also be useful to give a two sentence context of the period: e.g., "This speech was given in Nauvoo to a small congregation one month prior to the destruction of the Nauvoo Expositor. These and other tensions may have been on the prophet's mind when he gave the speech."
				
				\begin{description}
					\item[Original Source]
				\end{description}
				
					% brief description of original source material, with reference to JSP or somewhere else for more info
					% Background material: it might be helpful to have a note that says whether a given document was presented as a speech, was written in a journal, etc.
					% Also a comment here about how reliable the source might be, or what shape it is in, if there are large omissions or parts that are hard to understand, and so on.
				
				\begin{description}
					\item[Text]
				\end{description}
				
						\hypertarget{EMS-1841-JS-16-MAY-1841-TS-a}%
					He
						\marginnote{a}%
					commenced his observations by remarking that the kindness of our Heavenly Father, called for our heartfelt gratitude. He then observed that satan was generally blamed for the evils which we did, but if he was the cause of all our wickedness, men could not be condemned. The devil cannot compel mankind to evil, all was voluntary.--- Those who resist the spirit of God, are liable to be led into temptation, and then the association of heaven is withdrawn from those who refuse to be made partakers of such great glory---God would not exert any compulsory means and the Devil could not; and such ideas as were entertained by many were absurd.
						\hypertarget{EMS-1841-JS-16-MAY-1841-TS-b}%
					The
						\marginnote{b}%
					creature was made subject to vanity, not willingly, but Christ subjected the same in hope---we are all subject to vanity while we travel through the crooked paths, and difficulties which surround us. Where is the man that is free from vanity? None ever were perfect but Jesus, and why was he perfect? because he was the son of God, and had the fulness of the Spirit, and greater power than any man.--- But, notwithstanding our vanity, we look forward with hope, (because ``we are subjected in hope,'') to the time of our deliverance.

						\hypertarget{EMS-1841-JS-16-MAY-1841-TS-c}%
					He
						\marginnote{c}%
					then made some observations on the first principles of the gospel, observing that many of the saints who had come from different States and Nations, had only a very superficial knowledge of these principles, not having heard them fully investigated. He then briefly stated the principles of faith, repentance, and baptism for the remission of sins, which were believed by some of the religious societies of the day, but the doctrine of laying on of hands for the gift of the holy ghost, was discarded by them.

						\hypertarget{EMS-1841-JS-16-MAY-1841-TS-d}%
					The
						\marginnote{d}%
					speaker then referred them to the 6th chap. of Heb. 1. and 2. verses, ``not laying again the foundation of repentance from dead works \&c., but of the doctrines of baptism, laying on of hands, the resurrection and eternal judgment \&c.'' The doctrine of eternal judgment was perfectly understood by the apostle, is evident from several passages of scripture. Peter preached repentance and baptism for the remission of sins to the Jews, who had been led to acts of violence and blood, by their leaders, but to the Rulers he said, ``I would that through ignorance ye did it, as did also \textit{those ye ruled}.''--- Repent, therefore, and be converted that your sins may be blotted out, when the times of refreshing (redemption) shall come from the presence of the Lord, for he shall send Jesus Christ, who before was preached unto you \&c.''
						\hypertarget{EMS-1841-JS-16-MAY-1841-TS-e}%
					The
						\marginnote{e}%
					time of \textit{redemption} here had reference to the time, when Christ should come; then and not till then would their sins be blotted out. Why? Because they were murderers, and no murderer hath eternal life. Even David, must wait for those times of refreshing, before he can come forth and his sins be blotted out; for Peter speaking of him says, ``David hath not yet ascended into Heaven, for his sepulchre is with us to this day:' his remains were then in the tomb. Now we read that many bodies of the saints arose, at Christ's resurrection, probably all the saints, but it seems that David did not. Why? because he had been a murderer.

						\hypertarget{EMS-1841-JS-16-MAY-1841-TS-f}%
					If
						\marginnote{f}%
					the ministers of religion had a proper understanding of the doctrine of eternal judgment, they would not be found attending the man who had forfeited his life to the injured laws of his country by shedding innocent blood; for such characters cannot be forgiven, until they have paid the last farthing. The prayers of all the ministers in the world could never close the gates of hell against a murderer.

						\hypertarget{EMS-1841-JS-16-MAY-1841-TS-g}%
					The
						\marginnote{g}%
					speaker then spoke on the subject of election, and read the 9th chap. in Romans, from which it was evident that the election there spoken of was pertaining to the flesh, and had reference to the seed of Abraham, according to the promise God made to Abraham, saying, ``In thee and in thy seed all the families of the earth shall be blessed.'' To them belonged the adoption, and the covenants \&c. Paul said; when he saw their unbelief I wish myself accursed---according to the flesh---not according to the spirit.

						\hypertarget{EMS-1841-JS-16-MAY-1841-TS-h}%
					Why
						\marginnote{h}%
					did God say to Pharoah, ``for this cause have I raised thee up?'' Because Pharoah was a fit instrument---a wicked man, and had committed acts of cruelty of the most atrocious nature.

					The election of the promised seed still continues, and in the last days, they shall have the priesthood restored unto them, and they shall be the ``Saviors on mount Zion'' the ``ministers of our God,'' if it were not for the remnant which was left, then might we be as Sodom and as Gomorah.

						\hypertarget{EMS-1841-JS-16-MAY-1841-TS-i}%
					The
						\marginnote{i}%
					whole of the chapter had reference to the priesthood and the house of Israel; and unconditional election of individuals to eternal life was not taught by the apostles.

					God did elect or predestinate, that all those who would be saved, should be saved in Christ Jesus, and through obedience to the gospel; but he passes over no man's sins, but visits them with correction, and if his children will not repent of their sins, he will discard them.
				
			% --
			\paragraph{As Reported by Unidentified Scribe}
			% --
				\label{EMS-1841-JS-16-MAY-1841-US}
				
				\begin{description}
					\item[Print Sources]
				\end{description}

				\begin{tabular}{p{0.5in}p{1in}p{0.5in}p{1in}}
					JSP	 	& ---		& JSR 		& --- \\
					DC	 	& ---		& HC 		& --- \\
					HDDC 	& ---		& TCHC 		& --- \\
				\end{tabular}
				% HDDC = woodford's DC diss; JSR = Marquardt's JS Revelations; TCHC = Vogel HC
				
				\begin{description}
					\item[Online Access] \href{https://www.josephsmithpapers.org/paper-summary/discourse-16may-1841-as-reported-by-unidentified-scribe/1}{Here}
					% \item[Analysis] \hyperref[XX]{Here}
				\end{description}
				
				\begin{description}
					\item[Background]
				\end{description}
				
					% It might also be useful to give a two sentence context of the period: e.g., "This speech was given in Nauvoo to a small congregation one month prior to the destruction of the Nauvoo Expositor. These and other tensions may have been on the prophet's mind when he gave the speech."
				
				\begin{description}
					\item[Original Source]
				\end{description}
				
					% brief description of original source material, with reference to JSP or somewhere else for more info
					% Background material: it might be helpful to have a note that says whether a given document was presented as a speech, was written in a journal, etc.
					% Also a comment here about how reliable the source might be, or what shape it is in, if there are large omissions or parts that are hard to understand, and so on.
				
				\begin{description}
					\item[Text]
				\end{description}
				
					\begin{tightcenter}
							\hypertarget{EMS-1841-JS-16-MAY-1841-US-a}%
						Remarks
							\marginnote{a}%
						by Joseph May 16\supersu{th}\supers{.} 1841
					\end{tightcenter}
					
					There are three independant principles; the spirit of God. the spirit of Man. \& the spirit of the Devil. All men have power to resist the Devil. They who have Taber[n]acles have power over those who have not. The docterine of eternal judgement acts 2--- 41 Peter preached repent \& be baptised in the name of jesus Christ for the remision of sins \&c but in act 3 12 he says repent and be converted that your sins my be \sout{forgiven} blotted out when the times of \uline{refreshing} <\uline{redemption}> shall come \& he shall send jesus \&c remission of sins by baptisom was not to be preached to murderers.
						\hypertarget{EMS-1841-JS-16-MAY-1841-US-b}%
					all
						\marginnote{b}%
					the preasts in crisendom might pray for a murderer on the scaffold forever but could not avail so much as a gnat towards their forgiveness; there is no forgiveness for murderers they will have to wait untill the times of redemtion shall come and that in hell. Peter had the Keys of eternal Judgement \& he saw David in hell \& knew for what reason. \& that David would have to remain there untill the reserrection, at the coming of Christ. Rommans 9 All election that can be found in the scriptures is according to the flesh and pertaining to the Preasthood
				
			% --
			\paragraph{As Reported by William Clayton}
			% --
				\label{EMS-1841-JS-16-MAY-1841-WC}
				
				\begin{description}
					\item[Print Sources]
				\end{description}

				\begin{tabular}{p{0.5in}p{1in}p{0.5in}p{1in}}
					JSP	 	& ---		& JSR 		& --- \\
					DC	 	& ---		& HC 		& --- \\
					HDDC 	& ---		& TCHC 		& --- \\
				\end{tabular}
				% HDDC = woodford's DC diss; JSR = Marquardt's JS Revelations; TCHC = Vogel HC
				
				\begin{description}
					\item[Online Access] \href{https://www.josephsmithpapers.org/paper-summary/discourse-16-may-1841-as-reported-by-william-clayton/1}{Here}
					% \item[Analysis] \hyperref[XX]{Here}
				\end{description}
				
				\begin{description}
					\item[Background]
				\end{description}
				
					% It might also be useful to give a two sentence context of the period: e.g., "This speech was given in Nauvoo to a small congregation one month prior to the destruction of the Nauvoo Expositor. These and other tensions may have been on the prophet's mind when he gave the speech."
				
				\begin{description}
					\item[Original Source]
				\end{description}
				
					% brief description of original source material, with reference to JSP or somewhere else for more info
					% Background material: it might be helpful to have a note that says whether a given document was presented as a speech, was written in a journal, etc.
					% Also a comment here about how reliable the source might be, or what shape it is in, if there are large omissions or parts that are hard to understand, and so on.
				
				\begin{description}
					\item[Text]
				\end{description}
				
					\begin{tightcenter}
							\hypertarget{EMS-1841-JS-16-MAY-1841-WC-a}%
						Remarks
							\marginnote{a}%
						by Joseph, May 16\supersu{th}, 1841.
					\end{tightcenter}
					
					There are three independent principles--- the spirit of God, the spirit of man, and the spirit of the devil. All men have power to resist the devil. They who have tabernacles have power over those who have not. The doctrine of eternal judgment Acts 2--- 41 Peter preached repent and be baptized in the name of Jesus Christ for the remission of sins, \&\supersu{c}; but in Acts 3--- 19 he says ``Repent and be converted that your sins may be blotted out when the time of redemption shall come and he shall send Jesus,'' \&c. Remission of sins by baptism was not to be preached to murderers.
						\hypertarget{EMS-1841-JS-16-MAY-1841-WC-b}%
					All
						\marginnote{b}%
					the priests in Christendom might pray for a murderer on the scaffold forever, but could not avail so much as a gnat towards their forgiveness. There is no forgiveness for murderers. They will have to wait until the time of redemption shall come and that in hell. Peter had the keys of eternal judgment and he saw David in hell and knew for what reason, and that David would have to remain there until the resurrection at the coming of Christ. Romans 9.--- all election that can be found in the scripture is according to the flesh and pertaining to the priesthood.
					
		% -----
		\subsubsection{Letter to the Saints Abroad, 24 May 1841}
		% -----
			\label{EMS-1841-JS-24-MAY-1841}
			% \label{EMS-1840-JS-DAY-3LETTERMONTH-YEAR}
			
			\begin{description}
				\item[Print Sources]
			\end{description}

			\begin{tabular}{p{0.5in}p{1in}p{0.5in}p{1in}}
				JSP	 	& ---		& JSR 		& --- \\
				DC	 	& ---		& HC 		& --- \\
				HDDC 	& ---		& TCHC 		& --- \\
			\end{tabular}
			% HDDC = woodford's DC diss; JSR = Marquardt's JS Revelations; TCHC = Vogel HC
			
			\begin{description}
				\item[Online Access] \href{https://www.josephsmithpapers.org/paper-summary/letter-to-the-saints-abroad-24may-1841/1}{Here}
				% \item[Analysis] \hyperref[XX]{Here}
			\end{description}
			
			\begin{description}
				\item[Background]
			\end{description}
			
				% It might also be useful to give a two sentence context of the period: e.g., "This speech was given in Nauvoo to a small congregation one month prior to the destruction of the Nauvoo Expositor. These and other tensions may have been on the prophet's mind when he gave the speech."
			
			\begin{description}
				\item[Original Source]
			\end{description}
			
				% brief description of original source material, with reference to JSP or somewhere else for more info
				% Background material: it might be helpful to have a note that says whether a given document was presented as a speech, was written in a journal, etc.
				% Also a comment here about how reliable the source might be, or what shape it is in, if there are large omissions or parts that are hard to understand, and so on.
			
			\begin{description}
				\item[Text]
			\end{description}
			
				\begin{tightcenter}
					TOvTHE SAINTS ABROAD.
				\end{tightcenter}

				The First Presidency of the Church of Jesus Christ of Latter Day Saints, anxious to promote the prosperity of said church, feel it their duty to call upon the saints who reside out of this county, to make preparations to come in, without delay. This is important, and should be attended to by all who feel an interest in the prosperity of this the corner stone of Zion. Here the Temple must be raised, the University be built, and other edifices erected which are necessary for the great work of the last days; and which can only be done by a concentration of energy, and enterprise. Let it therefore be understood, that all the stakes, excepting those in this county, and in Lee county, Iowa, are discontinued, and the saints instructed to settle in this county as soon as circumstances will permit.

				\begin{tightright}
					JOSEPH SMITH.
				\end{tightright}

				City of Nauvoo, Hancock co., Ill.,

				May 24th 1841.
				
		% -----
		\subsubsection{Letter to Thomas Sharp, 26 May 1841}
		% -----
			\label{EMS-1841-JS-26-MAY-1841}
			% \label{EMS-1840-JS-DAY-3LETTERMONTH-YEAR}
			
			\begin{description}
				\item[Print Sources]
			\end{description}

			\begin{tabular}{p{0.5in}p{1in}p{0.5in}p{1in}}
				JSP	 	& ---		& JSR 		& --- \\
				DC	 	& ---		& HC 		& --- \\
				HDDC 	& ---		& TCHC 		& --- \\
			\end{tabular}
			% HDDC = woodford's DC diss; JSR = Marquardt's JS Revelations; TCHC = Vogel HC
			
			\begin{description}
				\item[Online Access] \href{https://www.josephsmithpapers.org/paper-summary/letter-to-thomas-sharp-26may-1841/1}{Here}
				% \item[Analysis] \hyperref[XX]{Here}
			\end{description}
			
			\begin{description}
				\item[Background]
			\end{description}
			
				% It might also be useful to give a two sentence context of the period: e.g., "This speech was given in Nauvoo to a small congregation one month prior to the destruction of the Nauvoo Expositor. These and other tensions may have been on the prophet's mind when he gave the speech."
			
			\begin{description}
				\item[Original Source]
			\end{description}
			
				% brief description of original source material, with reference to JSP or somewhere else for more info
				% Background material: it might be helpful to have a note that says whether a given document was presented as a speech, was written in a journal, etc.
				% Also a comment here about how reliable the source might be, or what shape it is in, if there are large omissions or parts that are hard to understand, and so on.
			
			\begin{description}
				\item[Text]
			\end{description}
			
				\begin{tightright}
					\textsc{Nauvoo}, Ill, May 26, 1841.
				\end{tightright}
				
				\textit{Mr. [Thomas] Sharp, Editor of the Warsaw Signal:}
				
				\textsc{Sir}---You will discontinue my paper---its contents are calculated to pollute me, and to patronize the filthy sheet---that tissue of lies---that sink of iniquity---is disgraceful to any moral man.
				
				\begin{tightright}
					Yours, with utter contempt,
				
					JOSEPH SMITH.
				\end{tightright}
				
				P. S. Please publish the above in your contemptible paper.
				
				J.S.
				
		% -----
		\subsubsection{Discourse, circa May 1841}
		% -----
			\label{EMS-1841-JS-CIR-MAY-1841}
			% \label{EMS-1840-JS-DAY-3LETTERMONTH-YEAR}
			
			% --
			\paragraph{As Reported By Unidentified Scribe}
			% --
				\label{EMS-1841-JS-CIR-MAY-1841-US}
				
				\begin{description}
					\item[Print Sources]
				\end{description}

				\begin{tabular}{p{0.5in}p{1in}p{0.5in}p{1in}}
					JSP	 	& ---		& JSR 		& --- \\
					DC	 	& ---		& HC 		& --- \\
					HDDC 	& ---		& TCHC 		& --- \\
				\end{tabular}
				% HDDC = woodford's DC diss; JSR = Marquardt's JS Revelations; TCHC = Vogel HC
				
				\begin{description}
					\item[Online Access] \href{https://www.josephsmithpapers.org/paper-summary/discourse-circa-may-1841-as-reported-by-unidentified-scribe/1}{Here}
					% \item[Analysis] \hyperref[XX]{Here}
				\end{description}
				
				\begin{description}
					\item[Background]
				\end{description}
				
					% It might also be useful to give a two sentence context of the period: e.g., "This speech was given in Nauvoo to a small congregation one month prior to the destruction of the Nauvoo Expositor. These and other tensions may have been on the prophet's mind when he gave the speech."
				
				\begin{description}
					\item[Original Source]
				\end{description}
				
					% brief description of original source material, with reference to JSP or somewhere else for more info
					% Background material: it might be helpful to have a note that says whether a given document was presented as a speech, was written in a journal, etc.
					% Also a comment here about how reliable the source might be, or what shape it is in, if there are large omissions or parts that are hard to understand, and so on.
				
				\begin{description}
					\item[Text]
				\end{description}
				
					\begin{tightcenter}
						By Joseph
					\end{tightcenter}
					
					Everlasting covenent was made between 3 personages before the organization of this earth; \& relates to their dispensation of things <to men> on the earth. Theses personages according to Abrahams record are called God the first, the creator God the second, the redeemer, \& God the third, the witness or Testator.---
					
			% --
			\paragraph{As Reported By William Clayton}
			% --
				\label{EMS-1841-JS-CIR-MAY-1841-WC}
				
				\begin{description}
					\item[Print Sources]
				\end{description}

				\begin{tabular}{p{0.5in}p{1in}p{0.5in}p{1in}}
					JSP	 	& ---		& JSR 		& --- \\
					DC	 	& ---		& HC 		& --- \\
					HDDC 	& ---		& TCHC 		& --- \\
				\end{tabular}
				% HDDC = woodford's DC diss; JSR = Marquardt's JS Revelations; TCHC = Vogel HC
				
				\begin{description}
					\item[Online Access] \href{https://www.josephsmithpapers.org/paper-summary/discourse-circa-may-1841-as-reported-by-william-clayton/1}{Here}
					% \item[Analysis] \hyperref[XX]{Here}
				\end{description}
				
				\begin{description}
					\item[Background]
				\end{description}
				
					% It might also be useful to give a two sentence context of the period: e.g., "This speech was given in Nauvoo to a small congregation one month prior to the destruction of the Nauvoo Expositor. These and other tensions may have been on the prophet's mind when he gave the speech."
				
				\begin{description}
					\item[Original Source]
				\end{description}
				
					% brief description of original source material, with reference to JSP or somewhere else for more info
					% Background material: it might be helpful to have a note that says whether a given document was presented as a speech, was written in a journal, etc.
					% Also a comment here about how reliable the source might be, or what shape it is in, if there are large omissions or parts that are hard to understand, and so on.
				
				\begin{description}
					\item[Text]
				\end{description}
				
					\begin{tightcenter}
						By Joseph.
					\end{tightcenter}
					
					Everlasting covenant was made between three personages before the organization of this earth and relates to their dispensation of things to men on the earth. These personages according to Abrahams record are <called> God the first, the Creator; God the second, the Redeemer; and God the third, the Witness or Testator.
					
		% -----
		\subsubsection{Account of Meeting, 3 July 1841}
		% -----
			\label{EMS-1841-JS-3-JUL-1841}
			% \label{EMS-1840-JS-DAY-3LETTERMONTH-YEAR}
				
			\begin{description}
				\item[Print Sources]
			\end{description}

			\begin{tabular}{p{0.5in}p{1in}p{0.5in}p{1in}}
				JSP	 	& ---		& JSR 		& --- \\
				DC	 	& ---		& HC 		& --- \\
				HDDC 	& ---		& TCHC 		& --- \\
			\end{tabular}
			% HDDC = woodford's DC diss; JSR = Marquardt's JS Revelations; TCHC = Vogel HC
			
			\begin{description}
				\item[Online Access] \href{https://www.josephsmithpapers.org/paper-summary/account-of-meeting-3july-1841/1}{Here}
				% \item[Analysis] \hyperref[XX]{Here}
			\end{description}
			
			\begin{description}
				\item[Background]
			\end{description}
			
				% It might also be useful to give a two sentence context of the period: e.g., "This speech was given in Nauvoo to a small congregation one month prior to the destruction of the Nauvoo Expositor. These and other tensions may have been on the prophet's mind when he gave the speech."
			
			\begin{description}
				\item[Original Source]
			\end{description}
			
				% brief description of original source material, with reference to JSP or somewhere else for more info
				% Background material: it might be helpful to have a note that says whether a given document was presented as a speech, was written in a journal, etc.
				% Also a comment here about how reliable the source might be, or what shape it is in, if there are large omissions or parts that are hard to understand, and so on.
			
			\begin{description}
				\item[Text]
			\end{description}
			
				On the the 3\supers{rd} day of July 1841 the Legion was called out to celebrate our National Independance <the 4\supers{th} being Sunday>, and was reviewed by Lieutenant-General Joseph Smith. Who made an Eloquent and patriotic Speech to the troops and strongly testified of his regard for our national welfare and his willingness to lay down his life \sout{for his} in defence of his country \sout{if need be} and closed with these <remarkable> words \sout{emphatically spoken}. ``I would ask no greater boon than to lay down my life for my country.''
				
		% -----
		\subsubsection{Revelation, 9 July 1841 [D\&C 126]}
		% -----
			\label{EMS-1841-JS-9-JUL-1841}
			% \label{EMS-1840-JS-DAY-3LETTERMONTH-YEAR}
				
			\begin{description}
				\item[Print Sources]
			\end{description}

			\begin{tabular}{p{0.5in}p{1in}p{0.5in}p{1in}}
				JSP	 	& ---		& JSR 		& --- \\
				DC	 	& ---		& HC 		& --- \\
				HDDC 	& ---		& TCHC 		& --- \\
			\end{tabular}
			% HDDC = woodford's DC diss; JSR = Marquardt's JS Revelations; TCHC = Vogel HC
			
			\begin{description}
				\item[Online Access] \href{https://www.josephsmithpapers.org/paper-summary/revelation-9july-1841-dc-126/1}{Here}
				% \item[Analysis] \hyperref[XX]{Here}
			\end{description}
			
			\begin{description}
				\item[Background]
			\end{description}
			
				% It might also be useful to give a two sentence context of the period: e.g., "This speech was given in Nauvoo to a small congregation one month prior to the destruction of the Nauvoo Expositor. These and other tensions may have been on the prophet's mind when he gave the speech."
			
			\begin{description}
				\item[Original Source]
			\end{description}
			
				% brief description of original source material, with reference to JSP or somewhere else for more info
				% Background material: it might be helpful to have a note that says whether a given document was presented as a speech, was written in a journal, etc.
				% Also a comment here about how reliable the source might be, or what shape it is in, if there are large omissions or parts that are hard to understand, and so on.
			
			\begin{description}
				\item[Text]
			\end{description}
			
				\begin{tightright}
					``Nauvoo City. July 9\supersu{th}\supers{.} 1841.
				\end{tightright}
				
				Dear \& well beloved Brother, Brigham Young, Verily thus saith the Lord unto you my servant Brigham it is no more required at your hand to leave your family as in times past for your offering is acceptable to me I have seen your labor and toil in journeyings for my name. I therefore command you to send my word abroad and take special care of your family from this time henceforth and forever, Amen. Given to Joseph Smith this day.''
				
		% -----
		\subsubsection{Minutes, 16 August 1841}
		% -----
			\label{EMS-1841-JS-16-AUG-1841}
			% \label{EMS-1840-JS-DAY-3LETTERMONTH-YEAR}
				
			\begin{description}
				\item[Print Sources]
			\end{description}

			\begin{tabular}{p{0.5in}p{1in}p{0.5in}p{1in}}
				JSP	 	& ---		& JSR 		& --- \\
				DC	 	& ---		& HC 		& --- \\
				HDDC 	& ---		& TCHC 		& --- \\
			\end{tabular}
			% HDDC = woodford's DC diss; JSR = Marquardt's JS Revelations; TCHC = Vogel HC
			
			\begin{description}
				\item[Online Access] \href{https://www.josephsmithpapers.org/paper-summary/minutes-16august-1841/1}{Here}
				% \item[Analysis] \hyperref[XX]{Here}
			\end{description}
			
			\begin{description}
				\item[Background]
			\end{description}
			
				% It might also be useful to give a two sentence context of the period: e.g., "This speech was given in Nauvoo to a small congregation one month prior to the destruction of the Nauvoo Expositor. These and other tensions may have been on the prophet's mind when he gave the speech."
			
			\begin{description}
				\item[Original Source]
			\end{description}
			
				% brief description of original source material, with reference to JSP or somewhere else for more info
				% Background material: it might be helpful to have a note that says whether a given document was presented as a speech, was written in a journal, etc.
				% Also a comment here about how reliable the source might be, or what shape it is in, if there are large omissions or parts that are hard to understand, and so on.
			
			\begin{description}
				\item[Text]
			\end{description}
			
					\hypertarget{EMS-1841-JS-16-AUG-1841-a}%
				...Met
					\marginnote{a}%
				at 2 P.M. When the conference \sout{were} was address'd by Elders Lorenzo Barnes and H. G. Sherwood on the subject of preaching the gospel and building up of the kingdom of God in theese last days

				President Joseph Smith (who had been absent in consequence of the death of his child during the former part of the Day) \sout{being present} <on his arrival> then addressed the conference on the objects of <calling> \sout{the} a conference at this time and in addition to what had been stated by president Young said that \sout{some} the twelve should be authorized <wh> to assist in managing the affairs of th[e] kingdom in this place. which he said was the\sout{ir} duties of their office \&c.

					\hypertarget{EMS-1841-JS-16-AUG-1841-b}%
				Motioned
					\marginnote{b}%
				seconded and Carried \sout{the} that the quorum of the twelve be authorized to act in \sout{building} \sout{the} accordance with the instructions given by president Joseph Smith in regulating and superintending the affairs <of the Church.>

				Motioned, seconded \& carried unanimou[s]ly that every individual who shall hereafter be found trying to influence any emigrants \sout{coming} b[e]longing to the Church to either buy or sell property of them or to them (except the agent) \sout{such person} shall be immediately tried for fellowship and dealt with as offenders and unless they repent they shall be cut off from the Chu[rc]h

					\hypertarget{EMS-1841-JS-16-AUG-1841-c}%
				President
					\marginnote{c}%
				\phantom{ }
				[Sidney] Rigdon then made some appropriate remarks on speculation \&c

				Resolved on Motion \sout{of President Joseph Smith} Res. that the twelve be authorised to make the selection of elders independent of the conference and present them \sout{to the first Pres. him} to President Joseph Smith for <his> approval

				Resolved that this conference be adjourned to the time of the general Conference in Oct next

				Closed by singing and prayer by President You[n]g
				
		% -----
		\subsubsection{Discourse, 16 August 1841}
		% -----
			\label{EMS-1841-JS-16-AUG-1841-D}
			% \label{EMS-1840-JS-DAY-3LETTERMONTH-YEAR}
			
			% --
			\paragraph{As Published in Times and Seasons}
			% --
				\label{EMS-1841-JS-16-AUG-1841-D-TS}
				
				\begin{description}
					\item[Print Sources]	
				\end{description}

				\begin{tabular}{p{0.5in}p{1in}p{0.5in}p{1in}}
					JSP	 	& ---		& JSR 		& --- \\
					DC	 	& ---		& HC 		& --- \\
					HDDC 	& ---		& TCHC 		& --- \\
				\end{tabular}
				% HDDC = woodford's DC diss; JSR = Marquardt's JS Revelations; TCHC = Vogel HC
				
				\begin{description}
					\item[Online Access] \href{https://www.josephsmithpapers.org/paper-summary/discourse-16august-1841-as-published-in-times-and-seasons/1}{Here}
					% \item[Analysis] \hyperref[XX]{Here}
				\end{description}
				
				\begin{description}
					\item[Background]
				\end{description}
				
					% It might also be useful to give a two sentence context of the period: e.g., "This speech was given in Nauvoo to a small congregation one month prior to the destruction of the Nauvoo Expositor. These and other tensions may have been on the prophet's mind when he gave the speech."
				
				\begin{description}
					\item[Original Source]
				\end{description}
				
					% brief description of original source material, with reference to JSP or somewhere else for more info
					% Background material: it might be helpful to have a note that says whether a given document was presented as a speech, was written in a journal, etc.
					% Also a comment here about how reliable the source might be, or what shape it is in, if there are large omissions or parts that are hard to understand, and so on.
				
				\begin{description}
					\item[Text]
				\end{description}
				
						\hypertarget{EMS-1841-JS-16-AUG-1841-D-TS-a}%
					President
						\marginnote{a}%
					Joseph Smith now arriving proceeded to state to the conference at considerable length, the object of their present meeting, and in addition to what President [Brigham] Young had stated in the morning, said that the time had come when the twelve should be called upon to stand in their place next to the first presidency, and attend to the settling of emegrants and the business of the church at the stakes, and assist to bear off the kingdom victorious to the nations;
						\hypertarget{EMS-1841-JS-16-AUG-1841-D-TS-b}%
					and
						\marginnote{b}%
					as they had been faithful and had borne the burden in the heat of the day that it was right that they should have an opportunity of provididing something for themselves and families, and at the same time relieve him so that he might attend to the businesss of translating.
					
		% -----
		\subsubsection{Discourse, 22 August 1841}
		% -----
			\label{EMS-1841-JS-22-AUG-1841}
			% \label{EMS-1840-JS-DAY-3LETTERMONTH-YEAR}
				
			\begin{description}
				\item[Print Sources]
			\end{description}

			\begin{tabular}{p{0.5in}p{1in}p{0.5in}p{1in}}
				JSP	 	& ---		& JSR 		& --- \\
				DC	 	& ---		& HC 		& --- \\
				HDDC 	& ---		& TCHC 		& --- \\
			\end{tabular}
			% HDDC = woodford's DC diss; JSR = Marquardt's JS Revelations; TCHC = Vogel HC
			
			\begin{description}
				\item[Online Access] \href{https://www.josephsmithpapers.org/paper-summary/discourse-22-august-1841/1}{Here}
				% \item[Analysis] \hyperref[XX]{Here}
			\end{description}
			
			\begin{description}
				\item[Background]
			\end{description}
			
				% It might also be useful to give a two sentence context of the period: e.g., "This speech was given in Nauvoo to a small congregation one month prior to the destruction of the Nauvoo Expositor. These and other tensions may have been on the prophet's mind when he gave the speech."
			
			\begin{description}
				\item[Original Source]
			\end{description}
			
				% brief description of original source material, with reference to JSP or somewhere else for more info
				% Background material: it might be helpful to have a note that says whether a given document was presented as a speech, was written in a journal, etc.
				% Also a comment here about how reliable the source might be, or what shape it is in, if there are large omissions or parts that are hard to understand, and so on.
			
			\begin{description}
				\item[Text]
			\end{description}
			
					\hypertarget{EMS-1841-JS-22-AUG-1841-a}%
				My
					\marginnote{a}%
				purpose was to have given you a few extracts, or rather specimens, of Jo's sermon; but my brief space will scarcely admit of it in this epistle. It would be a perfect \textit{curiosity}, I'll assure you. He disclaimed all ambition or vain glory; but I think it evident he possesses a sufficiency. In attempting to exhibit to the `brethren' how much more ready the world is to believe error than truth, `Jo' often made himself the `hero' of the text, with not a little self-complacency either, notwithstanding his \textit{professed} humility.
					\hypertarget{EMS-1841-JS-22-AUG-1841-b}%
				`Brethren,'
					\marginnote{b}%
				says he, `when Moses lived, that we read about, there was not half so many rascally lying editors as there is now-a-days. Old Moses was commanded of God to slay a man, but no such requirement is made of `Old Jo Smith, O no, not as long as the lying editors continue their abuse, the world is ready to swallow it.''--
				
					\hypertarget{EMS-1841-JS-22-AUG-1841-c}%
				He
					\marginnote{c}%
				then entertained his hearers by relating an anecdote, a case in point. At the time he was traveling in the stage from Missouri to Kirkland [Kirtland], Ohio, how a young \textit{upstart} of a lawyer told a `\textit{long yarn}' about the death of `Jo Smith'--dead, buried--having got shot; and how the `brethren' held a pow wow over the dead body; but it availed not; Joe would die because he was not strong in the faith.
					\hypertarget{EMS-1841-JS-22-AUG-1841-d}%
				This
					\marginnote{d}%
				was his own version. In concluding this story, which was quite ingenious, he says `Well brethren,' what do you think? Do you suppose the passengers in that stage believed me, or do you suppose they believed the lying language? Well, I'll tell you they believed--they thought and said I was a gross impostor! and for three days continuous traveling, I could not make them believe that I was the veritable Jo Smith;%
					\hypertarget{EMS-1841-JS-22-AUG-1841-e}%
				--thus
					\marginnote{e}%
				it is `brethren,' says he, mankind are prone to receive error rather than truth, and so they believe the thousands of falsehoods circulated about our religion, without coming to see and examine for themselves. Well, lie away, ye editors, I see ye are determined to make a \textit{big} man of me anyhow, but I will leave the Prophet's sermon.
				
		% -----
		\subsubsection{Letter to Horace Hotchkiss, 25 August 1841}
		% -----
			\label{EMS-1841-JS-25-AUG-1841}
			% \label{EMS-1840-JS-DAY-3LETTERMONTH-YEAR}
				
			\begin{description}
				\item[Print Sources]
			\end{description}

			\begin{tabular}{p{0.5in}p{1in}p{0.5in}p{1in}}
				JSP	 	& ---		& JSR 		& --- \\
				DC	 	& ---		& HC 		& --- \\
				HDDC 	& ---		& TCHC 		& --- \\
			\end{tabular}
			% HDDC = woodford's DC diss; JSR = Marquardt's JS Revelations; TCHC = Vogel HC
			
			\begin{description}
				\item[Online Access] \href{https://www.josephsmithpapers.org/paper-summary/letter-to-horace-hotchkiss-25august-1841/1}{Here}
				% \item[Analysis] \hyperref[XX]{Here}
			\end{description}
			
			\begin{description}
				\item[Background]
			\end{description}
			
				% It might also be useful to give a two sentence context of the period: e.g., "This speech was given in Nauvoo to a small congregation one month prior to the destruction of the Nauvoo Expositor. These and other tensions may have been on the prophet's mind when he gave the speech."
			
			\begin{description}
				\item[Original Source]
			\end{description}
			
				% brief description of original source material, with reference to JSP or somewhere else for more info
				% Background material: it might be helpful to have a note that says whether a given document was presented as a speech, was written in a journal, etc.
				% Also a comment here about how reliable the source might be, or what shape it is in, if there are large omissions or parts that are hard to understand, and so on.
			
			\begin{description}
				\item[Text]
			\end{description}
			
				\begin{tightright}
						\hypertarget{EMS-1841-JS-25-AUG-1841-a}%
					Nauvoo
						\marginnote{a}%
					August 25\supers{th.} 1841
				\end{tightright}

				Horace R Hotchkiss Esqr

				\begin{tightcenter}
					Dr Sir
				\end{tightcenter}

				Yours of the 24\supers{th} ult\supers{o.}, came to hand this day, The contents of which I duly appreciate I presume you are well aware of the difficulties that occurred before, and at, The execution of the writings in regard of the landed transactions between us, touching the annual payments of the interest, If you have forgotten I will here remind you; You verbally agreed, on our refusal, and hesitancy to execute the notes for the payment of the Land, That you would not exact the payment, of the interest that would accrue on them under five years, and that you would not coerce the payment even then,
					\hypertarget{EMS-1841-JS-25-AUG-1841-b}%
				To
					\marginnote{b}%
				all this you pledged your honor, and upon an after arrangement, you verbally agreed, to take Land in some one of the Atlantic States that would yield Six per cent interest (to you,) both for the principal and interest, and in view of that matter I deligated my Bro. Hyrum [Smith] \& Doctor Isaac Galland to go east and negociate for Lands, with our friends, and pay you off for the whole purchase that we made of you, But upon an interview with you they learned that you were unwilling to enter into an arrangement according to the powers that I had deligated to them, That you would not receive any of the principal at all, but the interest alone, which we never considered ourselves, in \uline{honor} or in \uline{justice} bound to pay under the expiration of five years,
					\hypertarget{EMS-1841-JS-25-AUG-1841-c}%
				I
					\marginnote{c}%
				presumed you are no stranger, to the part of the City plat we bought of you, it being a \uline{deathly} \uline{sickly} hole, and that we have not been able in consequence to realize any valuable consideration from it, although we have been keeping up appearances, and holding out inducements, to encourage emigration that we scarcely think justifiable in consequence of the mortality, that almost invariably awaits those who come from far distant parts, <and> \sout{at} that with a view to enable us to meet our engagements, And now to be goaded by you for a breach of good faith and neglect, and dishonorable conduct seems to me to be almost beyond endurance, you are aware that we came from Missouri destitute of every thing but physical form, and had nothing but our energies and perseverence to rely upon to meet the payment of the \uline{extortionate} sum, that you exacted for the Land that we had of you,
					\hypertarget{EMS-1841-JS-25-AUG-1841-d}%
				Have
					\marginnote{d}%
				you no feelings of commisseration, or is it your design to crush us, with a ponderous load before we are able to walk, or can you better dispose of the property than we are doing it. for your interest? If so, to the alternative; I therefore propose in order to avoid the perplexity and annoyance that has hitherto attended the transaction, that you come and take the premises and make the best you can of it. Or stand off and give us an oppertunity, that we may manage the concern, and enable ourselves by the management thereof to meet our engagement as was originally contemplated. We have taken a city plat at Warsaw at the head of navigation for vessels of heavy towage on the most advantagious terms,
					\hypertarget{EMS-1841-JS-25-AUG-1841-e}%
				The
					\marginnote{e}%
				proprietors waiting on us for the payment of the plat until we can reallize the money from the Sales, leaving to ourselves, a large and liberal net proffit. We have been making every exertion and used all the means at our command [to lay] a foundation that will now begin to enable us to meet our pecuniary engagements, and no doubt in our minds, to the entire satisfaction of all those concerned, if they will but exercise a small degree of patience, and stay a resort to Coersive Measures, which would kill us in the germ, even before we can (by reason of the season) begin to bud and blossom, in order to bring forth a plentiful yield of fruit.

				\begin{tightright}
						\hypertarget{EMS-1841-JS-25-AUG-1841-f}%
					I
						\marginnote{f}%
					am with considerations of high respect your Ob\supers{t.} Serv\supers{t.}

					Joseph Smith
				\end{tightright}

				P.S. Since writing the above, I have had a conference with my Bro. Hyrum, who informs me that when he left Pennsylvania that he left with Do\supers{c.} Galland nearly enough of real estate (in the hands of the Doctor) to liquidate the am\supers{t.} due you, my Bro. having been compelled to return, in consequence of ill heal[t]h, expecting that the doctor would have the matter arranged long before this. Therefore so soon as we learn the particulars from D\supers{r} Galland we will take such measures, as will most likely meet your approval,
					\hypertarget{EMS-1841-JS-25-AUG-1841-g}%
				We
					\marginnote{g}%
				have learned that Doctor Galland has been partially blind, which may be the reason why the business has not been arranged as stipulated between you and them. as, he, was to remain, with a view to accomplish that object, after my Bro. returned.--- I will now give you an account of some recent \uline{deaths} that have taken place here, My Brother D. C. [Don Carlos] Smith, Robert B Thompson. and one of my sons, together with many other valuable citizens--- In fact we are in the midst of death--- yours

				\begin{tightright}
					JS---

					<25>
				\end{tightright}

				<NAUVOO Ill AUG 28>

				Horace R. Hotchkiss Esqr

				New Haven

				Con.
				
		% -----
		\subsubsection{Proclamation, between 19 January and 27 August 1841}
		% -----
			\label{EMS-1841-JS-19-JAN-27-AUG-1841}
			% \label{EMS-1840-JS-DAY-3LETTERMONTH-YEAR}
				
			\begin{description}
				\item[Print Sources]
			\end{description}

			\begin{tabular}{p{0.5in}p{1in}p{0.5in}p{1in}}
				JSP	 	& ---		& JSR 		& --- \\
				DC	 	& ---		& HC 		& --- \\
				HDDC 	& ---		& TCHC 		& --- \\
			\end{tabular}
			% HDDC = woodford's DC diss; JSR = Marquardt's JS Revelations; TCHC = Vogel HC
			
			\begin{description}
				\item[Online Access] \href{https://www.josephsmithpapers.org/paper-summary/proclamation-between-19january-and-27august-1841/1}{Here}
				% \item[Analysis] \hyperref[XX]{Here}
			\end{description}
			
			\begin{description}
				\item[Background]
			\end{description}
			
				% It might also be useful to give a two sentence context of the period: e.g., "This speech was given in Nauvoo to a small congregation one month prior to the destruction of the Nauvoo Expositor. These and other tensions may have been on the prophet's mind when he gave the speech."
			
			\begin{description}
				\item[Original Source]
			\end{description}
			
				% brief description of original source material, with reference to JSP or somewhere else for more info
				% Background material: it might be helpful to have a note that says whether a given document was presented as a speech, was written in a journal, etc.
				% Also a comment here about how reliable the source might be, or what shape it is in, if there are large omissions or parts that are hard to understand, and so on.
			
			\begin{description}
				\item[Text]
			\end{description}
			
				\begin{tightcenter}
						\hypertarget{EMS-1841-JS-19-JAN-27-AUG-1841-a}%
					A
						\marginnote{a}%
					Religious Proclamation
				\end{tightcenter}

				From Joseph Smith, President of the Church of Jesus Christ of Latter Day Saints, and Prophet, Seer, and Revelator of the Most High God, to the President of the United States of North America--- the Governors of the several States--- the Emperors, Kings, and Princes of the earth--- the Executives of all nations--- the Chiefs of all tribes--- and all occupying high places in the administration of governments.

					\hypertarget{EMS-1841-JS-19-JAN-27-AUG-1841-b}%
				``Thus
					\marginnote{b}%
				saith the \uline{Lord God}, Behold I will lift up mine hand to the Gentiles, and set up my standard to the people: and they shall bring thy sons in \uline{their} arms, and thy daughters shall be carried upon \uline{their} shoulders. And kings shall be thy nursing-fathers, and their queens thy nursing-mothers: they shall bow down to thee with \uline{their} face toward the earth, and lick up the dust of thy feet: and thou shalt know that I \uline{am} the \uline{Lord}: for they shall not be ashamed that wait for me.''

				(Isaiah--- LX, LXI, LXII.)

					\hypertarget{EMS-1841-JS-19-JAN-27-AUG-1841-c}%
				Now
					\marginnote{c}%
				in obedience to a revelation given January 19th, [A]D. 1841, I proceed to call upon you to yield yourselves as obedient subjects to the requisitions of heaven, in \sout{fulfilling the} contributing to the fulfilment of the predictions of the prophets---to believe on the Lord Jesus Christ, repent of and abandon your sins, be immersed for the remission of your sins, receive to imposition of hands for the gift of the Holy Ghost, and, in fine, to embrace the gospel in its beauty \& fulness.
					\hypertarget{EMS-1841-JS-19-JAN-27-AUG-1841-d}%
				``And
					\marginnote{d}%
				now, why tarriest thou? arise, and be baptised, and wash away thy sins, calling on th[e] name of the Lord.'' This is more honorable than the diadems of kings. or th[e] pearls of princes--- it will confer upon you an excellence unsurpassed, a glory which you never knew. ``God came from Teman, and the Holy one from mount Paran;'' and, as his glory covered the heavens, so now will it environ the earth, and illumine the holy city, until all the obedient shall \sout{bask in the} rejoice in the brightness of his coming, and bask in the sun-shine of God's benignity.
					\hypertarget{EMS-1841-JS-19-JAN-27-AUG-1841-e}%
				Hasten
					\marginnote{e}%
				then to Zion! and contribute to the erection of temples, sanctuaries, and palaces, such as this world never saw--- \sout{decorated with gold and pearls,} \sout{and precious stones}, with their walls finished with th[e] pencil of Raphahel, decorated with gold, and pearls, and precious stones, beautified by the finger of God. Tho' your minds are yet darkened, and your eyes dim of sight, by the traditions, superstitions, and follies of the age, imposed upon you by the Papal See, and hierarchy, of Rome; th[e] Patriarch, and \sout{council} ecclesiastical council, of Constantinople; and th[e] priesthood of th[e] protestant sects;
					\hypertarget{EMS-1841-JS-19-JAN-27-AUG-1841-f}%
				the
					\marginnote{f}%
				God of heaven addresses you as intelligent beings, and directs you to come out from among them, that you may become the \uline{elite} of the kingdom--- bright, and shining lights in your Father's house.

				[\textit{page [3] blank}]

				[\textit{page [4] blank}]
				
		% -----
		\subsubsection{Letter to Oliver Granger, 30 August 1841}
		% -----
			\label{EMS-1841-JS-30-AUG-1841}
			% \label{EMS-1840-JS-DAY-3LETTERMONTH-YEAR}
				
			\begin{description}
				\item[Print Sources]
			\end{description}	

			\begin{tabular}{p{0.5in}p{1in}p{0.5in}p{1in}}
				JSP	 	& ---		& JSR 		& --- \\
				DC	 	& ---		& HC 		& --- \\
				HDDC 	& ---		& TCHC 		& --- \\
			\end{tabular}
			% HDDC = woodford's DC diss; JSR = Marquardt's JS Revelations; TCHC = Vogel HC
			
			\begin{description}
				\item[Online Access] \href{https://www.josephsmithpapers.org/paper-summary/letter-to-oliver-granger-30august-1841/1}{Here}
				% \item[Analysis] \hyperref[XX]{Here}
			\end{description}
			
			\begin{description}
				\item[Background]
			\end{description}
			
				% It might also be useful to give a two sentence context of the period: e.g., "This speech was given in Nauvoo to a small congregation one month prior to the destruction of the Nauvoo Expositor. These and other tensions may have been on the prophet's mind when he gave the speech."
			
			\begin{description}
				\item[Original Source]
			\end{description}
			
				% brief description of original source material, with reference to JSP or somewhere else for more info
				% Background material: it might be helpful to have a note that says whether a given document was presented as a speech, was written in a journal, etc.
				% Also a comment here about how reliable the source might be, or what shape it is in, if there are large omissions or parts that are hard to understand, and so on.
			
			\begin{description}
				\item[Text]
			\end{description}
			
				\begin{tightcenter}
						\hypertarget{EMS-1841-JS-30-AUG-1841-a}%
					Nauvoo
						\marginnote{a}%
					Agust 30th 1841
				\end{tightcenter}

				Brother O[liver] Granger Dear Sir

				We was cal[le]d upon a few Days since by a Gentleman from Newyork Citty by the Name of Devenport with a Judgement In favor of [Ray] \uline{Boynton} \& [Harry] \uline{Hyde} of that Place \& as I Had no recollection of the Debt I cited him to you beleiving you would understand all a bout it for the Judgement was rendered at Shardon [Chardon]

					\hypertarget{EMS-1841-JS-30-AUG-1841-b}%
				Mr
					\marginnote{b}%
				Devenport appeared to be qite a Gentleman \& said he would call, on you \& If there is any means Possible to Settle it we wish you to do so if the debt is a Just one he will be willing to take any Kind of property in a most any place on the Earth--- all things Prosper in this Place Except the Loss we have sustaind in the death of two of our most valuable men Brother D C Smth [Don Carlos Smith] \& Robert B Thompson
					\hypertarget{EMS-1841-JS-30-AUG-1841-c}%
				Both
					\marginnote{c}%
				have recently Died of what I call a qick Consumption theyr Desease was upon theyr Lungs they wasted a way in one week \& Spit up theyr verry vitals--- they are gone. theyr Loss is Irreparable but we must be submisive to the will of god \& try to Stand in our Lot bothe \sout{wow} <now> \& at this End we have not heard from you of Late we would be glad to hear from you often Let us Know how you Prosper in all your\sout{s} affairs \& ours also

					\hypertarget{EMS-1841-JS-30-AUG-1841-d}%
				I
					\marginnote{d}%
				Saw Sister Agness [Agnes Coolbrith Smith] Carloss widow <this morning> She wished me to say to you She hoped you would remember the widow \& fartherless \& be shure to deed the House \& Lot in Kirtland to her for it is all She has it is like the widows Mite She is Left qite dest[i]tute we wish you to do so if you have not. we will be responsible for the moneys you have Expended for the Place we have recently received a letter from Brother [William W.] Phelps reqesting Council\sout{l} I gave him Council when I was there Last March
					\hypertarget{EMS-1841-JS-30-AUG-1841-e}%
				I
					\marginnote{e}%
				think that if he had had as \sout{mall} mutch faith in my Council as he had in Brother A Babbits [Almon Babbitt's] when he came home from Pheladelphia he would not have wanted any further information on the Subject the commandments \& revelations are before him the Proclamation to the Church to gether with the reqest of the first Presidency \& the word of the Lord saying there remained a \sout{Stourge} Scourge for the People \sout{of} or for the in hab[i]tants of Kirtland \uline{\& C}

					\hypertarget{EMS-1841-JS-30-AUG-1841-f}%
				Hiram
					\marginnote{f}%
				Kimball is qite sick at this time we think he is some better on the mend Sarah [Granger Kimball] is qite smart the general health of <the> Place <is> as good as can be Expectted considering the Long <spell of> Dry weather there has been no rain here for three month to \sout{wit} wet the ground in two Inches the Inglish grain \& corn crops are good but the Potatoes are good for nothing the Lords House \& the Nauvoo House Prospers beyond Expectation

				My family \& Brother Josephs are well at Present there is but a verry few Sick in this Place at this time

					\hypertarget{EMS-1841-JS-30-AUG-1841-g}%
				Brother
					\marginnote{g}%
				Granger this Is an Eventful Period a Day of Darkness \& of thick darkness \& Mourning \sout{\&} weeping \& Lemantation \& what Ever the saints find to do Let them do it qickly \& above all things Else Let them strive to be Prepard to die---

				Please write to us on the receipt of this \& Let us Know all the Particulars of all the affairs Brother Joseph sends his Love to you \& famaly My Love to all \&.C.

				yours in the Bonds of the Everlasting Covenant

				\begin{tightright}
					Hyrum Smith

					Joseph Smith
				\end{tightright}

				[\textit{page [3] blank}]

				\begin{tightright}
					<25>
				\end{tightright}

				<NAUVOO Ill. SEP 4>

				Oliver Granger \uline{Esq}

				Kirtland Lake C\supersu{o}\supers{.}

				\uline{Ohio}
				
		% -----
		\subsubsection{Minutes and Discourse, 1--5 October 1841}
		% -----
			\label{EMS-1841-JS-1-5-OCT-1841}
			% \label{EMS-1840-JS-DAY-3LETTERMONTH-YEAR}
			
			% --
			\paragraph{As Published in \textit{Times and Seasons}}
			% --
				\label{EMS-1841-JS-1-5-OCT-1841-TS}
				
				\begin{description}
					\item[Print Sources]
				\end{description}	

				\begin{tabular}{p{0.5in}p{1in}p{0.5in}p{1in}}
					JSP	 	& ---		& JSR 		& --- \\
					DC	 	& ---		& HC 		& --- \\
					HDDC 	& ---		& TCHC 		& --- \\
				\end{tabular}
				% HDDC = woodford's DC diss; JSR = Marquardt's JS Revelations; TCHC = Vogel HC
				
				\begin{description}
					\item[Online Access] \href{https://www.josephsmithpapers.org/paper-summary/minutes-and-discourse-1-5october-1841/1}{Here}
					% \item[Analysis] \hyperref[XX]{Here}
				\end{description}
				
				\begin{description}
					\item[Background]
				\end{description}
				
					% It might also be useful to give a two sentence context of the period: e.g., "This speech was given in Nauvoo to a small congregation one month prior to the destruction of the Nauvoo Expositor. These and other tensions may have been on the prophet's mind when he gave the speech."
				
				\begin{description}
					\item[Original Source]
				\end{description}
				
					% brief description of original source material, with reference to JSP or somewhere else for more info
					% Background material: it might be helpful to have a note that says whether a given document was presented as a speech, was written in a journal, etc.
					% Also a comment here about how reliable the source might be, or what shape it is in, if there are large omissions or parts that are hard to understand, and so on.
				
				\begin{description}
					\item[Text]
				\end{description}
				
						\hypertarget{EMS-1841-JS-1-5-OCT-1841-TS-a}%
					...P. M.
						\marginnote{a}%
					Pres. Joseph Smith opened by calling on the choir to sing a Hymn---sung 18th Hymn. The President then read a letter from Br. O[rson] Hyde giving an account of his journeys and success in his mission, which was listened to with intense interest; and the cenference, by vote, expressed their approbation of the style and spirit of said letter. The President then made remarks on the inclemency of the weather and the uncomfortable situation of the saints with regard to a place of worship, and a place of public entertainment...
					
						\hypertarget{EMS-1841-JS-1-5-OCT-1841-TS-b}%
					Pres.
						\marginnote{b}%
					Joseph Smith then made some corrections of doctrine in quoting a passage from \hyperref[1CORINTHIANS-12-28]{1 Cor. 12, 28}, showing it to be a principle of order or gradation in rising from one office to another in the Priesthood...
					
						\hypertarget{EMS-1841-JS-1-5-OCT-1841-TS-c}%
					President
						\marginnote{c}%
					Joseph Smith, by request of some of the Twelve, gave instructions on the doctrine of Baptism for the Dead; which was listened to with intense interest by the large assembly. The speaker presented ``Baptism for the Dead'' as the only way that men can appear as saviors on mount Zion. The proclamation of the first principles of the gospel was a means of salvatien to men individually, and it was the truth, not men that saved them; but men, by actively engaging in rites of salvation substitutionally, became instrumental in bringing multitudes of their kin into the kingdom of God.
						\hypertarget{EMS-1841-JS-1-5-OCT-1841-TS-d}%
					He
						\marginnote{d}%
					explained a difference between an angel and a ministering spirit; the one a resurrected or translated body, with its spirit, ministering to embodied spirits---the other a disembodied spirit, visiting and ministering to disembodied spirits. Jesus Christ became a minestering spirit, while his body laying in the sepulchre, to the spirits in prison; to fulfil an important part of his mission, without which he could not have perfected his work or entered into his rest. After his resurrection, he appeared as an angel to his disciples \&c. Translated bodies cannot enter into rest until they have undergone a change equivalent to death. Translated bodies are designed for future missions. The angel that appeared to John on the Isle of Patmos was a translated or resurrected body.---
						\hypertarget{EMS-1841-JS-1-5-OCT-1841-TS-e}%
					Jesus
						\marginnote{e}%
					Christ went in body, after his resurrection, to minister to translated and resurrected bodies. There has been a chain of authority and power from Adam down to the present time. The only way to obtain truth and wisdom, is not to ask it from books, but to go to God in prayer and obtain divine teaching. It is no more incredible that God should \textit{save} the dead, than that he should \textit{raise} the dead. There is never a time when the spirit is too old to approach God.
						\hypertarget{EMS-1841-JS-1-5-OCT-1841-TS-f}%
					All
						\marginnote{f}%
					are within the reach of pardoning mercy, who have not committed the unpardonable sin, which hath no forgiveness, neither in this world, nor in the world to come. There is a way to release the spirit of the dead; that is, by the power and authority of the Priest[h]ood---by binding and loosing on earth

						\hypertarget{EMS-1841-JS-1-5-OCT-1841-TS-g}%
					This
						\marginnote{g}%
					doctrine appears glorious, inasmuch as it exhibits the greatness of divine compassion and benevolence in the extent of the plan of human salvation. This glorious truth is well calculated to enlarge the understanding, and to sustain the soul under troubles, difficulties, and distresses.

						\hypertarget{EMS-1841-JS-1-5-OCT-1841-TS-h}%
					For
						\marginnote{h}%
					illustration the speaker presented, by supposition, the case of too men, brothers, equally intelligent, learned, virtuous and lovely, walking in uprightness and in all good conscience, so far as they had been able to discern duty from the muddy stream of tradition, or from the blotted page of the book of nature. One dies, and is buried, having never heard the gospel of reconciliation; to the other the message of salvation is sent, he hears and embraces it, and is made the heir of eternal life. Shall the one become a partaker of glory, and the other be consigned to hopeless perdition? Is there no chance for his escape? Sectarianism answers, ``none! none!! none!!!'' Such an idea is worse than atheism.
						\hypertarget{EMS-1841-JS-1-5-OCT-1841-TS-i}%
					The
						\marginnote{i}%
					truth shall break down and dash in pieces all such bigoted Pharisaism; the sects shall be sifted, the honest in heart brought out and their priests left in the midst of their corruption. The speaker then answered the objections urged against the Latter Day Saints for not admitting the validity of sectarian baptism, and for withholding fellowship from sectarian churches[.]
						\hypertarget{EMS-1841-JS-1-5-OCT-1841-TS-j}%
					It
						\marginnote{j}%
					was like putting new wine into old bottles and putting old wine into new bottles. What, new revelations in the old churches! New revelatiens knock out the bottom of their bottomless pit. New wine into old bottles!---the bottles burst and the wine runs out. What, Sadducees in the new church! Old wine in new leathern bottles will leak through the pores and escape; so the Saddacee saints mock at authority, kick out of the traces, and run to the mountains of perdition, leaving the long echo of their braying behind them.
						
						\hypertarget{EMS-1841-JS-1-5-OCT-1841-TS-k}%
					The
						\marginnote{k}%
					speaker then contrasted the charity of the sects, in denouncing all who disagree with them in opinion, and in joining in persecuting t[h]e saints, with the faith of the saints, who believe that even such may be saved in this world and in the world to come, (murderers and apostates excepted.)

						\hypertarget{EMS-1841-JS-1-5-OCT-1841-TS-l}%
					This
						\marginnote{l}%
					doctrine, he said, presented in a clear light, the wisdom and mercy of God, in preparing an ordinance for the salvation of the dead, being baptised by proxy, their names recorded in heaven, and they judged according to the deeds done in the body. This doctrine was the burden of the scriptures. Those saints who neglect it, in behalf of their deceased relatives, do it at the peril of their own salvation.

					The dispensation of the fulness of times will bring to light the things that have been revealed in all former dispensations, also other things that have not been before revealed. He shall send Elijah the prophe[t] \&c., and restore all things in Christ.

						\hypertarget{EMS-1841-JS-1-5-OCT-1841-TS-m}%
					The
						\marginnote{m}%
					speaker then announced, ``There shall be no more baptisms for the dead, until the ordinance can be attended to in the font of the Lord's House; and the church shall not hold another general conference, until they can meet in said house. \textit{For thus saith the Lord!}''...
					
					Pres't. Joseph Smith made a lengthy exposition of the condition of the temporal affairs of the church, the agency of which had been committed to him at a general conference in Quincy---explaining the manner that he had discharged the duties involved in that agency, and the condition of the lands and other property of the church...
					
			% --
			\paragraph{As Reported by Willard Richards}
			% --
				\label{EMS-1841-JS-1-5-OCT-1841-WR}
				
				\begin{description}
					\item[Print Sources]
				\end{description}	

				\begin{tabular}{p{0.5in}p{1in}p{0.5in}p{1in}}
					JSP	 	& ---		& JSR 		& --- \\
					DC	 	& ---		& HC 		& --- \\
					HDDC 	& ---		& TCHC 		& --- \\
				\end{tabular}
				% HDDC = woodford's DC diss; JSR = Marquardt's JS Revelations; TCHC = Vogel HC
				
				\begin{description}
					\item[Online Access] \href{https://www.josephsmithpapers.org/paper-summary/discourse-3october-1841-as-reported-by-willard-richards/1}{Here}
					% \item[Analysis] \hyperref[XX]{Here}
				\end{description}
				
				\begin{description}
					\item[Background]
				\end{description}
				
					% It might also be useful to give a two sentence context of the period: e.g., "This speech was given in Nauvoo to a small congregation one month prior to the destruction of the Nauvoo Expositor. These and other tensions may have been on the prophet's mind when he gave the speech."
				
				\begin{description}
					\item[Original Source]
				\end{description}
				
					% brief description of original source material, with reference to JSP or somewhere else for more info
					% Background material: it might be helpful to have a note that says whether a given document was presented as a speech, was written in a journal, etc.
					% Also a comment here about how reliable the source might be, or what shape it is in, if there are large omissions or parts that are hard to understand, and so on.
				
				\begin{description}
					\item[Text]
				\end{description}
				
						\hypertarget{EMS-1841-JS-1-5-OCT-1841-WR-a}%
					Saviors
						\marginnote{a}%
					shall come up on Mt Zion \&.--- what is it. \&\&--- To preach only? No. In every dispensation something else to do beside pr$\diamond\diamond$ do all--- neglect the Baptism for the dead.--- Cant that man enter into the fullness of his rest? \textit{[illegible]} cannot be saviors of man upon any any other principle than by re'ving revelation of all the things To trace our family by geneoloy [genealogy] \& Rev--- 1st principles of the Gospel no more to the B[aptism] for the Dead. then the dimist [dimmest] Star is to the Glo[r]y of Yonder Son.
						\hypertarget{EMS-1841-JS-1-5-OCT-1841-WR-b}%
					All
						\marginnote{b}%
					the kindred from the days of Adam dow[n] upon that princple, Universalism is nothing. Jesus contemplated \&.--- Jesus went as minstering angel Spirits--- according to what men do fame in the flesh \& live \&c--- Jesus could not have eneterd \& if he had not done this work. Paul--- not be made perferct.--- Jesus ministry spirit \& \& angel--- Spirits of Just men made perfect \& command a co of angels--- spirits \& bodies translated--- spirits--- minister to spirits--- angels to angels--- translated bodies cannot enter into rest--- till instructed <by them who had the Key> John you shall instruct or prophecy.--- prophets rose at the Resurrection---

						\hypertarget{EMS-1841-JS-1-5-OCT-1841-WR-c}%
					Jesus
						\marginnote{c}%
					went with his body \& Spirit---

					Dispensation of fulness \&c--- all dispensations--- things before since the world.---

					go to Books \textit{[illegible]} Almighty god rend the vails Said Doctrine incredible--- why any more incredible that God should save the Dead after the \textit{[illegible]}? men never is to ad living or Dead to service the Gospel.--- in or out of the world. unpardonable sin--- sin this world or the world

						\hypertarget{EMS-1841-JS-1-5-OCT-1841-WR-d}%
					Spirits
						\marginnote{d}%
					after the dissolution of the body--- God is bound to fall of that Soul.--- Seal on Earth \&c--- seal of Baptism for my father.--- sealed in heaven.--- Books opened--- of I get my father name Entered. seal him though. \&c--- 2. Sons born alike--- one dies character alike--- tomorrow Gospel come.--- M. Brunstons Negros Generation where did Sectarians get their Doctrine.--- from heathen refuse--- brayed common sense out of the world.--- Mahomet, Juyomet, Toroyes Doctrines.--- Old Joe--- \textit{[illegible]} of this thesis--- Blasphemy--- voice of God from heavens.--- creeds of men.--- why so tenacous are the mormons? old wine, mend bottles \&. New Rev in old churches--- old Sadness Presbytrian--- into a new church--- drawn ga$\diamond\diamond$ for charity--- man gone to flames of hell

					[\textit{Text in top third portion of page 3 is illegible}]

						\hypertarget{EMS-1841-JS-1-5-OCT-1841-WR-e}%
					born
						\marginnote{e}%
					of water by proxy--- A man is judged according to \textit{[2 words illegible]}

					Pulling them out of the fire---

					Martin Luther \& Alexander Campbell.

					2\supers{d.} acts--- no Holy Ghost

						\hypertarget{EMS-1841-JS-1-5-OCT-1841-WR-f}%
					Dispensation
						\marginnote{f}%
					of the fullness of times which reveals in that dispensations all those dispensat[ions] all things not before revealed--- Adam did not attend to his progenitors <partriarchal>--- Eligah \textit{[illegible]} Israel out of the groves. by the Keys what Paul did not.--- hearts of the children to their fathers \&c--- ordinances \&--- all things attended to from Adam down--- every genten [generation] that have rejected the children fathers must be denied. with \textit{[2 words illegible]} fathers--- Baptismal font Geneoly the Dead neglected their dead. doc$\diamond\diamond\diamond$ed \textit{[3 words illegible]}

						\hypertarget{EMS-1841-JS-1-5-OCT-1841-WR-g}%
					\phantom{ }
						\marginnote{g}%
					\textit{[illegible]} days \textit{[illegible]} attend to \textit{[illegible]} just men Spirit of \textit{[3 words illegible]} of angels--- \textit{[illegible]} the reaping \textit{[illegible]} of angels--- false spirits.--- \textit{[illegible]} of just men made perfect--- Angels \textit{[illegible]}.--- \textit{[2 words illegible]} death \textit{[2 words illegible]}---

						\hypertarget{EMS-1841-JS-1-5-OCT-1841-WR-h}%
					\phantom{ }
						\marginnote{h}%
					\textit{[4 words illegible]} prepare your temple \textit{[illegible]}.--- Do \textit{[2 words illegible]} a man has a perfect Right to be \textit{[illegible]} by any man who is a kin to him \textit{[4 words illegible]} if he can \textit{[illegible]} to Adam \textit{[4 words illegible]} 144 \textit{[illegible]}--- pray God send spirits \& angels to inform them Pray God <for others> \textit{[illegible]} of your hands \textit{[illegible]} you.--- Be baptized. not without

					\textit{[4 words illegible]}.--- angels \textit{[4 words illegible]} seek out \textit{[4 words illegible]} general conference [p. 4] \textit{[illegible]} held here until ye hold it in my house [\textit{3/4 page blank}]

				
		% -----
		\subsubsection{Memorandum, 2 October 1841}
		% -----
			\label{EMS-1841-JS-2-OCT-1841}
			% \label{EMS-1840-JS-DAY-3LETTERMONTH-YEAR}
				
			\begin{description}
				\item[Print Sources]
			\end{description}	

			\begin{tabular}{p{0.5in}p{1in}p{0.5in}p{1in}}
				JSP	 	& ---		& JSR 		& --- \\
				DC	 	& ---		& HC 		& --- \\
				HDDC 	& ---		& TCHC 		& --- \\
			\end{tabular}
			% HDDC = woodford's DC diss; JSR = Marquardt's JS Revelations; TCHC = Vogel HC
			
			\begin{description}
				\item[Online Access] \href{https://www.josephsmithpapers.org/paper-summary/memorandum-2october-1841/1}{Here}
				% \item[Analysis] \hyperref[XX]{Here}
			\end{description}
			
			\begin{description}
				\item[Background]
			\end{description}
			
				% It might also be useful to give a two sentence context of the period: e.g., "This speech was given in Nauvoo to a small congregation one month prior to the destruction of the Nauvoo Expositor. These and other tensions may have been on the prophet's mind when he gave the speech."
			
			\begin{description}
				\item[Original Source]
			\end{description}
			
				% brief description of original source material, with reference to JSP or somewhere else for more info
				% Background material: it might be helpful to have a note that says whether a given document was presented as a speech, was written in a journal, etc.
				% Also a comment here about how reliable the source might be, or what shape it is in, if there are large omissions or parts that are hard to understand, and so on.
			
			\begin{description}
				\item[Text]
			\end{description}
			
					\hypertarget{EMS-1841-JS-2-OCT-1841-a}%
				The
					\marginnote{a}%
				following was copied from a memorandum which was made upon the accasion of which it treats.

				``The Corner Stone of the Nauvoo House, in which the following records, were deposited, was laid by Joseph Smith, the President of the Church, to wit:---

				A Printed copy of the `Book of Mormon'

				A Revelation, given January 19\supersu{th} 1841

				The ``Times \& Seasons'', containing the charter of the Nauvoo House---

				A Journal of Heber [C.] Kimball.

				The Memorial of Lyman Wight to the Senate of the United States.

				A Book of Covenants

				No. 35 of the ``Times \& Seasons''

				The original Manuscript of the ``Book of Mormon''.

					\hypertarget{EMS-1841-JS-2-OCT-1841-b}%
				The
					\marginnote{b}%
				Persecution of the Church, in the State of Missouri, published in the Times \& Seasons.

				The Holy Bible---

				Silver Coin, to wit

				1 Half Dollar

				1 Quarter Dollar

				2 Dimes--- 2 Half Dimes---

				and

				1 Cop[p]er Coin.

				\begin{tightright}
					Deposited on the 2\supersu{nd} October 1841
				\end{tightright}
				
		% -----
		\subsubsection{Letter to Smith Tuttle, 9 October 1841}
		% -----
			\label{EMS-1841-JS-9-OCT-1841}
			% \label{EMS-1840-JS-DAY-3LETTERMONTH-YEAR}
				
			\begin{description}
				\item[Print Sources]
			\end{description}	

			\begin{tabular}{p{0.5in}p{1in}p{0.5in}p{1in}}
				JSP	 	& ---		& JSR 		& --- \\
				DC	 	& ---		& HC 		& --- \\
				HDDC 	& ---		& TCHC 		& --- \\
			\end{tabular}
			% HDDC = woodford's DC diss; JSR = Marquardt's JS Revelations; TCHC = Vogel HC
			
			\begin{description}
				\item[Online Access] \href{https://www.josephsmithpapers.org/paper-summary/letter-to-smith-tuttle-9october-1841/1}{Here}
				% \item[Analysis] \hyperref[XX]{Here}
			\end{description}
			
			\begin{description}
				\item[Background]
			\end{description}
			
				% It might also be useful to give a two sentence context of the period: e.g., "This speech was given in Nauvoo to a small congregation one month prior to the destruction of the Nauvoo Expositor. These and other tensions may have been on the prophet's mind when he gave the speech."
			
			\begin{description}
				\item[Original Source]
			\end{description}
			
				% brief description of original source material, with reference to JSP or somewhere else for more info
				% Background material: it might be helpful to have a note that says whether a given document was presented as a speech, was written in a journal, etc.
				% Also a comment here about how reliable the source might be, or what shape it is in, if there are large omissions or parts that are hard to understand, and so on.
			
			\begin{description}
				\item[Text]
			\end{description}
				
					\hypertarget{EMS-1841-JS-9-OCT-1841-a}%
				Nauvoo,
					\marginnote{a}%
				Ill. October 9\supersu{th} 1841

				Smith Tuttle Esqr.

				Dear Sir, Your kind letter of Sept. last was rec'd. during our Conference which is just over, containing a full \& particular explanation of every thing which gave rise to some feelings of disappointment in relation to our business transactions; and I will assure you, it has allayed, on our part, every prejudice. It breath[e]s the spirit of kindness \& truth. I will assure you that we exceedingly regret that there have been any grounds for hardness and disappointment. But so far as I am concerned, I must plead innocence; and you will consider me so, when you come to know all the facts--- I have done all that I could on my part. I will still do all that I can. I will not leave one stone unturned.---

					\hypertarget{EMS-1841-JS-9-OCT-1841-b}%
				Now
					\marginnote{b}%
				the facts are these. I sent my Brother Hyrum Smith \& Dr. [Isaac] Galland with means in their hands.--- say, not money, but with power to obtain every property or money which was necessary to enable them to fulfil the contract I made with Mr. [Horace] Hotchkiss. My brother was under the necessity of returning, in consequence of ill health, to this place, leaving the business in the hands of Dr. Galland, with the fullest expectation that he would make over the property or money to Mr. Hotchkiss, and make every thing square, so far as the interest is concerned, if not the principal. He was instructed to pay the interest that had accrued \& would accrue up to the fall of 1842, so as to be in advance of our indebtedness.
					\hypertarget{EMS-1841-JS-9-OCT-1841-c}%
				I
					\marginnote{c}%
				had also made arrangements with the eastern churches, \& had it in my power to fork over land for the whole debt; \& had expected that an arrangement of that kind would have been entered into. I am well assured that D\supers{r.} Galland did not lack <any> means whatever, to pay the interest, at any rate, if not the principal; \& why he has not done according to my instructions, God only knows. I do not feel to charge him with having done wrong, until I can investigate the matter, and ascertain for a certainty where the fault lies. It may be, that through sickness or disaster, this strange neglect has happened, I would to God the thing had not happened, When I read Mr. Hotchkiss' letter, I learned that he was dissatisfied.
					\hypertarget{EMS-1841-JS-9-OCT-1841-d}%
				I
					\marginnote{d}%
				thought he meant to oppress me, \& felt exceedingly mortified \& sorrowful in the midst of affliction, to think that he should distrust me for a moment, that I would not do all that was within my power. But upon hearing an explanation of the whole matter, my feelings are changed; and I think you all have had cause for complaining; but you will in the magnanimity of yr. good feelings, certainly not blame me when you find I have discharged an honorable duty on my part. I regret exceedingly that I did not know some months since, what I now know, so that I could have made another exertion before it got so late.

					\hypertarget{EMS-1841-JS-9-OCT-1841-e}%
				Cold
					\marginnote{e}%
				weather is now rolling in upon us. I have been confined here this season by sickness \& various other things that were beyond my control; such as having been demanded by the Governor of Mo. of the Governor of this State, \& he not having moral courage enough to resist the demand, although it was founded in injustice \& cruelty. I accordingly was taken prisoner, \& they put me to some ten or eleven hundred dollars expense \& trouble before I could be redeemed from under the difficulty; lawyers fees, witnesses--- \&c. \&c.--- But I am now clear from them once more, \& now in contemplating the face of the whole subject, I find that I am under the necessity of asking a little farther indulgence, <say> until next spring so that I may be enabled to recover myself; and then if God spares my life, \& gives me power to do so, I will come in person to yr. country, \& will never cease my labors until the whole matter is completely adjusted to the fullest satisfaction of all of you.
					\hypertarget{EMS-1841-JS-9-OCT-1841-f}%
				The
					\marginnote{f}%
				subject of your debt was presented fairly before our general Conference, on the first of the month, of some ten thousand people for their decission for the wisest \& best course, in relation to meeting your demands. The Twelve, as they are denominated in the ``Times and Seasons,'' were ordered by the Conference to make arrangements in the eastern branches of the church, ordering them to go to you \& turn over their property as you \& they could agree, \& take up our obligations \& bring them here, \& receive property here for them. \& I have been ordered by the Gen\supers{l.} Conference to write this letter to you, informing you of the measures which are about being taken to make all things right. I would inform you that Dr. Galland has not returned to the Western Country as yet.
					\hypertarget{EMS-1841-JS-9-OCT-1841-g}%
				He
					\marginnote{g}%
				has a considerable amt. of our money in his hands, which was to have been paid to you as was intended. He is on his way for aught we know, \& is retarded in his journey by some misfortune or other He may return, however, as yet, \& give a just and honorable account of himself. We hope this may be the case. I am sorrowful on account of yr. disappointments. It is a great disappointment to me as well as to yourselves. As to the growth of our place, it is very rapid; \& it would be more so, were it not for sickness \& death. There have been many deaths which leaves a melancholly reflection, but we can not help it. When God speaks from the heavens to call us home, we must submit to his mandate.--- And for your sincerity \& friendship, gentlemen, we have not the most distant doubt. We will not harbour any. We know it is for your interest to do us good, \& for our happiness \& welfare, to be punctual in the fulfilment of all our vows.
					\hypertarget{EMS-1841-JS-9-OCT-1841-h}%
				And
					\marginnote{h}%
				we think for the future you will have no cause of complaint. We intend to struggle with all the misfortunes of life, \& shoulder them all up handsomely \& honorably, even like men. We ask nothing, therefore, but what ought to be \sout{granted} <required> between man \& man, \& by those principles which bind man to man by kindred blood, in bearing our own part in every thing which duty calls us to do or not inferior to any of the human race, \& will be treated as such, although differing with some in matters of opinion in things, (viz:--- religious matters,) for which we only feel ourselves amenable to the Eternal God. And may God forbid that pride, ambition, a want of humility or any degree of importance, unjustly should have any dominion in our bosoms.
					\hypertarget{EMS-1841-JS-9-OCT-1841-i}%
				We
					\marginnote{i}%
				are the sons of Adam. We are the \uline{free} \uline{born} \uline{sons} \uline{of} \uline{America}, and as having been trampled upon, \& our constitutional rights taken from us, by a great many who boast themselves of being valiant in freedom's cause, while their hearts possess not a spark of its benign \& enlivening influence. This will afford a sufficient excuse, we hope, for any harsh remarks that may have been dropped by us, when we thought there was an assumption of superiority designed to gall our feelings. We are very sensative, as a people; we confess it, but we want to be pardoned for our sins, if any we have committed.

					\hypertarget{EMS-1841-JS-9-OCT-1841-j}%
				With
					\marginnote{j}%
				regard to the time when the first payment of interest should be made, it appears we did not understand each other, but it is a matter, which I hope we can amicably adjust when we see each other. I do not intend that it shall prevent our making an arrangement concerning the \uline{whole} \uline{matter}, for unless we can accomplish this, it will be useless for us to think of profitting by the purchase.

				With sentiments of respect, I remain as ever Yours \&c.--- \uline{Joseph Smith}
			
				<NAUVOO Ills. OCT 12>

				\begin{tightright}
					<Paid 25>
				\end{tightright}

				Smith Tuttle Esqr.

				New Haven

				C.t.
				
		% -----
		\subsubsection{Interview, 3 November 1841}
		% -----
			\label{EMS-1841-JS-3-NOV-1841}
			% \label{EMS-1840-JS-DAY-3LETTERMONTH-YEAR}
				
			\begin{description}
				\item[Print Sources]
			\end{description}	

			\begin{tabular}{p{0.5in}p{1in}p{0.5in}p{1in}}
				JSP	 	& ---		& JSR 		& --- \\
				DC	 	& ---		& HC 		& --- \\
				HDDC 	& ---		& TCHC 		& --- \\
			\end{tabular}
			% HDDC = woodford's DC diss; JSR = Marquardt's JS Revelations; TCHC = Vogel HC
			
			\begin{description}
				\item[Online Access] \href{https://www.josephsmithpapers.org/paper-summary/interview-3november-1841/1}{Here}
				% \item[Analysis] \hyperref[XX]{Here}
			\end{description}
			
			\begin{description}
				\item[Background]
			\end{description}
			
				% It might also be useful to give a two sentence context of the period: e.g., "This speech was given in Nauvoo to a small congregation one month prior to the destruction of the Nauvoo Expositor. These and other tensions may have been on the prophet's mind when he gave the speech."
			
			\begin{description}
				\item[Original Source]
			\end{description}
			
				% brief description of original source material, with reference to JSP or somewhere else for more info
				% Background material: it might be helpful to have a note that says whether a given document was presented as a speech, was written in a journal, etc.
				% Also a comment here about how reliable the source might be, or what shape it is in, if there are large omissions or parts that are hard to understand, and so on.
			
			\begin{description}
				\item[Text]
			\end{description}
			
					\hypertarget{EMS-1841-JS-3-NOV-1841-a}%
				A.
					\marginnote{a}%
				``You have the beginning of a great city here, Mr. Smith.''

				-[Here came in the more prominent objects of the city. The expense of the temple, Mr. Smith thought, would be \$200,000 or \$300,000. The temple is 127 feet side, by 88 feet front; and by its plan, which was kindly shown us, will fall short of some of our public buildings. As yet, only the foundations are laid. Mr. Smith then spoke of the ``false'' reports current about himself, and ``supposed we had heard enough of them.''

					\hypertarget{EMS-1841-JS-3-NOV-1841-b}%
				A.
					\marginnote{b}%
				``You know, sir, persecution sometimes drives `the \textit{wise} man mad.'''5

				Mr. S. (laughing.) ``Ah, sir, you must not put me among the wise men; my place is not there. I make no pretensions to piety, either. If you give me credit for any thing, let it be for being a \textit{good manager}. A good manager I do claim to be.''

					\hypertarget{EMS-1841-JS-3-NOV-1841-c}%
				A.
					\marginnote{c}%
				``You have great influence here, Mr. Smith.''

				Mr. S. ``Yes; I have. I bought 900 acres here, a few years ago, and they all have their lands of me. My influence, however, is ecclesiastical only; in civil affairs, I am but a common citizen. To be sure, I am a member of the City Council, and Lieutenant General of the Nauvoo Legion. I can command a thousand men to the field, at any moment, to support the laws. I had hard work to make them turn out and form the `Legion,' until I shouldered my musket, and entered the ranks myself.
					\hypertarget{EMS-1841-JS-3-NOV-1841-d}%
				Now,
					\marginnote{d}%
				they have nearly all provided themselves with a good uniform, poor as they are. By the way, we had a regular `set to' up here, a day or two since. The City Council ordered a liquor seller to leave the place, when his time was up; and, as he still remained, they directed that his house should be pulled down about his ears. They gave me a hand in the scrape; and I had occasion to knock a man down more than once. They mustered so strong an opposition, that it was either `knock down,' or `be knocked down.' We beat him off, at last; and are determined to have no grog shops in or about our grounds.''

				-[The conversation flowed on pleasantly, until my friend, to fill a pause that occurred, referred to my calling as a preacher.]-

					\hypertarget{EMS-1841-JS-3-NOV-1841-e}%
				Mr. S.
					\marginnote{e}%
				``Well I suppose (turning from me) he is one of the craft trained to his creed.''

				A. ``My creed, sir, is the New Testament.''

				Mr. S. ``Then, sir, we shall see truth just alike; for the scripture says, `They shall see, eye to eye.' All who are true men, must read the bible alike, must they not?''

				A. ``True, Mr. Smith; and yet I doubt if they will see it \textit{precisely} alike. If no two blades of grass are precisely alike, for a higher reason, it seems that no two intellects are.''

				Mr. S. (getting warm.) ``There---I told you so. You don't come here to seek truth. You begin with taking the place of opposition. Now, say what I may, you have but to answer, `No two men can see alike.'''

					\hypertarget{EMS-1841-JS-3-NOV-1841-f}%
				A.
					\marginnote{f}%
				``Mr. Smith, I said not that no two could see \textit{alike}; but that no two could see, on the whole, \textit{precisely} alike.''

				Mr. S. ``Does not the scripture say, `They shall see, eye to eye?'''

				A. ``Granted, sir; but be good enough to take a case; The words `all' and `all things' were brought up as meaning, at one time, universal creation. And again: `One believeth that he may \textit{eat} all \textit{things},' i.e. \textit{any thing}, or, as we say, \textit{every thing}.''

				Mr. S. ``You may explain away the bible, sir, as much as you please. I ask you, have you ever been baptized?''

				A. ``Yes, sir; I think I have.''

					\hypertarget{EMS-1841-JS-3-NOV-1841-g}%
				Mr.
					\marginnote{g}%
				S. ``Can you prophesy?''

				A. ``Well, sir, that depends on the meaning you give the word. I grant that it generally means to foretell; but I believe that it often means, \textit{to preach the gospel}. In this sense, sir, I can prophesy.''

				Mr. S. ``You lie, sir; and you know it.''

				A. ``It is as easy for me to impugn your motives, Mr. Smith, as for you to impugn mine.''

				Mr. S. ``I tell you, you don't seek to know the truth. You are a hypocrite: I saw it when you first began to speak.''

					\hypertarget{EMS-1841-JS-3-NOV-1841-h}%
				A.
					\marginnote{h}%
				``It is plain, Mr. Smith, that we differ in opinion. Now, one man's opinion is as good as another's, until some third party comes in to strike a balance between them.''

				Mr. S. ``I want no third party, sir. You are a fool, sir, to talk as you do. Have I not seen twice the years that you have? -[Joseph Smith is 36 years old; the speaker, A., was 10 years younger.]- I say, sir, you are no gentleman. I would'nt trust you with my purse across the street.''

					\hypertarget{EMS-1841-JS-3-NOV-1841-i}%
				\phantom{ }
					\marginnote{i}%
				-[Here my friend interposed, saying: ``I don't believe, Mr. Smith, that this gentleman came to your house to insult you. He had heard all sorts of accounts of your people, and came simply to see with his own eyes.'']-

				Mr. S. ``I have no ill feelings towards the gentleman. He is welcome to my house; but what I see to be the truth, I must speak out; I flatter no man. I tell you, sir, that man is a hypocrite. You'll find him out, if you're long enough with him. I tell you, I would'nt trust him as far as I could see him. What right has he to speak so to me? Am I not the leader of a great people? He, himself, will not blame me for speaking the truth plainly.''
				
		% -----
		\subsubsection{Discourse, 7 November 1841}
		% -----
			\label{EMS-1841-JS-7-NOV-1841}
			% \label{EMS-1840-JS-DAY-3LETTERMONTH-YEAR}
			
			% --
			\paragraph{As Reported by Wilford Woodruff}
			% --
				\label{EMS-1841-JS-7-NOV-1841-WW}
				
				\begin{description}
					\item[Print Sources]
				\end{description}	

				\begin{tabular}{p{0.5in}p{1in}p{0.5in}p{1in}}
					JSP	 	& ---		& JSR 		& --- \\
					DC	 	& ---		& HC 		& --- \\
					HDDC 	& ---		& TCHC 		& --- \\
				\end{tabular}
				% HDDC = woodford's DC diss; JSR = Marquardt's JS Revelations; TCHC = Vogel HC
				
				\begin{description}
					\item[Online Access] \href{https://www.josephsmithpapers.org/paper-summary/discourse-7november-1841-as-reported-by-wilford-woodruff/1}{Here}
					% \item[Analysis] \hyperref[XX]{Here}
				\end{description}
				
				\begin{description}
					\item[Background]
				\end{description}
				
					% It might also be useful to give a two sentence context of the period: e.g., "This speech was given in Nauvoo to a small congregation one month prior to the destruction of the Nauvoo Expositor. These and other tensions may have been on the prophet's mind when he gave the speech."
				
				\begin{description}
					\item[Original Source]
				\end{description}
				
					% brief description of original source material, with reference to JSP or somewhere else for more info
					% Background material: it might be helpful to have a note that says whether a given document was presented as a speech, was written in a journal, etc.
					% Also a comment here about how reliable the source might be, or what shape it is in, if there are large omissions or parts that are hard to understand, and so on.
				
				\begin{description}
					\item[Text]
				\end{description}
				
						\hypertarget{EMS-1841-JS-7-NOV-1841-WW-a}%
					Elder
						\marginnote{a}%
					W\supersu{m} Clark preached about 2 hours when Br Joseph arose \& reproved him as pharisaical \& hypocritical \& not edifying the people Br Joseph then deliverd unto us an edifying adress showing us what temperance faith virtue, charity \& truth was he also said if we did not accuse one another God would not accuse us \& if we had no accuser we should enter heaven he would take us there as his backload if we would not accuse him he would not accuse us
						\hypertarget{EMS-1841-JS-7-NOV-1841-WW-b}%
					\&
						\marginnote{b}%
					if we would throw a cloak of charity over his sins he would over ours for charity coverd a multitude of sins \& what many people called sin was not sin \& he did many things to break down superstition \& he would break it down he spoke of the curse of ham for laughing at Noah while in his wine but doing his harm.
					
			% --
			\paragraph{As Reported by Willard Richards}
			% --
				\label{EMS-1841-JS-7-NOV-1841-WR}
				
				\begin{description}
					\item[Print Sources]
				\end{description}	

				\begin{tabular}{p{0.5in}p{1in}p{0.5in}p{1in}}
					JSP	 	& ---		& JSR 		& --- \\
					DC	 	& ---		& HC 		& --- \\
					HDDC 	& ---		& TCHC 		& --- \\
				\end{tabular}
				% HDDC = woodford's DC diss; JSR = Marquardt's JS Revelations; TCHC = Vogel HC
				
				\begin{description}
					\item[Online Access] \href{https://www.josephsmithpapers.org/paper-summary/discourse-7november-1841-as-reported-by-willard-richards/1}{Here}
					% \item[Analysis] \hyperref[XX]{Here}
				\end{description}
				
				\begin{description}
					\item[Background]
				\end{description}
				
					% It might also be useful to give a two sentence context of the period: e.g., "This speech was given in Nauvoo to a small congregation one month prior to the destruction of the Nauvoo Expositor. These and other tensions may have been on the prophet's mind when he gave the speech."
				
				\begin{description}
					\item[Original Source]
				\end{description}
				
					% brief description of original source material, with reference to JSP or somewhere else for more info
					% Background material: it might be helpful to have a note that says whether a given document was presented as a speech, was written in a journal, etc.
					% Also a comment here about how reliable the source might be, or what shape it is in, if there are large omissions or parts that are hard to understand, and so on.
				
				\begin{description}
					\item[Text]
				\end{description}
				
						\hypertarget{EMS-1841-JS-7-NOV-1841-WR-a}%
					In
						\marginnote{a}%
					my last I believe, I gave you some sketches of Joseph's Sermon on baptism for the dead. I heard him last Sabbath, on Superstition \&c.---highly interesting.
					
					Brother Wm. O. Clark had been preaching to the congregation to pursuade them to Charity, temperence and every thing that is good and lovely, \&c. and without telling them what was good and lovely, \&c. and Brother Joseph shewed that this was the sectarian method of preaching and by such preaching no one could learn the principles of righteousness, and the world is ignorant of those principles;
						\hypertarget{EMS-1841-JS-7-NOV-1841-WR-b}%
					and
						\marginnote{b}%
					to do the people good we ought to tell them what those principles are instead of trying continually to enforce something they know nothing about. And he would tell them what virtue was, viz. to keep all the commandments of God without doubting or querying about it. God gives laws to suit the circumstances of his creatures. Laws in themselves contradictory; ``Thou shalt not kill;'' then to Abraham ``Slay thy son Isaac.'' Abraham rendered obedience, nothing doubting. This was virtue, perfecting his faith by works.
						\hypertarget{EMS-1841-JS-7-NOV-1841-WR-c}%
					Men
						\marginnote{c}%
					got up and pretended to preach by the power of the Holy Ghost, and tell nothing new nothing but what the people understood. This is false. When men preach by the power of the Holy Ghost they will tell something new, and if a man has nothing new to tell let him sit down to give room to others (this was in Nauvoo) lest he should come under some condemnation he would till something new.
					
						\hypertarget{EMS-1841-JS-7-NOV-1841-WR-d}%
					``No
						\marginnote{d}%
					man will be condemned before God who has no accuser.'' Like the woman who was taken in the very act and they came accusing \&c. ``Let him without sin cast the first stone.'' ``Where are thine accusers? hath no man condemned thee? No man Lord. Neither do I.'' Where two or three are agreed--- suppose it to be to take a glass of wine in the secret chamber and enjoy themselves for an hour and harm no one. They are agreed; who shall condemn them? No one; but one of the company in not agreed; he turns accuser. Condemnation follows.
						\hypertarget{EMS-1841-JS-7-NOV-1841-WR-e}%
					Drunkenness
						\marginnote{e}%
					is not good; but in such a case God might take no notice of it, if no one entered a complaint or accused the parties. Spirit of accusing is a spirit of evil and many may be condemned by it which otherwise might go clear. Our actions are between us and God if we infringe not on the rights of others. If the brethren would not acuse him he would not acuse them, but would take them all on his back and bear them safe through the gates into the kingdom, \&c. \&c. \&c.
					
		% -----
		\subsubsection{Letter to John M. Bernhisel, 16 November 1841}
		% -----
			\label{EMS-1841-JS-16-NOV-1841}
			% \label{EMS-1840-JS-DAY-3LETTERMONTH-YEAR}
				
			\begin{description}
				\item[Print Sources]
			\end{description}	

			\begin{tabular}{p{0.5in}p{1in}p{0.5in}p{1in}}
				JSP	 	& ---		& JSR 		& --- \\
				DC	 	& ---		& HC 		& --- \\
				HDDC 	& ---		& TCHC 		& --- \\
			\end{tabular}
			% HDDC = woodford's DC diss; JSR = Marquardt's JS Revelations; TCHC = Vogel HC
			
			\begin{description}
				\item[Online Access] \href{https://www.josephsmithpapers.org/paper-summary/letter-to-johnm-bernhisel-16november-1841/1}{Here}
				% \item[Analysis] \hyperref[XX]{Here}
			\end{description}
			
			\begin{description}
				\item[Background]
			\end{description}
			
				% It might also be useful to give a two sentence context of the period: e.g., "This speech was given in Nauvoo to a small congregation one month prior to the destruction of the Nauvoo Expositor. These and other tensions may have been on the prophet's mind when he gave the speech."
			
			\begin{description}
				\item[Original Source]
			\end{description}
			
				% brief description of original source material, with reference to JSP or somewhere else for more info
				% Background material: it might be helpful to have a note that says whether a given document was presented as a speech, was written in a journal, etc.
				% Also a comment here about how reliable the source might be, or what shape it is in, if there are large omissions or parts that are hard to understand, and so on.
			
			\begin{description}
				\item[Text]
			\end{description}
			
				\begin{tightright}
						\hypertarget{EMS-1841-JS-16-NOV-1841-a}%
					Nauvoo
						\marginnote{a}%
					November 16, 1841
				\end{tightright}

				Dear Sir

				I received your kind present by the hand of E\supersu{r}\supers{.} [Wilford] Woodruff \& feel myself under many obligations for this mark of your esteem \& friendship which to me is the more interesting as it unfolds \& developes many things that are of great importance to this generation \& corresponds with \& supports the testimony of the Book of Mormon; I have read the volumnes with the greatest interest \& pleasure \& must say that of all histories that have been written pertaining to the antiquities of this country it is the most correct luminous \& comprihensive.---

					\hypertarget{EMS-1841-JS-16-NOV-1841-b}%
				In
					\marginnote{b}%
				regard to the land referred to by you I would simply state that I have lands both in and out of the City some of which I hold deeds for and others bonds for deeds when you come which I hope will be as soon as convenient you can make such a selection from among those as shall best meet with your veiws \& feelings. In gratefull remembrance of your kindness I remain your affectionate Brother in the bonds of the

				\begin{tightright}
					Everlasting Covenant

					Joseph Smith
				\end{tightright}

				To D\supersu{r}\supers{.} [John M.] Bernhisel

				[\textit{page [2] blank}]

				[\textit{page [3] blank}]

				<Nauvoo Ill>

				<Nov 23\supers{d}>

				<25>

				D\supersu{r}\supers{.} Bernhisel care of Lucian R. Foster

				No 13 Oliver St

				New York

				<left by Elder [James] Divine at the bequest of Eldr \uline{Foster}>
				
		% -----
		\subsubsection{Remarks, 28 November 1841}
		% -----
			\label{EMS-1841-JS-28-NOV-1841}
			% \label{EMS-1840-JS-DAY-3LETTERMONTH-YEAR}
				
			\begin{description}
				\item[Print Sources]
			\end{description}	

			\begin{tabular}{p{0.5in}p{1in}p{0.5in}p{1in}}
				JSP	 	& ---		& JSR 		& --- \\
				DC	 	& ---		& HC 		& --- \\
				HDDC 	& ---		& TCHC 		& --- \\
			\end{tabular}
			% HDDC = woodford's DC diss; JSR = Marquardt's JS Revelations; TCHC = Vogel HC
			
			\begin{description}
				\item[Online Access] \href{https://www.josephsmithpapers.org/paper-summary/remarks-28-november-1841/1}{Here}
				% \item[Analysis] \hyperref[XX]{Here}
			\end{description}
			
			\begin{description}
				\item[Background]
			\end{description}
			
				% It might also be useful to give a two sentence context of the period: e.g., "This speech was given in Nauvoo to a small congregation one month prior to the destruction of the Nauvoo Expositor. These and other tensions may have been on the prophet's mind when he gave the speech."
			
			\begin{description}
				\item[Original Source]
			\end{description}
			
				% brief description of original source material, with reference to JSP or somewhere else for more info
				% Background material: it might be helpful to have a note that says whether a given document was presented as a speech, was written in a journal, etc.
				% Also a comment here about how reliable the source might be, or what shape it is in, if there are large omissions or parts that are hard to understand, and so on.
			
			\begin{description}
				\item[Text]
			\end{description}
			
				28[\supersu{th}] Sunday I spent the day at B[righam] Young in company with Joseph \& the Twelve in conversing upon a variety of subjects it was an interesting day Elder Joseph Fielding was present he had been in England four years we also saw a number of english Brethren Joseph said the Book of Mormon was the most correct of any Book on earth \& the keystone of our religion \& a man would get nearer to God by abiding by its precepts than any other Book
				
		% -----
		\subsubsection{Affidavit, 29 November 1841}
		% -----
			\label{EMS-1841-JS-29-NOV-1841}
			% \label{EMS-1840-JS-DAY-3LETTERMONTH-YEAR}
				
			\begin{description}
				\item[Print Sources]
			\end{description}	

			\begin{tabular}{p{0.5in}p{1in}p{0.5in}p{1in}}
				JSP	 	& ---		& JSR 		& --- \\
				DC	 	& ---		& HC 		& --- \\
				HDDC 	& ---		& TCHC 		& --- \\
			\end{tabular}
			% HDDC = woodford's DC diss; JSR = Marquardt's JS Revelations; TCHC = Vogel HC
			
			\begin{description}
				\item[Online Access] \href{https://www.josephsmithpapers.org/paper-summary/affidavit-29november-1841/1}{Here}
				% \item[Analysis] \hyperref[XX]{Here}
			\end{description}
			
			\begin{description}
				\item[Background]
			\end{description}
			
				% It might also be useful to give a two sentence context of the period: e.g., "This speech was given in Nauvoo to a small congregation one month prior to the destruction of the Nauvoo Expositor. These and other tensions may have been on the prophet's mind when he gave the speech."
			
			\begin{description}
				\item[Original Source]
			\end{description}
			
				% brief description of original source material, with reference to JSP or somewhere else for more info
				% Background material: it might be helpful to have a note that says whether a given document was presented as a speech, was written in a journal, etc.
				% Also a comment here about how reliable the source might be, or what shape it is in, if there are large omissions or parts that are hard to understand, and so on.
			
			\begin{description}
				\item[Text]
			\end{description}
			
				\begin{tightcenter}
						\hypertarget{EMS-1841-JS-29-NOV-1841-a}%
					PRES'T.
						\marginnote{a}%
					J. SMITH'S AFFIDAVIT.
				\end{tightcenter}

				\begin{tightright}
					\textit{City of Nauvoo, Ill}.,

					\textit{Nov.} 29\textit{th A. D. 1841}.
				\end{tightright}

				\textsc{To the Public}:---

				The transpiration of recent events makes it criminal for me to remain longer silent. The tongue of the vile yet speaks, and sends forth the poison of asps---the ears of the spoiler yet hear, and he puts forth his hands to iniquity. It has been proclaimed upon the house-top and in the secret chamber, in the public walks and private circle, throughout the length and breadth of this vast continent, that stealing by the Latter Day Saints has received my approval; nay, that I have taught the doctrine, encouraged them in plunder, and led on the van---than which nothing is more foreign from my heart.
					\hypertarget{EMS-1841-JS-29-NOV-1841-b}%
				I
					\marginnote{b}%
				disfellowship the perpetrators of all such abominations---they are devils and not saints, totally unfit for the society of Christians, or men. It is true that some professing to be Latter Day Saints have taught such vile heresies, but all are not Israel that are of Israel;10 and I wish it to be distinctly understood in all coming time, that the church over which I have the honor of presiding will ever set its brows like brass, and its face like steel, against all such abominable acts of villany and crime; and to this end I append my affidavit of disavowal taken this day before General [John C.] Bennett, that there may be no mistake here after as to my real sentiments, or those of the leaders of the church, in relation to this important matter,------

				\begin{tabular}{p{5in}p{1in}}
				\textsc{State of Illinois},) & ss \\
				Hancock County.) &  
				\end{tabular}

					\hypertarget{EMS-1841-JS-29-NOV-1841-c}%
				Before
					\marginnote{c}%
				me, John C. Bennett, Mayor of the City of Nauvoo, personally came Joseph Smith, President of the Church of Jesus Christ of Latter Day Saints, (commonly called Mormons,) who being duly sworn according to law, deposeth and saith, that he has never directly or indirectly encouraged the purloining of property, or taught the doctrine of stealing, or any other evil practice, and that all such vile and unlawful acts will ever receive his unqualified and unreserved disapproval, and the most vigorous opposition of the church over which he presides, and further this deponent saith not.

				\begin{tightright}
					JOSEPH SMITH, President of the Church of Jesus Christ of Latter Day Saints.
				\end{tightright}

				Sworn to, and subscribed before me, at my office, in the City of Nauvoo, this twenty ninth day of November, Anno Domini 1841.

				\begin{tightright}
					-[L. S.]- JOHN C. BENNETT,

					Mayor of the City of Nauvoo.
				\end{tightright}

					\hypertarget{EMS-1841-JS-29-NOV-1841-d}%
				Now
					\marginnote{d}%
				it is to be hoped that none will hereafter be so reckless as to state that I, or the church to which I belong, approve of thieving---but that all the friends of law and order will join in ferreting out thieves wherever, and whenever, they may be found, and assist in bringing them to that condign punishment which such infamous crimes so richly merit.

				\begin{tightright}
					JOSEPH SMITH, President of the Church of Jesus Christ of Latter Day Saints.
				\end{tightright}
				
		% -----
		\subsubsection{Revelation, 2 December 1841}
		% -----
			\label{EMS-1841-JS-2-DEC-1841}
			% \label{EMS-1840-JS-DAY-3LETTERMONTH-YEAR}
				
			\begin{description}
				\item[Print Sources]
			\end{description}	

			\begin{tabular}{p{0.5in}p{1in}p{0.5in}p{1in}}
				JSP	 	& ---		& JSR 		& --- \\
				DC	 	& ---		& HC 		& --- \\
				HDDC 	& ---		& TCHC 		& --- \\
			\end{tabular}
			% HDDC = woodford's DC diss; JSR = Marquardt's JS Revelations; TCHC = Vogel HC
			
			\begin{description}
				\item[Online Access] \href{https://www.josephsmithpapers.org/paper-summary/revelation-2-december-1841/1}{Here}
				% \item[Analysis] \hyperref[XX]{Here}
			\end{description}
			
			\begin{description}
				\item[Background]
			\end{description}
			
				% It might also be useful to give a two sentence context of the period: e.g., "This speech was given in Nauvoo to a small congregation one month prior to the destruction of the Nauvoo Expositor. These and other tensions may have been on the prophet's mind when he gave the speech."
			
			\begin{description}
				\item[Original Source]
			\end{description}
			
				% brief description of original source material, with reference to JSP or somewhere else for more info
				% Background material: it might be helpful to have a note that says whether a given document was presented as a speech, was written in a journal, etc.
				% Also a comment here about how reliable the source might be, or what shape it is in, if there are large omissions or parts that are hard to understand, and so on.
			
			\begin{description}
				\item[Text]
			\end{description}
			
					\hypertarget{EMS-1841-JS-2-DEC-1841-a}%
				A.
					\marginnote{a}%
				Revelation Given Dc\supersu{r}\supers{.} 2\supersu{d}\supers{.} \sout{1842} 1841. N. M. Hyde [Marinda Nancy Johnson Hyde]
				
				Verily thus saith the Lord unto you my servant Joseph. that in as much as you have called upon me to know my will concerning my handmaid Nancy Marinda Hyde Behold it is my will that she should have a better place prepared for her than that in which she now lives, in order that her life may be spared unto her; Therefore go and say unto my servant Ebenezer Robinson. \& To my handmaid his wife, Let them open their doors and take her and her children into their house.
					\hypertarget{EMS-1841-JS-2-DEC-1841-b}%
				and
					\marginnote{b}%
				take care of them faithfully and kindly until my Servant Orson Hyde returns from his mission or until some other provision can be made for her welfare \& safety: Let them do these things and spare not. and I the Lord will bless them \& heal them. if they do it not grudgingly saith the Lord God. and she shall be a blessing unto them,--- and let my handmaid Nancy Marinda Hyde hearken to the counsel of my servant Joseph in all things whatsoever he shall teach unto her, and it shall be a blessing upon her and upon her children after her. unto her \hyperref[JUSTIFICATION]{Justification} saith the Lord. [\textit{6 lines blank}]
				
		% -----
		\subsubsection{Letter to Church Leaders in Kirtland, Ohio, 15 December 1841}
		% -----
			\label{EMS-1841-JS-15-DEC-1841}
			% \label{EMS-1840-JS-DAY-3LETTERMONTH-YEAR}
				
			\begin{description}
				\item[Print Sources]
			\end{description}	

			\begin{tabular}{p{0.5in}p{1in}p{0.5in}p{1in}}
				JSP	 	& ---		& JSR 		& --- \\
				DC	 	& ---		& HC 		& --- \\
				HDDC 	& ---		& TCHC 		& --- \\
			\end{tabular}
			% HDDC = woodford's DC diss; JSR = Marquardt's JS Revelations; TCHC = Vogel HC
			
			\begin{description}
				\item[Online Access] \href{https://www.josephsmithpapers.org/paper-summary/letter-to-church-leaders-in-kirtland-ohio-15-december-1841/1}{Here}
				% \item[Analysis] \hyperref[XX]{Here}
			\end{description}
			
			\begin{description}
				\item[Background]
			\end{description}
			
				% It might also be useful to give a two sentence context of the period: e.g., "This speech was given in Nauvoo to a small congregation one month prior to the destruction of the Nauvoo Expositor. These and other tensions may have been on the prophet's mind when he gave the speech."
			
			\begin{description}
				\item[Original Source]
			\end{description}
			
				% brief description of original source material, with reference to JSP or somewhere else for more info
				% Background material: it might be helpful to have a note that says whether a given document was presented as a speech, was written in a journal, etc.
				% Also a comment here about how reliable the source might be, or what shape it is in, if there are large omissions or parts that are hard to understand, and so on.
			
			\begin{description}
				\item[Text]
			\end{description}
			
					\hypertarget{EMS-1841-JS-15-DEC-1841-a}%
				``It
					\marginnote{a}%
				remains for Almon Babbitt to offer satisfaction. if he wishes so to do, according to the minutes of the Conference, \uline{You} \uline{are} \uline{doubtless} \uline{all} \uline{well} \uline{aware} \uline{that} \uline{all} \uline{the} \uline{stakes} \uline{except} \uline{those} \uline{in} \uline{Hancock}, \uline{Co}.--- \uline{Illinois}, \uline{\&} \uline{Lee} \uline{county} \uline{Iowa}. \uline{were} \uline{discontinued} \uline{some} \uline{time} \uline{Since} \uline{by} \uline{the} \uline{First} \uline{Presidency}, \uline{as} \uline{published} \uline{in} \uline{the} \uline{Times} \uline{and} \uline{Season}; but as it appears that there are many in Kirtland who desire to remain there.
					\hypertarget{EMS-1841-JS-15-DEC-1841-b}%
				and
					\marginnote{b}%
				build up that place, and as you have made great exertions, according to your letter, to establish a printing press, \& take care of the poor, \&\supers{c.} since that period, you may as well continue operations according to your designs, \& go on with your printing, \& do what you can in Righteousness to build up Kirtland but do not suffer yourselves to harbor the Idea that Kirtland will rise on the ruins of Nauvoo. It is the privilege of brethren emigrating from any quarter To come to this place, and it is not right to attempt to persuade those who desire it, to stop short,'' (Extract from the letter of the presidents in reply--- Decr 15[\supersu{th}] 1841.) 
				
		% -----
		\subsubsection{Discourse, 19 December 1841}
		% -----
			\label{EMS-1841-JS-19-DEC-1841}
			% \label{EMS-1840-JS-DAY-3LETTERMONTH-YEAR}
			
			% --
			\paragraph{As Reported by Wilford Woodruff}
			% --
				\label{EMS-1841-JS-19-DEC-1841-WW}
				
				\begin{description}
					\item[Print Sources]
				\end{description}	

				\begin{tabular}{p{0.5in}p{1in}p{0.5in}p{1in}}
					JSP	 	& ---		& JSR 		& --- \\
					DC	 	& ---		& HC 		& --- \\
					HDDC 	& ---		& TCHC 		& --- \\
				\end{tabular}
				% HDDC = woodford's DC diss; JSR = Marquardt's JS Revelations; TCHC = Vogel HC
				
				\begin{description}
					\item[Online Access] \href{https://www.josephsmithpapers.org/paper-summary/letter-to-church-leaders-in-kirtland-ohio-15-december-1841/1}{Here}
					% \item[Analysis] \hyperref[XX]{Here}
				\end{description}
				
				\begin{description}
					\item[Background]
				\end{description}
				
					% It might also be useful to give a two sentence context of the period: e.g., "This speech was given in Nauvoo to a small congregation one month prior to the destruction of the Nauvoo Expositor. These and other tensions may have been on the prophet's mind when he gave the speech."
				
				\begin{description}
					\item[Original Source]
				\end{description}
				
					% brief description of original source material, with reference to JSP or somewhere else for more info
					% Background material: it might be helpful to have a note that says whether a given document was presented as a speech, was written in a journal, etc.
					% Also a comment here about how reliable the source might be, or what shape it is in, if there are large omissions or parts that are hard to understand, and so on.
				
				\begin{description}
					\item[Text]
				\end{description}
				
						\hypertarget{EMS-1841-JS-19-DEC-1841-WW-a}%
					Joseph
						\marginnote{a}%
					the Seer arose \& read a chapter in the New Testiment containing the parable of the vine \& its branches \& explained it much to our edification \& said ``if we kept the commandments of God we should bring forth fruit \& be the friends of God \& know what our Lord did, ``Some say Joseph is a fallen Prophet because he does not bring forth more of the word of the Lord'' ``Why does \uline{he} not do it'' are we able to receive it No (says he) not one in this room,
						\hypertarget{EMS-1841-JS-19-DEC-1841-WW-b}%
					He
						\marginnote{b}%
					then chastized us for our wickedness \& unbelief knowing that whom the Lord loveth he chasteneth \& scourgeth evry son \& daughter whom He receiveth \& if we do not receive chastizements then are we Bastards \& not Sons.
						\hypertarget{EMS-1841-JS-19-DEC-1841-WW-c}%
					On
						\marginnote{c}%
					Revelation He said ``A man would command his son to dig potatoes, saddle his horse but before he had done either tell him to do somthing els, this is all consider[ed] right ``But as soon as the Lord gives a commandment \& revokes that decree \& commands somthing els then the prophet is considerd fallen \&c'' Because we will not receive chastizment at the hand of the Prophet \& Apostles the Lord chastizeth us with sickness \& death ``Let not any man publish his own righteousness for others can do that for him.''
						\hypertarget{EMS-1841-JS-19-DEC-1841-WW-d}%
					sooner
						\marginnote{d}%
					let him confess his sins \& then he will be forgiven \& he will bring forth more fruit when A man is chastized he gets angry \& will not endure it. The reason we do not have the secrets of the Lord revealed unto us is because we do not keep them but reveal them, we do not keep our own secrets but reveal our difficulties to the world even to our enemies then how would we keep the secrets of the Lord Joseph says I can keep a secret till dooms day
					
						\hypertarget{EMS-1841-JS-19-DEC-1841-WW-e}%
					He
						\marginnote{e}%
					spoke of love what greater love hath any man than that he lay down his life for his friend then why not fight for our friend untill we die \& many other things of interest was spoken.
					
		% -----
		\subsubsection{Revelation, circa 22 December 1841--A}
		% -----
			\label{EMS-1841-JS-22-DEC-1841-A}
			% \label{EMS-1840-JS-DAY-3LETTERMONTH-YEAR}
				
			\begin{description}
				\item[Print Sources]
			\end{description}	

			\begin{tabular}{p{0.5in}p{1in}p{0.5in}p{1in}}
				JSP	 	& ---		& JSR 		& --- \\
				DC	 	& ---		& HC 		& --- \\
				HDDC 	& ---		& TCHC 		& --- \\
			\end{tabular}
			% HDDC = woodford's DC diss; JSR = Marquardt's JS Revelations; TCHC = Vogel HC
			
			\begin{description}
				\item[Online Access] \href{https://www.josephsmithpapers.org/paper-summary/revelation-circa-22-december-1841-a/1}{Here}
				% \item[Analysis] \hyperref[XX]{Here}
			\end{description}
			
			\begin{description}
				\item[Background]
			\end{description}
			
				% It might also be useful to give a two sentence context of the period: e.g., "This speech was given in Nauvoo to a small congregation one month prior to the destruction of the Nauvoo Expositor. These and other tensions may have been on the prophet's mind when he gave the speech."
			
			\begin{description}
				\item[Original Source]
			\end{description}
			
				% brief description of original source material, with reference to JSP or somewhere else for more info
				% Background material: it might be helpful to have a note that says whether a given document was presented as a speech, was written in a journal, etc.
				% Also a comment here about how reliable the source might be, or what shape it is in, if there are large omissions or parts that are hard to understand, and so on.
			
			\begin{description}
				\item[Text]
			\end{description}
			
				Verily thus saith the Lord. unto my serva[n]ts the 12 let them appoint unto my servat A[mos] Fuller a mission to preach my gospel. unto the children of man. as it shall be manifested unto them by my Holy Spi[ri]t, Amen.------
				
		% -----
		\subsubsection{Revelation, circa 22 December 1841--B}
		% -----
			\label{EMS-1841-JS-22-DEC-1841-B}
			% \label{EMS-1840-JS-DAY-3LETTERMONTH-YEAR}
				
			\begin{description}
				\item[Print Sources]
			\end{description}	

			\begin{tabular}{p{0.5in}p{1in}p{0.5in}p{1in}}
				JSP	 	& ---		& JSR 		& --- \\
				DC	 	& ---		& HC 		& --- \\
				HDDC 	& ---		& TCHC 		& --- \\
			\end{tabular}
			% HDDC = woodford's DC diss; JSR = Marquardt's JS Revelations; TCHC = Vogel HC
			
			\begin{description}
				\item[Online Access] \href{https://www.josephsmithpapers.org/paper-summary/revelation-circa-22-december-1841-b/1}{Here}
				% \item[Analysis] \hyperref[XX]{Here}
			\end{description}
			
			\begin{description}
				\item[Background]
			\end{description}
			
				% It might also be useful to give a two sentence context of the period: e.g., "This speech was given in Nauvoo to a small congregation one month prior to the destruction of the Nauvoo Expositor. These and other tensions may have been on the prophet's mind when he gave the speech."
			
			\begin{description}
				\item[Original Source]
			\end{description}
			
				% brief description of original source material, with reference to JSP or somewhere else for more info
				% Background material: it might be helpful to have a note that says whether a given document was presented as a speech, was written in a journal, etc.
				% Also a comment here about how reliable the source might be, or what shape it is in, if there are large omissions or parts that are hard to understand, and so on.
			
			\begin{description}
				\item[Text]
			\end{description}
			
				And the word of the Lord came \sout{say} verily thus Saith the Lord. Let my Serv[a]nt John Snider take a mission. To the Eastern Continent. unto all the confernces that are now sitting in that region \& let him carry a package of Epistles, that shall be writtn by my Servants the 12, making known unto them th[e]ir duties concer[n]ing the buidig [building] of my houses which I have appoi[n]ted unto you saith the Lord. that they may bri[n]g their Go[ld] \& their Si[lver] \& their precious Stones. \& the box tree \& th[e] Fir tree. \& all fine wood. \& beautify the place of [m]y sanctuary saith the Lord. \& let him return with all means speedily which shall be put into his hands even so Amen.
				
		% -----
		\subsubsection{Discourse, 26 December 1841}
		% -----
			\label{EMS-1841-JS-26-DEC-1841}
			% \label{EMS-1840-JS-DAY-3LETTERMONTH-YEAR}
			
			% --
			\paragraph{As Reported by Willard Richards}
			% --
				\label{EMS-1841-JS-26-DEC-1841-WR}
				
				\begin{description}
					\item[Print Sources]
				\end{description}	

				\begin{tabular}{p{0.5in}p{1in}p{0.5in}p{1in}}
					JSP	 	& ---		& JSR 		& --- \\
					DC	 	& ---		& HC 		& --- \\
					HDDC 	& ---		& TCHC 		& --- \\
				\end{tabular}
				% HDDC = woodford's DC diss; JSR = Marquardt's JS Revelations; TCHC = Vogel HC
				
				\begin{description}
					\item[Online Access] \href{https://www.josephsmithpapers.org/paper-summary/discourse-26-december-1841-as-reported-by-willard-richards/1}{Here}
					% \item[Analysis] \hyperref[XX]{Here}
				\end{description}
				
				\begin{description}
					\item[Background]
				\end{description}
				
					% It might also be useful to give a two sentence context of the period: e.g., "This speech was given in Nauvoo to a small congregation one month prior to the destruction of the Nauvoo Expositor. These and other tensions may have been on the prophet's mind when he gave the speech."
				
				\begin{description}
					\item[Original Source]
				\end{description}
				
					% brief description of original source material, with reference to JSP or somewhere else for more info
					% Background material: it might be helpful to have a note that says whether a given document was presented as a speech, was written in a journal, etc.
					% Also a comment here about how reliable the source might be, or what shape it is in, if there are large omissions or parts that are hard to understand, and so on.
				
				\begin{description}
					\item[Text]
				\end{description}
				
						\hypertarget{EMS-1841-JS-26-DEC-1841-WR-a}%
					President Joseph read the 13\supersu{th}\supers{.} chap of 1st corinthians and a part of the 14 chap, and remarked that the gift of Tongues was necessary in the church; <but> That if satan could not speak in tongues he could not tempt a Dutchman, or any other nation, but the English, for he can tempt the Englishman, for he has tempted me, \& I am an Englishman;
						\hypertarget{EMS-1841-JS-26-DEC-1841-WR-b}%
					but
						\marginnote{b}%
					the Gift of Tongues, by the power of the Holy Ghost, in the church, is for the benefit of the servants of God to preach to unbelievers, as on the days of Pentecost, when devout men from evry nation shall assemble to hear of the things of God. let the <elders> preach to them in their own Mother tongue. whither it is German, French, Spanish or Irish. or any other. \& let those interpret who understand the \sout{tongue} <Language> spoken. \sout{by} in their mother tongue. \& this is what the Apostle meant. in 1s[t] corinthians 14:27.